% !TEX program = pdfLaTeX
% !TEX encoding = UTF-8
\documentclass[11pt,a4paper]{article}
\usepackage[utf8]{inputenc}
\usepackage{amsmath,amssymb,amsthm,mathtools,enumitem}
\usepackage[a4paper,margin=2.3cm]{geometry}
\usepackage{xcolor}
\usepackage[hypertexnames=false,colorlinks=true,linkcolor=StudyBlue,bookmarksopen=true]{hyperref}
\usepackage{compact-math-notes}
\usepackage{needspace}
\setlength{\parindent}{0pt}
\setlength{\emergencystretch}{1em}
\setlist{topsep=3pt,itemsep=2pt,parsep=0pt,font=\color{StudyBlue}\bfseries}
\newcommand{\R}{\mathbb R}
\newcommand{\N}{\mathbb N}
\newcommand{\Q}{\mathbb Q}
\providecommand{\SF}{SF^+}
\providecommand{\mS}{m\Sigma}
\providecommand{\ind}[1]{\mathbf{1}_{#1}}
\renewcommand{\qedsymbol}{\textnormal{QED}}
\raggedbottom
% Study hierarchy: colors and weight are confined to titles and part labels.
\definecolor{StudyRed}{HTML}{B51F2C}
\definecolor{StudyBlue}{HTML}{12356D}
\renewcommand{\contentsname}{\textcolor{StudyBlue}{Contents}}
\newcommand{\studyheading}[1]{{\color{StudyBlue}\normalfont\bfseries\boldmath #1}}
\newcommand{\studylabel}[1]{\textcolor{StudyBlue}{\textnormal{\bfseries #1}}}
\linespread{1.03}
\newcommand{\StudyParagraphSpacing}{\setlength{\parskip}{2pt plus .7pt minus .4pt}}
\AddToHook{cmd/normalsize/after}{\StudyParagraphSpacing}
\AtBeginDocument{\StudyParagraphSpacing}
\clubpenalty=10000
\widowpenalty=10000
\displaywidowpenalty=10000
\makeatletter
\renewcommand{\section}{\@startsection{section}{1}{0pt}
  {12pt plus 3pt minus 2pt}{7pt plus 2pt}{\normalfont\Large\bfseries\boldmath\color{StudyRed}}}
\renewcommand{\subsection}{\@startsection{subsection}{2}{0pt}
  {9pt plus 2pt minus 1pt}{3pt plus 1pt}{\normalfont\normalsize\bfseries\boldmath\color{StudyBlue}}}
\renewcommand{\subsubsection}{\@startsection{subsubsection}{3}{0pt}
  {7pt plus 2pt minus 1pt}{3pt plus 1pt}{\normalfont\normalsize\bfseries\boldmath\color{StudyBlue}}}
\renewcommand{\paragraph}{\@startsection{paragraph}{4}{0pt}
  {6pt plus 2pt minus 1pt}{2pt plus 1pt}{\normalfont\normalsize\bfseries\boldmath\color{StudyBlue}}}
\renewenvironment{proof}[1][\proofname]{\par
  \pushQED{\qed}\normalfont\topsep4pt\relax
  \trivlist\item[\hskip\labelsep\color{StudyBlue}\bfseries\boldmath #1\@addpunct{.}]
  \leavevmode\par\nobreak
}{\popQED\endtrivlist\@endpefalse}
\makeatother
\pdfstringdefDisableCommands{\def\studylabel#1{#1}\def\studyheading#1{#1}}


% Explicit navigation entries are declared beside each visible title, including
% inline headings. Editing a Contents caption never changes the study text.
\newcommand{\manualtoc}[3]{\addcontentsline{toc}{#1}{#2}\label{#3}}
\newcommand{\inlinetoc}[3]{\phantomsection\manualtoc{#1}{#2}{#3}}
\setcounter{tocdepth}{4}
\hypersetup{bookmarksdepth=4,linktoc=all}
\makeatletter
\renewcommand{\l@section}[2]{\addpenalty{\@secpenalty}\addvspace{5pt}%
  \begingroup\hypersetup{linkcolor=StudyRed}\bfseries\color{StudyRed}%
  \@dottedtocline{1}{0em}{0em}{#1}{#2}\endgroup}
\renewcommand{\l@subsection}[2]{\begingroup\bfseries\color{StudyBlue}%
  \@dottedtocline{2}{1em}{0em}{#1}{#2}\endgroup}
\renewcommand{\l@subsubsection}[2]{\begingroup\bfseries\color{StudyBlue}%
  \@dottedtocline{3}{2em}{0em}{#1}{#2}\endgroup}
\renewcommand{\l@paragraph}[2]{\begingroup\bfseries\color{StudyBlue}%
  \@dottedtocline{4}{3em}{0em}{#1}{#2}\endgroup}
\makeatother

\makeatletter
\def\toclevel@frontmatter{1}
\newcommand{\l@frontmatter}[2]{\begingroup\bfseries\color{StudyBlue}%
  \@dottedtocline{1}{0em}{0em}{#1}{#2}\endgroup}
\makeatother

\hypersetup{pdftitle={Probability for Finance - Studiabile through HW 53}}

\begin{document}
\begin{center}
{\Large\bfseries\color{StudyBlue} Probability for Finance}\\[4pt]
{\large\bfseries\color{StudyBlue} Studiabile --- Chapters 1--5, through HW 53}
\end{center}
% !TEX root = ../main.tex
% Full disclaimer explicitly requested by the user.
\begin{center}
{\large\studyheading{\inlinetoc{frontmatter}{Theory, Expanded Proofs, Worked Homework, Past and AI-Generated Exams, and More}{studiabile-nav-disclaimer}Theory, Expanded Proofs, Worked Homework, Past and AI-Generated Exams, and More}}
\end{center}
\noindent These unrevised notes were generated using artificial intelligence, following the student's preferred style and, where possible, the style of Professor Tiziano De Angelis and former tutor Andrea Bovo. Additional proofs and explanations not presented or proved by the professor have been included at the student's discretion. The material has not undergone a complete independent review and may contain errors.

\noindent These notes follow the lecturer's order and explain the hypotheses
and steps of each proof. Every homework retains its own statement and solution
beside the relevant theory. A later homework is anticipated only when the
lecturer's own proof takes that exact fact as given; its proof remains at the
homework itself. All mathematical displays are unnumbered.
%
\clearpage
\begingroup
\setlength{\parskip}{0pt}
\tableofcontents
\endgroup
\clearpage
\section*{1. Measure Spaces}
\manualtoc{section}{\texorpdfstring{1. Measure Spaces}{1. Measure Spaces}}{studiabile-nav-0001}
% !TEX root = ../main.tex
\subsection*{Measure theory and continuous functions}%
\manualtoc{subsection}{\texorpdfstring{Measure theory and continuous functions}{Measure theory and continuous functions}}{studiabile-nav-0002}%

A function $f$ from a topological set $X$ into another topological set $Y$ (denoted as $f:X\to Y$) is said to be continuous if the pre-images of open sets in $Y$ are open sets in $X$.
In formulae $f^{-1}(A)=\{x\in X:f(x)\in A\}$ is open in $X$ whenever $A$ is open in $Y$.

In Measure Theory the focus is on measurable functions.
A function $f:X\to Y$ is measurable if $f^{-1}(A)=\{x\in X:f(x)\in A\}$ is a measurable set in $X$ whenever $A$ is a measurable set in $Y$.
%
% !TEX root = ../main.tex
\subsection*{Def 1.1. Algebra}%
\manualtoc{subsection}{\texorpdfstring{Def 1.1. Algebra}{Def 1.1. Algebra}}{studiabile-nav-0003}%

Let $S$ be a set. A collection $\Sigma_0$ of subsets of $S$ is called an algebra if
\[
\begin{aligned}
\text{\studylabel{(i)}}\quad&S\in\Sigma_0,\\
\text{\studylabel{(ii)}}\quad&A\in\Sigma_0\Rightarrow A^c\in\Sigma_0,\\
\text{\studylabel{(iii)}}\quad&A,B\in\Sigma_0\Rightarrow A\cup B\in\Sigma_0.
\end{aligned}
\]
%
% !TEX root = ../main.tex
\subsection*{HW 1 --- Closure under finite intersections}%
\manualtoc{subsection}{\texorpdfstring{HW 1 --- Closure under finite intersections}{HW 1 --- Closure under finite intersections}}{studiabile-nav-0004}%

Show that $A,B\in\Sigma_0\Rightarrow A\cap B\in\Sigma_0$.

\par\noindent\studyheading{Proof.}\par\nobreak
Let $A,B\in\Sigma_0$. It is enough to show that $(A\cap B)^c\in\Sigma_0$ since $\Sigma_0$ is an algebra.
\[
(A\cap B)^c=\{x:x\notin(A\cap B)\}=\{x:x\notin A\text{ or }x\notin B\}=A^c\cup B^c.
\]
By the properties \studylabel{(ii)} and \studylabel{(iii)} of an algebra we have that $(A\cap B)^c\in\Sigma_0$ and so $(A\cap B)\in\Sigma_0$ by \studylabel{(ii)}.
%
% !TEX root = ../main.tex
\subsection*{HW 2 --- Finite unions and intersections}%
\manualtoc{subsection}{\texorpdfstring{HW 2 --- Finite unions and intersections}{HW 2 --- Finite unions and intersections}}{studiabile-nav-0006}%

Show that $(A_n)_{n=1}^N\subseteq\Sigma_0$ implies $\bigcup_{n=1}^N A_n\in\Sigma_0$ and $\bigcap_{n=1}^N A_n\in\Sigma_0$.

\par\noindent\studyheading{Proof.}\par\nobreak
By definition of algebra we know that $A,B\in\Sigma_0\Rightarrow A\cup B\in\Sigma_0$.
We prove the result by induction on $N$.
\studyheading{Base case:} $N=1$ is trivial.
\studyheading{Inductive hypothesis:} Assume that for some $N\ge1$, $\bigcup_{n=1}^N A_n\in\Sigma_0$.
\studyheading{Inductive step:} We want to prove that $\bigcup_{n=1}^{N+1}A_n\in\Sigma_0$.
Notice that
\[
\bigcup_{n=1}^{N+1}A_n=\left(\bigcup_{n=1}^{N}A_n\right)\cup A_{N+1}\in\Sigma_0.
\]
Therefore by inductive hypothesis and by assumption $A_{N+1}\in\Sigma_0$ the closure of an algebra under finite unions holds.
For finite intersection by De Morgan's law
\[
\bigcap_{n=1}^N A_n=\left(\bigcup_{n=1}^N A_n^c\right)^c.
\]
Hence by the closure under complements $A_n^c\in\Sigma_0$ for $n=1,\ldots,N$ and by the closure under finite unions also $\bigcup_{n=1}^N A_n^c\in\Sigma_0$ and therefore
\[
\left(\bigcup_{n=1}^N A_n^c\right)^c=\bigcap_{n=1}^N A_n\in\Sigma_0
\]
by closure under complements.
%
% !TEX root = ../main.tex
\subsection*{Def 1.2. $\sigma$-algebra}%
\manualtoc{subsection}{\texorpdfstring{Def 1.2. $\sigma$-algebra}{Def 1.2. sigma-algebra}}{studiabile-nav-0011}%

Let $S$ be a set. A collection $\Sigma$ of subsets of $S$ is called a $\sigma$-algebra if
\[
\begin{aligned}
\text{\studylabel{(i)}}\quad&S\in\Sigma,\\
\text{\studylabel{(ii)}}\quad&A\in\Sigma\Rightarrow A^c\in\Sigma,\\
\text{\studylabel{(iii)}}\quad&(A_n)_{n=1}^\infty\subseteq\Sigma\Rightarrow\bigcup_{n=1}^\infty A_n\in\Sigma.
\end{aligned}
\]
%
% !TEX root = ../main.tex
\subsection*{HW 3 --- Countable intersections}%
\manualtoc{subsection}{\texorpdfstring{HW 3 --- Countable intersections}{HW 3 --- Countable intersections}}{studiabile-nav-0012}%

Show that $(A_n)_{n=1}^\infty\subseteq\Sigma\Rightarrow\bigcap_{n=1}^\infty A_n\in\Sigma$.

\par\noindent\studyheading{Proof.}\par\nobreak
Let $(A_n)_{n\in\mathbb N}\subseteq\Sigma$. Since $\Sigma$ is a $\sigma$-algebra it is closed under complements. Therefore $A_n^c\in\Sigma$ $\forall n\in\mathbb N$. By closure under countable unions $\bigcup_{n=1}^\infty A_n^c\in\Sigma$ and therefore its complement $\left(\bigcup_{n=1}^\infty A_n^c\right)^c\in\Sigma$, which is $\bigcap_{n=1}^\infty A_n$ by De Morgan's law.
%
% !TEX root = ../main.tex
\subsection*{Measurable spaces, generated sigma-algebras and Borel sets}%
\manualtoc{subsection}{\texorpdfstring{Measurable spaces, generated sigma-algebras and Borel sets}{Measurable spaces, generated sigma-algebras and Borel sets}}{studiabile-nav-0014}%

In the probabilistic formulation of a random experiment we will use $\sigma$-algebras to describe all possible events. However, working with $\sigma$-algebras is difficult because it is often difficult to describe ALL sets in a $\sigma$-algebra. It is often easier to handle algebras and we will soon learn about $\pi$-systems.

\studyheading{\inlinetoc{subsubsection}{\texorpdfstring{Notation:}{Notation:}}{studiabile-nav-0015}Notation:} From now on $S$ will always denote a set, $\Sigma_0$ an algebra of subsets of $S$ and $\Sigma$ a $\sigma$-algebra of subsets of $S$.

\studyheading{\inlinetoc{subsubsection}{\texorpdfstring{Measurable sets:}{Measurable sets:}}{studiabile-nav-0016}Measurable sets:} any $A\in\Sigma$ is said to be a $\Sigma$-measurable set.

\studyheading{\inlinetoc{subsubsection}{\texorpdfstring{Measurable space:}{Measurable space:}}{studiabile-nav-0017}Measurable space:} a pair $(S,\Sigma)$ is called a measurable space.

Letting $\mathcal C$ be a class of subsets of $S$ we denote by $\sigma(\mathcal C)$ the smallest $\sigma$-algebra that contains all sets in $\mathcal C$. We call it the $\sigma$-algebra generated by $\mathcal C$. Notice that, in particular, if $\mathcal C\subseteq\Sigma$ then $\sigma(\mathcal C)\subseteq\Sigma$.

Given a topological space $X$ the Borel $\sigma$-algebra $\mathcal B(X)$ is the smallest $\sigma$-algebra that contains all the open sets of $X$. That is
\[
\mathcal B(X)=\sigma(\text{open sets in }X).
\]
In particular we will often work with $X=\mathbb R$ and $\mathcal B(\mathbb R)$ generated by open sets on the real line. In that context a very useful class of subsets is the following (example of $\pi$-system)
\[
\pi(\mathbb R)=\{(-\infty,x],\ x\in\mathbb R\}.
\]
Indeed it is easy to describe all the set contained in $\pi(\mathbb R)$ and the next proposition allow us to connect with $\mathcal B(\mathbb R)$.
%
% !TEX root = ../main.tex
\subsection*{Prop 1.1}%
\manualtoc{subsection}{\texorpdfstring{Prop 1.1}{Prop 1.1}}{studiabile-nav-0018}%

It holds $\sigma(\pi(\mathbb R))=\mathcal B(\mathbb R)$.

\par\noindent\studyheading{Proof.}\par\nobreak
To prove the proposition we will show that $\sigma(\pi(\mathbb R))\subseteq\mathcal B(\mathbb R)$ and $\sigma(\pi(\mathbb R))\supseteq\mathcal B(\mathbb R)$.
A generic set in $\pi(\mathbb R)$ is of the form $(-\infty,x]$ for some $x\in\mathbb R$ and it can be written as
\[
(-\infty,x]=\bigcap_{n=1}^\infty(-\infty,x+\tfrac1n)\in\mathcal B(\mathbb R),
\]
A point $y\le x$ belongs to every open ray on the right. Conversely, if $y>x$, choose $n$ with $1/n<y-x$; then $y\notin(-\infty,x+1/n)$. This proves the equality. Each ray is open, and the Borel sigma-algebra is closed under countable intersections, proving the final membership.
\textit{Every element of $\pi(\mathbb R)$ is also an element of $\mathcal B(\mathbb R)$.}
Therefore $\pi(\mathbb R)\subseteq\mathcal B(\mathbb R)$ which implies $\sigma(\pi(\mathbb R))\subseteq\mathcal B(\mathbb R)$.
\emph{For the reverse inclusion, we show that every open set belongs to $\sigma(\pi(\mathbb R))$.} We notice that any open set in $\mathbb R$ can be written as countable union of disjoint open intervals. Therefore it suffices to show that any open interval $(a,b)$, allowing infinite endpoints, belongs to $\sigma(\pi(\mathbb R))$.
If both endpoints are infinite, the interval is $\mathbb R$, which belongs to every $\sigma$-algebra on $\mathbb R$. For the remaining cases the following identities apply, omitting a redundant ray when an endpoint is infinite. Notice that
\[
\begin{aligned}
(a,b)&=(-\infty,b)\cap(a,\infty),\\
(a,+\infty)&=(-\infty,a]^c\in\sigma(\pi(\mathbb R)),
\end{aligned}
\]
because it is the complement of a set in $\sigma(\pi(\mathbb R))$.
\[
(-\infty,b)=\bigcup_{n=1}^\infty(-\infty,b-\tfrac1n]\in\sigma(\pi(\mathbb R))
\]
For this equality, every point in a ray on the right is strictly below $b$. If $y<b$, choose $n$ such that $1/n\le b-y$; then $y\le b-1/n$ and belongs to the union. Each closed ray belongs to the generating class, so countable-union closure proves the membership.
Since $(a,+\infty)$ and $(-\infty,b)\in\sigma(\pi(\mathbb R))$ then $(a,b)\in\sigma(\pi(\mathbb R))$.
For an arbitrary open set $A\subseteq\mathbb R$ we have a sequence $(a_n,b_n)_{n\in\mathbb N}$ such that
\[
A=\bigcup_{n=1}^\infty(a_n,b_n)\in\sigma(\pi(\mathbb R)),
\]
where the final membership is by countable union of elements of a $\sigma$-alg.
Therefore any open set in $\mathbb R$ belongs to $\sigma(\pi(\mathbb R))$ and thus we conclude that $\mathcal B(\mathbb R)\subseteq\sigma(\pi(\mathbb R))$.
%
% !TEX root = ../main.tex
\subsection*{Set functions}%
\manualtoc{subsection}{\texorpdfstring{Set functions}{Set functions}}{studiabile-nav-0020}%

A set function is a mapping from an algebra to the positive extended real line $\mu_0:\Sigma_0\to[0,+\infty]$.
%
% !TEX root = ../main.tex
\subsection*{Additive and countably additive set functions}%
\manualtoc{subsection}{\texorpdfstring{Additive and countably additive set functions}{Additive and countably additive set functions}}{studiabile-nav-0021}%

A set function $\mu_0:\Sigma_0\to[0,+\infty]$ is additive if $\mu_0(\emptyset)=0$ and
\[
\mu_0(A\cup B)=\mu_0(A)+\mu_0(B)
\]
whenever $A,B\in\Sigma_0$ are disjoint.
A set function is countably additive if $\mu_0(\emptyset)=0$ and for any $(A_n)_{n=1}^\infty\subseteq\Sigma_0$ such that $A_i\cap A_j=\emptyset$ $\forall i\ne j$ and $\bigcup_{n=1}^\infty A_n\in\Sigma_0$ it holds
\[
\mu_0\left(\bigcup_{n=1}^\infty A_n\right)=\sum_{n=1}^\infty\mu_0(A_n).
\]
%
% !TEX root = ../main.tex
\subsection*{HW 4 --- Finite additivity by induction}%
\manualtoc{subsection}{\texorpdfstring{HW 4 --- Finite additivity by induction}{HW 4 --- Finite additivity by induction}}{studiabile-nav-0022}%

Let $\mu_0$ be an additive set function. Show that for any $(A_n)_{n=1}^N\subseteq\Sigma_0$ such that $A_i\cap A_j=\emptyset$ for all $i\ne j$, it holds
\[
\mu_0\left(\bigcup_{n=1}^N A_n\right)=\sum_{n=1}^N\mu_0(A_n).
\]
\par\noindent\studyheading{Proof.}\par\nobreak
We prove the statement by induction on $N$.
\studyheading{Base case:} $N=1$ is trivial $\mu_0(A_1)=\mu_0(A_1)$.
\studyheading{Inductive hypothesis:} Assume the statement holds for some $N\ge1$.
\studyheading{Inductive step:} We want to prove that the statement holds for $N+1$.

Notice that
\[
\bigcup_{n=1}^{N+1}A_n=\bigcup_{n=1}^N A_n\cup A_{N+1}\in\Sigma_0
\]
since $\Sigma_0$ is an algebra (finite unions).
\[
\left(\bigcup_{n=1}^N A_n\right)\cap A_{N+1}=\emptyset
\]
since sets are pairwise disjoint.
By additivity of $\mu_0$
\[
\begin{aligned}
\mu_0\left(\bigcup_{n=1}^{N+1}A_n\right)
&=\mu_0\left(\bigcup_{n=1}^N A_n\cup A_{N+1}\right) =\mu_0\left(\bigcup_{n=1}^N A_n\right)+\mu_0(A_{N+1})\\&=\sum_{n=1}^N\mu_0(A_n)+\mu_0(A_{N+1})
=\sum_{n=1}^{N+1}\mu_0(A_n).
\end{aligned}
\]
The first equality uses the decomposition of the union, the second finite additivity for the two disjoint measurable sets, and the third the inductive hypothesis. The final equality is the definition of the finite sum. This completes the induction.
%
% !TEX root = ../main.tex
\subsection*{HW 5 --- Countable additivity implies additivity}%
\manualtoc{subsection}{\texorpdfstring{HW 5 --- Countable additivity implies additivity}{HW 5 --- Countable additivity implies additivity}}{studiabile-nav-0027}%

Show that a countably additive set function is also additive.

\par\noindent\studyheading{Proof.}\par\nobreak
Let $\mu_0:\Sigma_0\to[0,+\infty]$ a countably additive set function.
By definition of countable additivity, $\mu_0(\emptyset)=0$.
Let $A,B\in\Sigma_0$ be two disjoint sets, we define the sequence $(C_n)_{n\in\mathbb N}$ as follows:
$C_1=A$, $C_2=B$, $C_n=\emptyset$ $\forall n\ge3$.
Notice that $(C_n)_{n\in\mathbb N}$ are pairwise disjoint and $\bigcup_{n=1}^\infty C_n=A\cup B\in\Sigma_0$.
Therefore by countable additivity
\[
\begin{aligned}
\mu_0(A\cup B)&=\mu_0\left(\bigcup_{n=1}^\infty C_n\right)
=\sum_{n=1}^\infty\mu_0(C_n) =\mu_0(A)+\mu_0(B)+\sum_{n=3}^\infty\mu_0(\emptyset)
=\mu_0(A)+\mu_0(B).
\end{aligned}
\]
The first equality uses the union of the $C_n$, the second countable additivity, and the remaining equalities use their definitions and $\mu_0(\emptyset)=0$.
%
% !TEX root = ../main.tex
\subsection*{Measure}%
\manualtoc{subsection}{\texorpdfstring{Measure}{Measure}}{studiabile-nav-0029}%

Given a measurable space $(S,\Sigma)$ a function $\mu:\Sigma\to[0,+\infty]$ is called measure if it is countably additive.

\studyheading{\inlinetoc{subsubsection}{\texorpdfstring{Measure space:}{Measure space:}}{studiabile-nav-0030}Measure space:} The triple $(S,\Sigma,\mu)$ define a measure space.

A measure is finite if $\mu(S)<\infty$.

A measure is $\sigma$-finite if there exists $(A_n)_{n\in\mathbb N}\subseteq\Sigma$ such that $\bigcup_{n=1}^\infty A_n=S$ and $\mu(A_n)<\infty$ for each $n\in\mathbb N$.

\studyheading{\inlinetoc{subsubsection}{\texorpdfstring{Null set:}{Null set:}}{studiabile-nav-0031}Null set:} a set $A\in\Sigma$ is a null set if $\mu(A)=0$. Actually, to be formal the collection of null sets is
$\mathcal N=\{A\subseteq S:A\subseteq B\in\Sigma,\ \mu(B)=0\}$.
Notice that in the definition of $\mathcal N$ the sets $A$ are not necessarily measurable but they are subsets of measurable sets.

\studyheading{\inlinetoc{subsubsection}{\texorpdfstring{Complete measure space:}{Complete measure space:}}{studiabile-nav-0032}Complete measure space:} a measure space is complete if $\mathcal N\subseteq\Sigma$.

A statement $H(s)$ about points $s\in S$ is said to hold almost everywhere (almost surely in case of probability measures) if its failure set is contained in a measurable set of measure zero. In particular, when the failure set is measurable, this means $\mu(F)=0$, where
$F=\{s\in S:H(s)\text{ is false}\}$.
%
% !TEX root = ../main.tex
\subsection*{Definition 1.3. $\pi$-system}%
\manualtoc{subsection}{\texorpdfstring{Definition 1.3. $\pi$-system}{Definition 1.3. pi-system}}{studiabile-nav-0033}%

A collection $\mathcal I$ of subsets of $S$ such that
$I_1,I_2\in\mathcal I\Rightarrow I_1\cap I_2\in\mathcal I$ is called $\pi$-system.
%
% !TEX root = ../main.tex
\subsection*{HW 6 --- Lower rays form a pi-system}%
\manualtoc{subsection}{\texorpdfstring{HW 6 --- Lower rays form a pi-system}{HW 6 --- Lower rays form a pi-system}}{studiabile-nav-0034}%

Prove that $\pi(\mathbb R)$ is a $\pi$-system.
Recall the definition of $\pi(\mathbb R)=\{(-\infty,x],\ x\in\mathbb R\}$.

\par\noindent\studyheading{Proof.}\par\nobreak
By definition a collection $\mathcal I$ is a $\pi$-system if $A,B\in\mathcal I\Rightarrow A\cap B\in\mathcal I$.
Let $A,B\in\pi(\mathbb R)$ then there exists $x,y\in\mathbb R$ such that $A=(-\infty,x]$, $B=(-\infty,y]$.
\[
A\cap B=(-\infty,\min\{x,y\}]\in\pi(\mathbb R).
\]
Indeed, a real number belongs to both lower rays exactly when it is at most both endpoints, equivalently at most their minimum. The resulting ray is again an element of $\pi(\mathbb R)$.
%
% !TEX root = ../main.tex
\subsection*{Lemma 1.2}%
\manualtoc{subsection}{\texorpdfstring{Lemma 1.2}{Lemma 1.2}}{studiabile-nav-0036}%

Let $\Sigma=\sigma(\mathcal I)$ where $\mathcal I$ is a $\pi$-system on $S$. Suppose that $\mu_1$ and $\mu_2$ are measures on $(S,\Sigma)$ such that $\mu_1(S)=\mu_2(S)<\infty$ and $\mu_1(I)=\mu_2(I)$ for every $I\in\mathcal I$. Then $\mu_1(A)=\mu_2(A)$ $\forall A\in\Sigma$.
%
% !TEX root = ../main.tex
\subsection*{Theorem 1.3 (Carath\'eodory's Theorem)}%
\manualtoc{subsection}{\texorpdfstring{Theorem 1.3 (Carath\'eodory's Theorem)}{Theorem 1.3 (Carathéodory's Theorem)}}{studiabile-nav-0037}%

Let $S$ be a set, $\Sigma_0$ an algebra on $S$ and $\Sigma=\sigma(\Sigma_0)$. If a set function $\mu_0:\Sigma_0\to[0,+\infty]$ is countably additive, then there exists a measure $\mu:\Sigma\to[0,+\infty]$ such that $\mu(A)=\mu_0(A)$ for any $A\in\Sigma_0$. Moreover if $\mu_0(S)<\infty$ then $\mu$ is uniquely defined.
We say that $\mu$ in the theorem above is an extension of $\mu_0$ to the $\sigma$-algebra $\Sigma$.
%
% !TEX root = ../main.tex
\subsection*{Lebesgue measure}%
\manualtoc{subsection}{\texorpdfstring{Lebesgue measure}{Lebesgue measure}}{studiabile-nav-0038}%

Let $S=(0,1]$, $\Sigma=\mathcal B((0,1])$ and
\[
\Sigma_0=\left\{F:F=\bigcup_{i=1}^n(a_i,b_i],\ \text{for pairwise disjoint sets }(a_i,b_i]\subseteq(0,1]\text{ and some }n\in\mathbb N\right\}.
\]
Then $\sigma(\Sigma_0)=\Sigma$. Define $\mu_0$ for $F\in\Sigma_0$ as $\mu_0(F)=\sum_{i=1}^n(b_i-a_i)$.
$\mu_0$ is countably additive on $\Sigma_0$ and $\mu_0(S)=1$ $\Rightarrow$ $\exists$ a measure $\mu$ on $\Sigma$ that extends $\mu_0$. We call it Lebesgue measure, denoted $\mathrm{Leb}$.
An analogous construction defines Lebesgue measure on $\mathbb R$ and on $\mathbb R^d$. The countable-additivity step is taken as given here; HW~8 asks only for finite additivity.
%
% !TEX root = ../main.tex
\subsection*{HW 7 --- The interval algebra}%
\manualtoc{subsection}{\texorpdfstring{HW 7 --- The interval algebra}{HW 7 --- The interval algebra}}{studiabile-nav-0039}%

Let $S=(0,1]$ and let $\Sigma_0$ be the collection of finite unions of pairwise disjoint intervals $(a_i,b_i]\subseteq S$, with $0\le a_i\le b_i\le1$. Prove that $\Sigma_0$ is an algebra on $S$.

\par\noindent\studyheading{Proof.}\par\nobreak
First,
\[
\emptyset\in\Sigma_0,\qquad S=(0,1]\in\Sigma_0,
\]
since $(x,x]=\emptyset$ for every $x\in(0,1]$, while $S$ is itself an interval of the required form.

Let $A\in\Sigma_0$ be nonempty. Without loss of generality, we can assume that
\[
A=\bigcup_{i=1}^{n}(a_i,b_i],
\]
with
\[
0\le a_1<b_1\le a_2<b_2\le\cdots\le a_n<b_n\le1,
\]
since the nonempty intervals are pairwise disjoint and can be ordered by their left endpoints. We obtain
\[
A^c=(0,a_1]\cup
\left(\bigcup_{i=1}^{n-1}(b_i,a_{i+1}]\right)
\cup(b_n,1].
\]
These intervals are precisely the gaps between the intervals forming $A$, including possible gaps at the endpoints of $S$. Since they are finitely many and pairwise disjoint, $A^c\in\Sigma_0$. If $A=\emptyset$, then $A^c=S\in\Sigma_0$.

Finally, let $A,B\in\Sigma_0$, with
\[
A=\bigcup_{i=1}^{n}(a_i,b_i],\qquad
B=\bigcup_{j=1}^{m}(c_j,d_j].
\]
The intervals from the two collections may overlap. To obtain a finite partition, collect all their endpoints $a_i,b_i,c_j,d_j$, together with $0$ and $1$, remove repetitions and arrange them in increasing order:
\[
0=t_0<t_1<\cdots<t_r=1.
\]
Define
\[
D_\ell=(t_{\ell-1},t_\ell],\qquad \ell=1,\ldots,r.
\]
These intervals are pairwise disjoint and satisfy
\[
S=\bigcup_{\ell=1}^{r}D_\ell,
\]
so they form a finite partition of $S$.

The same construction applies to any finite collection of sets in $\Sigma_0$. Thus, in subsequent arguments, we may work without loss of generality with a common finite partition of $S$ into intervals of this form, with each original set represented as a union of its pieces.

Moreover, every original nonempty interval $(a_i,b_i]$ is a union of consecutive $D_\ell$. Indeed, since all its endpoints were selected, there exist $p<q$ such that $a_i=t_p$ and $b_i=t_q$. Hence,
\[
(a_i,b_i]=\bigcup_{\ell=p+1}^{q}D_\ell.
\]
The equality holds because consecutive half-open intervals include their right endpoints and exclude their left endpoints, so no points are lost or counted twice.

The same argument applies to every interval forming $B$. Consequently, $A\cup B$ is a finite union of some of the $D_\ell$, which are pairwise disjoint intervals of the required form. Therefore,
\[
A\cup B\in\Sigma_0.
\]
Thus, $\Sigma_0$ contains $S$ and is closed under complements and finite unions, so it is an algebra on $(0,1]$.

%
% !TEX root = ../main.tex
\subsection*{HW 8 --- Finite additivity of interval length}%
\manualtoc{subsection}{\texorpdfstring{HW 8 --- Finite additivity of interval length}{HW 8 --- Finite additivity of interval length}}{studiabile-nav-0041}%

On the interval algebra $\Sigma_0$ from HW~7, define $\mu_0(A)=\sum_{i=1}^n(b_i-a_i)$ for a disjoint representation $A=\bigcup_{i=1}^n(a_i,b_i]$. Prove that $\mu_0$ is well defined and finitely additive.

\par\noindent\studyheading{Proof.}\par\nobreak
First, for every $A\in\Sigma_0$ represented by pairwise disjoint intervals,
\[
A=\bigcup_{i=1}^{n}(a_i,b_i],\qquad
\mu_0(A)=\sum_{i=1}^{n}(b_i-a_i)\ge0.
\]
The inequality follows because $b_i\ge a_i$ for every $i$. Moreover,
\[
\mu_0(\emptyset)=\mu_0((x,x])=x-x=0.
\]

\emph{We also verify that $\mu_0$ does not depend on the chosen disjoint interval representation.} Given two representations of $A$, their common endpoint partition, constructed as in HW~7, yields
\[
\sum_{i=1}^{n}(b_i-a_i)
=\sum_{\ell:\,D_\ell\subseteq A}(t_\ell-t_{\ell-1})
=\sum_{j=1}^{m}(d_j-c_j).
\]
Both equalities follow by splitting the intervals of each representation into consecutive partition intervals and adding their lengths.

Finally, let $A,B\in\Sigma_0$ be disjoint, with
\[
A=\bigcup_{i=1}^{n}(a_i,b_i],\qquad
B=\bigcup_{j=1}^{m}(c_j,d_j].
\]
Since $A\cap B=\emptyset$, every interval of $A$ is disjoint from every interval of $B$. Define
\[
e_k=\begin{cases}
a_k,&k=1,\ldots,n,\\
c_{k-n},&k=n+1,\ldots,n+m,
\end{cases}
\qquad
f_k=\begin{cases}
b_k,&k=1,\ldots,n,\\
d_{k-n},&k=n+1,\ldots,n+m.
\end{cases}
\]
Then
\[
A\cup B=\bigcup_{k=1}^{n+m}(e_k,f_k],
\]
where all intervals are pairwise disjoint. Hence,
\[
\begin{aligned}
\mu_0(A\cup B)
&=\sum_{k=1}^{n+m}(f_k-e_k) =\sum_{i=1}^{n}(b_i-a_i)+\sum_{j=1}^{m}(d_j-c_j) =\mu_0(A)+\mu_0(B).
\end{aligned}
\]
The first equality uses the disjoint representation just constructed; the second separates the first $n$ and last $m$ terms; the last follows from the definition of $\mu_0$. Repeating the argument proves additivity for any finite family of pairwise disjoint sets.

%
% !TEX root = ../main.tex
\subsection*{HW 9 --- Borel singletons}%
\manualtoc{subsection}{\texorpdfstring{HW 9 --- Borel singletons}{HW 9 --- Borel singletons}}{studiabile-nav-0043}%

Prove that $\{x\}\in\mathcal B(\mathbb R)$ $\forall x\in\mathbb R$.

\par\noindent\studyheading{Proof.}\par\nobreak
The Borel sigma algebra contains all open intervals and is closed under countable intersections.
Let $x\in\mathbb R$, we have that
$\{x\}=\bigcap_{n=1}^\infty(x-\frac1n,x+\frac1n)$.
We know that for all $n\in\mathbb N$, $(x-\frac1n,x+\frac1n)\in\mathcal B(\mathbb R)$ and countable intersection of measurable sets is a measurable set.
For the set equality, $x$ belongs to every interval. Conversely, if $y\ne x$, choose $n$ such that $1/n<|y-x|$; then $y$ is outside the corresponding interval. Thus the intersection is exactly $\{x\}$.
%
% !TEX root = ../main.tex
\subsection*{HW 10 --- Lebesgue measure of a singleton}%
\manualtoc{subsection}{\texorpdfstring{HW 10 --- Lebesgue measure of a singleton}{HW 10 --- Lebesgue measure of a singleton}}{studiabile-nav-0045}%

Prove that $\mathrm{Leb}(\{x\})=0$ $\forall x\in\mathbb R$.

\par\noindent\studyheading{Proof.}\par\nobreak
Let $x\in\mathbb R$. From HW~9 we know that $\{x\}\in\mathcal B(\mathbb R)$.
Notice that for every $n\in\mathbb N$, $\{x\}\subseteq(x-\frac1n,x]$ and $(x-\frac1n,x]=(x-\frac1n,x)\cup\{x\}$ where the two sets are disjoint and belong to $\mathcal B(\mathbb R)$.
By additivity of measure we have
\[
\frac1n=\mathrm{Leb}((x-\tfrac1n,x])
=\mathrm{Leb}((x-\tfrac1n,x))+\mathrm{Leb}(\{x\})\ge\mathrm{Leb}(\{x\}),
\]
where the last inequality follows from the non-negativity of measures.
Therefore $0\le\mathrm{Leb}(\{x\})\le\frac1n$.
Since this inequality holds for every $n\in\mathbb N$ by letting $n\to+\infty$ and applying the squeeze thm we obtain $\mathrm{Leb}(\{x\})=0$.
Moreover if a singleton had measure $a>0$, choosing $n>\frac1a$ would contradict $a\le\frac1n$ hence its measure is zero.
%
% !TEX root = ../main.tex
\subsection*{Lemma 1.4}%
\manualtoc{subsection}{\texorpdfstring{Lemma 1.4}{Lemma 1.4}}{studiabile-nav-0047}%

Let $(S,\Sigma,\mu)$ be a measure space. Then
\[
\begin{aligned}
\text{(a)}\quad&\mu(A\cup B)\le\mu(A)+\mu(B)\qquad\forall A,B\in\Sigma,\\
\text{(b)}\quad&\mu\left(\bigcup_{i=1}^n A_i\right)\le\sum_{i=1}^n\mu(A_i)\qquad\forall(A_i)_{i=1}^n\subseteq\Sigma,\\
\text{(c)}\quad&\text{If }\mu(S)<\infty\text{ then }\mu(A\cup B)=\mu(A)+\mu(B)-\mu(A\cap B)\quad\forall A,B\in\Sigma.
\end{aligned}
\]
%
% !TEX root = ../main.tex
\subsection*{HW 11 --- Subadditivity and the union formula}%
\manualtoc{subsection}{\texorpdfstring{HW 11 --- Subadditivity and the union formula}{HW 11 --- Subadditivity and the union formula}}{studiabile-nav-0048}%

Prove all three assertions of Lemma~1.4: binary subadditivity, finite subadditivity, and the union formula under $\mu(S)<\infty$.
\par\noindent\studyheading{Proof.}\par\nobreak
\studylabel{(a)} Let $A,B\in\Sigma$, we have that $A\setminus B=A\cap B^c\in\Sigma$, $A\cup B=(A\setminus B)\cup B$ and $(A\setminus B)\cap B=\emptyset$. It follows that
\[
\begin{aligned}
\mu(A\cup B)&=\mu((A\setminus B)\cup B)=\mu(A\setminus B)+\mu(B) \le\mu(A\setminus B)+\mu(A\cap B)+\mu(B)\\&=\mu((A\setminus B)\cup(A\cap B))+\mu(B)=\mu(A)+\mu(B).
\end{aligned}
\]
The first two equalities use the disjoint decomposition of $A\cup B$, the inequality uses $\mu(A\cap B)\ge0$, and the last two equalities use the disjoint decomposition of $A$.

\studylabel{(b)} Finite sub-additivity. We prove the statement by induction on $n$.
\studyheading{Base case:} $n=1$. Statement is trivial $\mu(A_1)=\mu(A_1)$.
\studyheading{Inductive hypothesis:} $n\ge1$, $\mu(\bigcup_{i=1}^n A_i)\le\sum_{i=1}^n\mu(A_i)$.
\studyheading{Inductive step:} $n+1$. Notice that $\bigcup_{i=1}^{n+1}A_i=(\bigcup_{i=1}^n A_i)\cup A_{n+1}$.
By prop (a) we have
\[
\begin{aligned}
\mu\left(\bigcup_{i=1}^{n+1}A_i\right)
&\le\mu\left(\bigcup_{i=1}^n A_i\right)+\mu(A_{n+1}) \le\sum_{i=1}^n\mu(A_i)+\mu(A_{n+1})
=\sum_{i=1}^{n+1}\mu(A_i),
\end{aligned}
\]
where the first inequality uses part (a), the second the inductive hypothesis, and the equality the definition of the finite sum.

\studylabel{(c)} $\mu(A\cup B)=\mu(A\setminus B)+\mu(B)$, since $A\setminus B$ and $B$ are disjoint.
$\mu(A)=\mu(A\setminus B)+\mu(A\cap B)$ since $A\setminus B$ and $A\cap B$ are disjoint.
Adding $\mu(B)$ we obtain
\[
\begin{aligned}
\mu(A)+\mu(B)&=\mu(A\setminus B)+\mu(A\cap B)+\mu(B) =\mu(A\cup B)+\mu(A\cap B),
\end{aligned}
\]
where the first equality uses the disjoint decomposition of $A$, and the second uses the disjoint decomposition of $A\cup B$ from part (a).
Since $\mu(S)<\infty$ for every $E\in\Sigma$, $\mu(E)<\infty$, in fact $\mu(S)=\mu(E)+\mu(E^c)$.
Therefore subtracting $\mu(A\cap B)$ from from both sidess we obtain
\[
\mu(A)+\mu(B)-\mu(A\cap B)=\mu(A\cup B).
\]
%
% !TEX root = ../main.tex
\subsection*{HW 12 --- Monotonicity of measures}%
\manualtoc{subsection}{\texorpdfstring{HW 12 --- Monotonicity of measures}{HW 12 --- Monotonicity of measures}}{studiabile-nav-0056}%

Prove that if $A,B\in\Sigma$, $A\subseteq B$ then $\mu(A)\le\mu(B)$.

\par\noindent\studyheading{Proof.}\par\nobreak
Notice that $B=(B\cap A)\cup(B\cap A^c)$ and that $(B\cap A)\cap(B\cap A^c)=\emptyset$. It follows
\[
\begin{aligned}
\mu(B)&=\mu((B\cap A)\cup(B\cap A^c))
=\mu(B\cap A)+\mu(B\cap A^c) \ge\mu(B\cap A)=\mu(A),
\end{aligned}
\]
because $A\subseteq B$.
The second equality uses additivity on the two disjoint measurable pieces, the inequality their nonnegative measures, and the last equality $A\cap B=A$.
%
% !TEX root = ../main.tex
\subsection*{HW 13 --- Measure of a difference}%
\manualtoc{subsection}{\texorpdfstring{HW 13 --- Measure of a difference}{HW 13 --- Measure of a difference}}{studiabile-nav-0058}%


% BEGIN PFF RESTORATION HW13_DIFFERENCE_DEFINITION
The operation of ``subtracting'' two sets is denoted by $A\setminus B$ and it means $A\cap B^c$.
% END PFF RESTORATION HW13_DIFFERENCE_DEFINITION
Show that for a finite measure $\mu$, if $B\subseteq A$ then $\mu(A\setminus B)=\mu(A)-\mu(B)$.

\par\noindent\studyheading{Proof.}\par\nobreak
We obtain
\[
\begin{aligned}
\mu(A)&=\mu((A\cap B)\cup(A\cap B^c)) =\mu(A\cap B)+\mu(A\cap B^c)=\mu(A\cap B)+\mu(A\setminus B).
\end{aligned}
\]
Knowing that $A\cap B=B$ it follows $\mu(A)=\mu(B)+\mu(A\setminus B)$ and so $\mu(A\setminus B)=\mu(A)-\mu(B)$ if $\mu(B)<\infty$.
The decomposition is disjoint and measurable. Finiteness of the measure ensures that the subtraction is between finite quantities; no $\infty-\infty$ expression is used.
%
% !TEX root = ../main.tex
\subsection*{Lemma 1.5. Continuity of measure from below}%
\manualtoc{subsection}{\texorpdfstring{Lemma 1.5. Continuity of measure from below}{Lemma 1.5. Continuity of measure from below}}{studiabile-nav-0060}%

Let $(S,\Sigma,\mu)$ be a measure space. If $(A_n)_{n\in\mathbb N}\subseteq\Sigma$ such that $A_n\subseteq A_{n+1}$ and $A_n\uparrow A$ as $n\to+\infty$ (in the sense that $A=\bigcup_{n=1}^\infty A_n$). Then $\mu(A_n)\uparrow\mu(A)$ as $n\to+\infty$.

\par\noindent\studyheading{Proof.}\par\nobreak
We relabel sets to make them disjoint with: $B_1=A_1$, $B_2=A_2\setminus B_1$, $\ldots$, $B_n=A_n\setminus\bigcup_{i=1}^{n-1}B_i$. Each $B_n$ is measurable, being a difference of measurable sets. As we will see in HW~14, the sets $(B_n)_{n\in\mathbb N}$ are pairwise disjoint. As we will see in HW~15, $\bigcup_{i=1}^nB_i=A_n$ and $\bigcup_{i=1}^\infty B_i=A$.
We obtain
\[
\begin{aligned}
\mu(A_n)&=\mu\left(\bigcup_{i=1}^n B_i\right)=\sum_{i=1}^n\mu(B_i) \le\sum_{i=1}^{n+1}\mu(B_i)=\mu\left(\bigcup_{i=1}^{n+1}B_i\right)=\mu(A_{n+1}).
\end{aligned}
\]

Here the union identities are supplied by HW~15, the finite-sum equalities use additivity on the disjoint sets from HW~14, and the inequality uses $\mu(B_{n+1})\ge0$.
Therefore the sequence $\mu(A_n)$, $n\in\mathbb N$ is non decreasing hence it admits a limit which we can now compute
\[
\lim_{n\to+\infty}\mu(A_n)=\lim_{n\to+\infty}\sum_{i=1}^n\mu(B_i)
=\sum_{i=1}^\infty\mu(B_i)=\mu\left(\bigcup_{i=1}^\infty B_i\right)=\mu(A),
\]
The first equality substitutes the finite-sum identity, the second is the definition of a nonnegative series, the third countable additivity for the disjoint measurable $B_i$, and the last $\bigcup_iB_i=A$ from HW~15.
%
% !TEX root = ../main.tex
\subsection*{HW 14 --- Disjoint pieces}%
\manualtoc{subsection}{\texorpdfstring{HW 14 --- Disjoint pieces}{HW 14 --- Disjoint pieces}}{studiabile-nav-0062}%

Prove that $B_i\cap B_j=\emptyset$ $\forall i\ne j$ where $B_1=A_1$, $B_n=A_n\setminus\bigcup_{k=1}^{n-1}B_k$, $n\ge2$.
\par\noindent\studyheading{Proof.}\par\nobreak
We prove by induction on $n$ that $B_1,\ldots,B_n$ are pairwise disjoint.
\studyheading{Base case:} $n=1$ is trivial.
\studyheading{Inductive hypothesis:} $n\ge1$. $B_1,\ldots,B_n$ are pairwise disjoint i.e. $B_i\cap B_j=\emptyset$ $\forall1\le i<j\le n$.
\studyheading{Inductive step:} By definition $B_{n+1}=A_{n+1}\setminus\bigcup_{k=1}^n B_k$.
Therefore $B_{n+1}\cap(\bigcup_{k=1}^n B_k)=\emptyset$.
Since $B_i\subseteq\bigcup_{k=1}^n B_k$ $\forall i=1,\ldots,n$ we have that $B_{n+1}\cap B_i=\emptyset$ $\forall i=1,\ldots,n$.
Moreover by inductive hyp $B_1,\ldots,B_n$ are already pairwise disjoint hence $B_1,\ldots,B_{n+1}$ are pairwise disjoint.
%
% !TEX root = ../main.tex
\subsection*{HW 15 --- Unions of the disjoint pieces}%
\manualtoc{subsection}{\texorpdfstring{HW 15 --- Unions of the disjoint pieces}{HW 15 --- Unions of the disjoint pieces}}{studiabile-nav-0067}%

Prove that $\bigcup_{i=1}^n B_i=A_n$ \studylabel{(i)} and $\bigcup_{i=1}^\infty B_i=A$.

\par\noindent\studyheading{Proof.}\par\nobreak
Recall that $(A_n)_{n\in\mathbb N}$ is non decreasing sequence of sets such that $A_n\subseteq A_{n+1}$; $A=\bigcup_{i=1}^\infty A_n$.
We define $B_1=A_1$, $B_n=A_n\setminus\bigcup_{j=1}^{n-1}B_j$, $n\ge2$.
We prove \studylabel{(i)} by induction: $N=1$ is trivial.
\studyheading{Inductive hypothesis:} assume that statement hold for $n\ge1$, i.e. $\bigcup_{i=1}^n B_i=A_n$.
\studyheading{Inductive step:} $B_{n+1}=A_{n+1}\setminus\bigcup_{i=1}^n B_i$, by inductive hyp we obtain $B_{n+1}=A_{n+1}\setminus A_n$. Therefore
\[
\bigcup_{i=1}^{n+1}B_i=\left(\bigcup_{i=1}^n B_i\right)\cup B_{n+1}
=A_n\cup(A_{n+1}\setminus A_n)=A_{n+1},
\]
where the last equality follows from $A_n\subseteq A_{n+1}$.
Therefore by induction $\bigcup_{i=1}^n B_i=A_n$ $\forall n\in\mathbb N$ \studylabel{(i)}.
For \studylabel{(ii)} we know that $A_n\uparrow A$ and $A=\bigcup_{i=1}^\infty A_n$ by ass.
Therefore
\[
\bigcup_{i=1}^\infty B_i=\bigcup_{n=1}^\infty\left(\bigcup_{i=1}^n B_i\right)
=\bigcup_{n=1}^\infty A_n=A,
\]
where the first equality follows the fact that $B_i$ belongs to all finite unions with $n\ge i$ and the second follow the first part of HW~15 \studylabel{(i)}. The last follow the assumption $A_n\uparrow A$.
%
% !TEX root = ../main.tex
\subsection*{Lemma 1.6. Continuity of measure from above}%
\manualtoc{subsection}{\texorpdfstring{Lemma 1.6. Continuity of measure from above}{Lemma 1.6. Continuity of measure from above}}{studiabile-nav-0071}%

Let $(S,\Sigma,\mu)$ be a measure space. If $(A_n)_{n\in\mathbb N}\subseteq\Sigma$ such that $A_n\supseteq A_{n+1}$ and $A_n\downarrow A$ as $n\to+\infty$ (in the sense that $A=\bigcap_{n=1}^\infty A_n$). Then $\mu(A_n)\downarrow\mu(A)$ as $n\to+\infty$ provided that $\mu(A_m)<\infty$ for some $m\in\mathbb N$.
%
% !TEX root = ../main.tex
\subsection*{HW 16 --- Continuity from above}%
\manualtoc{subsection}{\texorpdfstring{HW 16 --- Continuity from above}{HW 16 --- Continuity from above}}{studiabile-nav-0072}%


% BEGIN PFF RESTORATION HW16_STATEMENT_HINT
\studyheading{Statement.} Prove Lemma~1.6.
\studyheading{\inlinetoc{subsubsection}{\texorpdfstring{Hint:}{Hint:}}{studiabile-nav-0074}Hint:} take $B_n=A_k\setminus A_{k+n}$ so that
$B_n\uparrow A_k\setminus A$ and use HW~13.
% END PFF RESTORATION HW16_STATEMENT_HINT
\par\noindent\studyheading{Proof.}\par\nobreak
Let $A_m$ such that $\mu(A_m)<\infty$ and define a sequence $B_k=A_m\setminus A_{m+k}$.
Since $A_{n+1}\subseteq A_n$ $\forall n\in\mathbb N$ we have $B_{k+1}\supseteq B_k$ $\forall k\in\mathbb N$.
Moreover we have that $B_k\uparrow B$ with $B=A_m\setminus A=A_m\setminus\bigcap_{n=1}^\infty A_n$.
By Lemma 1.5 we have
\[
\lim_{k\to+\infty}\mu(B_k)=\mu(B)=\mu(A_m\setminus A)=\mu(A_m)-\mu(A).
\]

By definition we have
\[
\lim_{k\to+\infty}\mu(B_k)
=\lim_{k\to+\infty}\mu(A_m\setminus A_{m+k})
=\mu(A_m)-\lim_{k\to+\infty}\mu(A_{m+k}).
\]
Combining these two equations we obtain
$\mu(A_m)-\mu(A)=\mu(A_m)-\lim_{k\to+\infty}\mu(A_{m+k})$.

We need the condition $\mu(A_n)<\infty$, if we don't impose that we would have this counter-example:
Take $S=[0,+\infty)$, $\Sigma=\mathcal B([0,+\infty))$ and $\mu=\mathrm{Leb}$.
Take $A_n=[n,+\infty)$ so that $A_n\downarrow\emptyset$, notice that $\mu(A_n)=\infty$ $\forall n$ and so the lemma fails.
All subtractions are legitimate because $A\subseteq A_{m+k}\subseteq A_m$ and $\mu(A_m)<\infty$, so the difference formula from HW~13 applies. Cancelling the common finite quantity gives $\lim_{k\to\infty}\mu(A_{m+k})=\mu(A)$. Omitting the finitely many initial indices does not change the limit; monotonicity of $A_n$ and HW~12 show that $\mu(A_n)$ decreases to this value.
%
% !TEX root = ../main.tex
\subsection*{Lemma 1.7}%
\manualtoc{subsection}{\texorpdfstring{Lemma 1.7}{Lemma 1.7}}{studiabile-nav-0076}%

Countable union of null set is a null set.
\par\noindent\studyheading{Proof.}\par\nobreak
Let $(N_n)_{n\in\mathbb N}\subseteq\Sigma$ be such that $\mu(N_n)=0$ for all $n$'s.
Set $\bar N=\bigcup_{n=1}^\infty N_n$ and $\bar N^k=\bigcup_{n=1}^k N_n$. Then by Lemma 1.5 and sub additivity of the measure
\[
\mu(\bar N)=\lim_{k\to+\infty}\mu(\bar N^k)
\le\lim_{k\to+\infty}\sum_{n=1}^k\mu(N_n)=0.
\]
Here $\bar N^k\uparrow\bar N$, so the first equality uses continuity from below. The inequality uses finite subadditivity at each $k$, and the last equality uses that every summand is zero. Nonnegativity then gives $\mu(\bar N)=0$.

For null sets $N_n$ that are not assumed measurable, choose $B_n\in\Sigma$ such that $N_n\subseteq B_n$ and $\mu(B_n)=0$, as allowed by the definition of a null set. We obtain
\[
\bigcup_{n=1}^{\infty}N_n\subseteq\bigcup_{n=1}^{\infty}B_n,
\qquad
\mu\left(\bigcup_{n=1}^{\infty}B_n\right)=0.
\]
The inclusion follows from $N_n\subseteq B_n$ for every $n$, and the equality follows from the measurable case just proved. Thus $\bigcup_nN_n$ is a null set in the general sense, whether or not it is measurable; no measure is assigned to a non-measurable set.
%
% !TEX root = ../main.tex
\subsection*{HW 17 --- Rational points form a null set}%
\manualtoc{subsection}{\texorpdfstring{HW 17 --- Rational points form a null set}{HW 17 --- Rational points form a null set}}{studiabile-nav-0078}%

Take the measure space $([0,1],\mathcal B([0,1]),\mathrm{Leb})$.
Prove that $[0,1]\cap\mathbb Q\in\mathcal B([0,1])$ and $\mathrm{Leb}([0,1]\cap\mathbb Q)=0$.
\par\noindent\studyheading{Proof.}\par\nobreak
Recall that $\mathbb Q$ is countable and any singleton $\in\mathcal B(\mathbb R)$.
We can write $[0,1]\cap\mathbb Q=\{q_1,q_2,\ldots\}=\bigcup_{n=1}^\infty\{q_n\}$ where $q_n\in[0,1]\cap\mathbb Q$.
Since $\{q_n\}\in\mathcal B([0,1])$ and $\mathcal B([0,1])$ is closed under countable unions then $[0,1]\cap\mathbb Q=\bigcup_{n=1}^\infty\{q_n\}\in\mathcal B([0,1])$.
By HW~10 we know that $\mathrm{Leb}(\{q_n\})=0$ $\forall n\in\mathbb N$, moreover since $q_n$'s are distinct are pairwise disjoint. Therefore by countable additivity
\[
\mathrm{Leb}([0,1]\cap\mathbb Q)
=\mathrm{Leb}\left(\bigcup_{n=1}^\infty\{q_n\}\right)
=\sum_{n=1}^\infty\mathrm{Leb}(\{q_n\})=0.
\]
%
\clearpage
%
\section*{2. Events}
\manualtoc{section}{\texorpdfstring{2. Events}{2. Events}}{studiabile-nav-0080}
% !TEX root = ../main.tex
\subsection*{Probability spaces}%
\manualtoc{subsection}{\texorpdfstring{Probability spaces}{Probability spaces}}{studiabile-nav-0081}%

A mathematical model for an experiment with uncertain (random) outcome is given by probability space $(\Omega,\mathcal F,\mathbb P)$ where

$\Omega$ is the sample space (the set $S$ from the previous section). A sample point is a point in $\Omega$, $\omega\in\Omega$ and it represents (loosely speaking) a possible scenario in the experiment.

$\mathcal F$ is the collection of all events (the $\sigma$-algebra $\Sigma$ from the previous section). An event $E\in\mathcal F$ is a collection of sample points, i.e. is a collection of all possible scenarios that may lead to the observed outcome in the experiment.

$\mathbb P$ is the probability measure that assigns to each event a probability of occurring. The probability $\mathbb P$ is a measure on $\mathcal F$ with $\mathbb P(\Omega)=1$. (Notice that $\mathbb P:\mathcal F\to[0,1]$.)
%
% !TEX root = ../main.tex
\subsection*{Example 1 (Simple coin tossing models)}%
\manualtoc{subsection}{\texorpdfstring{Example 1 (Simple coin tossing models)}{Example 1 (Simple coin tossing models)}}{studiabile-nav-0082}%

\studylabel{(a)} Toss a coin twice:
$\Omega=\{HH,HT,TH,TT\}$ and $\mathcal F=\mathcal P(\Omega)$, the collection of all unions of singleton outcomes, closed under intersections and complements.

\studylabel{(b)} Toss a coin infinitely many times:
$\Omega=\{H,T\}^{\mathbb N}$ $\Rightarrow$ $\omega\in\Omega$ is a sequence $\omega=(\omega_1,\omega_2,\omega_3,\ldots)$ with $\omega_i=H$ or $T$.
In principle one could consider $\mathcal F=\mathcal P(\Omega)$ where $\mathcal P(\Omega)$ is the collection of ALL subsets of $\Omega$ and it is called the power set. It turns out that this is not necessary and we can take the easier $\sigma$-algebra
\[
\mathcal F=\sigma(\{\omega\in\Omega:\omega_n=W\},\ n\in\mathbb N,\ W\in\{H,T\}).
\]
To convince ourselves that this is a good choice we will work towards showing that the set
\[
F:=\left\{\omega\in\Omega:\lim_{n\to+\infty}\frac{\#(k\le n:\omega_k=H)}{n}=\frac12\right\}
\]
belongs to $\mathcal F$ and, under the independent fair-coin probability model, $\mathbb P(F)=1$. This is telling us that indeed the Law of Large Numbers (LLN) holds as it should in any sensible mathematical model for random experiments!

Now we prove a result that is in some sense a converse of Lemma 1.7. It tells us that countable intersection of sure events give us again a sure event.
We need a small result first.
%
% !TEX root = ../main.tex
\subsection*{HW 18 --- Sure events and null complements}%
\manualtoc{subsection}{\texorpdfstring{HW 18 --- Sure events and null complements}{HW 18 --- Sure events and null complements}}{studiabile-nav-0085}%

For any $A\in\mathcal F$ show that $\mathbb P(A)=1\Longleftrightarrow\mathbb P(A^c)=0$.

\par\noindent\studyheading{Proof.}\par\nobreak
Let $A\in\mathcal F$, $A^c\in\mathcal F$. Notice that $A\cup A^c=\Omega$ and $A\cap A^c=\emptyset$.
We obtain $1=\mathbb P(\Omega)=\mathbb P(A)+\mathbb P(A^c)$, where the first equality is normalization and the second additivity on the two disjoint complementary events.
\noindent\studyheading{($\Rightarrow$)} If we assume that $\mathbb P(A)=1$, then $\mathbb P(A^c)=0$.

\noindent\studyheading{($\Leftarrow$)} If we assume that $\mathbb P(A^c)=0$ then $\mathbb P(A)$ has to be $1$.
%
% !TEX root = ../main.tex
\subsection*{Proposition 2.1}%
\manualtoc{subsection}{\texorpdfstring{Proposition 2.1}{Proposition 2.1}}{studiabile-nav-0089}%

If $(A_n)_{n\in\mathbb N}\subseteq\mathcal F$ are such that $\mathbb P(A_n)=1$ for all $n\in\mathbb N$, then $\mathbb P(\bigcap_{n=1}^\infty A_n)=1$.
\par\noindent\studyheading{Proof.}\par\nobreak
Since $A_n\in\mathcal F\Rightarrow A_n^c\in\mathcal F$ then $(\bigcap_{n=1}^\infty A_n)^c=\bigcup_{n=1}^\infty A_n^c\in\mathcal F$ and
\[
\mathbb P\left(\bigcap_{n=1}^\infty A_n\right)=1
\Longleftrightarrow
\mathbb P\left(\bigcup_{n=1}^\infty A_n^c\right)=0.
\]
Since $\mathbb P(A_n)=1\Longleftrightarrow\mathbb P(A_n^c)=0$ then $\mathbb P(\bigcup_{n=1}^\infty A_n^c)=0$ by Lemma 1.7 which concludes the proof.
%
% !TEX root = ../main.tex
\subsection*{Def lim inf and lim sup}%
\manualtoc{subsection}{\texorpdfstring{Def lim inf and lim sup}{Def lim inf and lim sup}}{studiabile-nav-0091}%

Given a sequence $(x_n)_{n\in\mathbb N}\subseteq\mathbb R$ we define
\[
\liminf_{n\to+\infty}x_n=\sup_{n\ge1}\inf_{m\ge n}x_m
\]
(i.e. the largest minorant in the limit),
\[
\limsup_{n\to+\infty}x_n=\inf_{n\ge1}\sup_{m\ge n}x_m
\]
(i.e. the smallest majorant in the limit).
Notice that the sequence $y_n:=\sup_{m\ge n}x_m$ is decreasing as
$y_n=\sup_{m\ge n}x_m\ge\sup_{m\ge n+1}x_m=y_{n+1}$.
Then $\inf_{n\ge1}y_n=\lim_{n\to+\infty}y_n$ and
\[
\limsup_{n\to+\infty}x_n=\inf_{n\ge1}\sup_{m\ge n}x_m
=\lim_{n\to+\infty}\sup_{m\ge n}x_m.
\]
%
% !TEX root = ../main.tex
\subsection*{HW 19 --- Tail infima}%
\manualtoc{subsection}{\texorpdfstring{HW 19 --- Tail infima}{HW 19 --- Tail infima}}{studiabile-nav-0092}%

Show that $\liminf_{n\to+\infty}x_n=\lim_{n\to+\infty}\inf_{m\ge n}x_m$.
\par\noindent\studyheading{Proof.}\par\nobreak
Let $y_n=\inf_{m\ge n}x_m$.
Notice that $\{x_m:m\ge n+1\}\subseteq\{x_m:m\ge n\}$.
Therefore $\inf_{m\ge n}x_m\le\inf_{m\ge n+1}x_m$ which implies $y_n\le y_{n+1}$.
Hence $(y_n)_{n\in\mathbb N}$ is a non-decreasing sequence, by monotonicity the limit existss in the extended real line and $\lim_{n\to+\infty}y_n=\sup_{n\ge1}y_n$. Substituting the definition of $y_n$ we conclude the proof.
%
% !TEX root = ../main.tex
\subsection*{Definition of convergence through liminf and limsup}%
\manualtoc{subsection}{\texorpdfstring{Definition of convergence through liminf and limsup}{Definition of convergence through liminf and limsup}}{studiabile-nav-0094}%

Given $(x_n)_{n\in\mathbb N}\subseteq\mathbb R$, the values of $\limsup x_n$ and $\liminf x_n$ are always well-defined (not necessarily finite) whereas the value of $\lim x_n$ may not exist. We say that a limit exist if
\[
\liminf_{n\to+\infty}x_n=\limsup_{n\to+\infty}x_n=:\lim_{n\to+\infty}x_n.
\]
The common value may belong to $\overline{\mathbb R}=[-\infty,+\infty]$. Convergence to a finite real number additionally requires that this common value is finite.
%
% !TEX root = ../main.tex
\subsection*{HW 20 --- An oscillating sequence}%
\manualtoc{subsection}{\texorpdfstring{HW 20 --- An oscillating sequence}{HW 20 --- An oscillating sequence}}{studiabile-nav-0095}%

Consider the sequence $x_n=(-1)^n$, $n\in\mathbb N$, calculate its liminf and limsup. Does the limit exists?

\par\noindent\studyheading{Proof.}\par\nobreak
Notice that $x_n=\begin{cases}-1&n\text{ odd},\\1&n\text{ even}.\end{cases}$
Therefore $\inf_{m\ge n}x_m=-1$ and $\sup_{m\ge n}x_m=1$.
Therefore
$\liminf_{n\to+\infty}x_n=\sup_{n\ge1}\inf_{m\ge n}x_m=\sup_{n\ge1}(-1)=-1$
and in the same way $\limsup_{n\to+\infty}x_n=1$.
Therefore they don't coincide and the limit does not exist.
%
% !TEX root = ../main.tex
\subsection*{HW 21 --- Threshold properties of limits}%
\manualtoc{subsection}{\texorpdfstring{HW 21 --- Threshold properties of limits}{HW 21 --- Threshold properties of limits}}{studiabile-nav-0097}%

Prove the following properties of liminf and limsup.

\studylabel{(i)} If $z>\limsup_{n\to+\infty}x_n$, then there exists $N\in\mathbb N$ such that $x_n<z$ $\forall n\ge N$.

\studylabel{(ii)} If $z<\liminf_{n\to+\infty}x_n$ then there exists $N\in\mathbb N$ such that $x_n>z$ $\forall n\ge N$.

\studylabel{(iii)} If $z<\limsup_{n\to+\infty}x_n$ then there are infinitely many $n\in\mathbb N$ such that $x_n>z$.

\studylabel{(iv)} If $z>\liminf_{n\to+\infty}x_n$ then there are infinitely many $n\in\mathbb N$ such that $x_n<z$.

\par\noindent\studyheading{Proof.}\par\nobreak
\studylabel{(i)} Assume that
$z>\limsup_{n\to+\infty}x_n=\inf_{n\ge1}\sup_{m\ge n}x_m:=\inf_{n\ge1}y_n$,
where $y_n:=\sup_{m\ge n}x_m$.
By definition of infimum, if $\inf_{n\ge1}y_n\in\mathbb R$ we have $\forall\varepsilon>0$ $\exists N=N(\varepsilon)\in\mathbb N$ such that $y_N<\inf_{n\ge1}y_n+\varepsilon$.
Combining these two inequalities it exist $\varepsilon>0$ sufficiently small so that
$y_N\le\inf_{n\ge1}y_n+\varepsilon<z$.
Restoring the definition of $y_N$ we obtain
\[
x_m\le\sup_{j\ge N}x_j=y_N<z\quad\forall m\ge N
\quad\Rightarrow\quad x_m<z\quad\forall m\ge N,
\]
where the first inequality follows from the definition of supremum.

\studylabel{(ii)} Assume that $z<\liminf_{n\to+\infty}x_n=\sup_{n\ge1}\inf_{m\ge n}x_m:=\sup_{n\ge1}y_n$,
where $y_n:=\inf_{m\ge n}x_m$.
By definition of supremum if $\sup_{n\ge1}y_n\in\mathbb R$ we have $\forall\varepsilon>0$ $\exists N=N(\varepsilon)\in\mathbb N$ such that $y_N>\sup_{n\ge1}y_n-\varepsilon$.
Since $z<\sup_{n\ge1}y_n$ we can choose $\varepsilon>0$ sufficiently small such that $z<\sup_{n\ge1}y_n-\varepsilon$.
Combining these two inequalities we obtain $z\le\sup_{n\ge1}y_n-\varepsilon<y_N$.
Restoring the definition of $y_N$ we obtain
\[
x_m\ge\inf_{j\ge N}x_j=y_N>z\quad\forall m\ge N.
\]
Hence $x_m>z$ $\forall m\ge N$.

For the tail infima in \studylabel{(ii)}, if $\sup_{n\ge1}y_n=+\infty$ its definition directly guarantees an $N$ such that $y_N>z$ and the conclusion follows in the same way.

For the tail suprema $y_n=\sup_{m\ge n}x_m$ in \studylabel{(i)}, if $\inf_{n\ge1}y_n=-\infty$ then by definition of infimum $\forall z\in\mathbb R$ there exists $N\in\mathbb N$ such that $y_N<z$ and restoring def of $y_N$ we obtain $x_m<z$ $\forall m\ge N$.

\studylabel{(iii)} Assume that $z<\limsup_{n\to+\infty}x_n=\inf_{n\ge1}\sup_{m\ge n}x_m=\inf_{n\ge1}y_n$.
By the definition of infimum we have $z<\inf_{n\ge1}y_n\le y_n$ $\forall n\in\mathbb N$.
Restoring the definition of $y_n$ we obtain $z<\sup_{m\ge n}x_m$ $\forall n\in\mathbb N$.
Taking $n=1$ by the definition of supremum there exists $n_1\in\mathbb N$ such that $x_{n_1}>z$.
Now taking $n=n_1+1$ we have $z<\sup_{m\ge n_1+1}x_m$.
Hence by definition of supremum there exists $n_2>n_1$ such that $x_{n_2}>z$.
Repeating this argument we obtain a strictly increasing sequence of indices $n_1<n_2<n_3<\cdots$ such that $x_{n_j}>z$ $\forall j\in\mathbb N$.
Therefore $x_n>z$ for infinitely many $n\in\mathbb N$.

\studylabel{(iv)} Assume that $z>\liminf_{n\to+\infty}x_n=\sup_{n\ge1}\inf_{m\ge n}x_m:=\sup_{n\ge1}y_n$.
By the definition of supremum we have $z>\sup_{n\ge1}y_n\ge y_n$ $\forall n\in\mathbb N$.
Restoring the definition of $y_n$ we obtain $z>\inf_{m\ge n}x_m$ $\forall n\in\mathbb N$.
Taking $n=1$ by the definition of infimum there exists $n_1\in\mathbb N$ such that $x_{n_1}<z$.
Now taking $n=n_1+1$ we have $z>\inf_{m\ge n_1+1}x_m$.
Hence by the definition of infimum there exists $n_2>n_1$ such that $x_{n_2}<z$.
Repeating this argument we obtain a strictly increasing sequence of indices $n_1<n_2<n_3<\cdots$ such that $x_{n_j}<z$ $\forall j\in\mathbb N$.
Therefore $x_n<z$ for infinitely many $n\in\mathbb N$.
%
% !TEX root = ../main.tex
\subsection*{HW 22 --- Negation, suprema and infima}%
\manualtoc{subsection}{\texorpdfstring{HW 22 --- Negation, suprema and infima}{HW 22 --- Negation, suprema and infima}}{studiabile-nav-0107}%

Let $(x_\alpha)_{\alpha\in\mathcal I}$ be a family of real numbers ($\mathcal I\ne\emptyset$ is not necessarily countable, e.g., $\mathcal I=[0,1]$).
Prove that $\sup_{\alpha\in\mathcal I}x_\alpha=-\inf_{\alpha\in\mathcal I}(-x_\alpha)$ \studylabel{(i)}.
Deduce that $\liminf_{n\to+\infty}x_n=-\limsup_{n\to+\infty}(-x_n)$ \studylabel{(ii)}.

\par\noindent\studyheading{Proof.}\par\nobreak
\studylabel{(i)} Let $M=\inf_{\alpha\in\mathcal I}(-x_\alpha)$. We distinguish two cases.

$M\in\mathbb R$. By definition of infimum we have $M\le-x_\alpha$ $\forall\alpha\in\mathcal I$ and for every $\varepsilon>0$ there exists $\beta\in\mathcal I$ such that $-x_\beta\le M+\varepsilon$.
Multiplying both inequalities by $(-1)$ we obtain $-M\ge x_\alpha$ $\forall\alpha\in\mathcal I$ and $x_\beta\ge-M-\varepsilon$.
Therefore $-M$ is an upper bound of $(x_\alpha)_{\alpha\in\mathcal I}$ and for every $\varepsilon>0$ there exists $\beta\in\mathcal I$ such that $-M-\varepsilon\le x_\beta$.
Hence by definition of supremum $\sup_{\alpha\in\mathcal I}x_\alpha=-M=-\inf_{\alpha\in\mathcal I}(-x_\alpha)$.

$M=-\infty$. By definition of infimum, for every $K\in\mathbb R$ there exists $\beta\in\mathcal I$ such that $-x_\beta<-K$.
Multiplying by $(-1)$ we obtain $x_\beta>K$.
Therefore $(x_\alpha)_{\alpha\in\mathcal I}$ is not bounded above and hence $\sup_{\alpha\in\mathcal I}x_\alpha=+\infty$.
Since $M=-\infty$ we conclude that $\sup_{\alpha\in\mathcal I}x_\alpha=+\infty=-M=-\inf_{\alpha\in\mathcal I}(-x_\alpha)$.

$M=+\infty$ is impossible since $\mathcal I\ne\emptyset$ and every $x_\alpha\in\mathbb R$.

\studylabel{(ii)} By definition $\liminf_{n\to+\infty}x_n=\sup_{n\ge1}\inf_{m\ge n}x_m$.
Notice that
\[
\begin{aligned}
\liminf_{n\to+\infty}x_n
&=\sup_{n\ge1}\inf_{m\ge n}x_m
=\sup_{n\ge1}\left(-\sup_{m\ge n}(-x_m)\right) =-\inf_{n\ge1}\sup_{m\ge n}(-x_m)
=-\limsup_{n\to+\infty}(-x_n),
\end{aligned}
\]
where the second equality follows by \studylabel{(i)} taking $\inf_{m\ge n}x_m=-\sup_{m\ge n}(-x_m)$ and the third equality follows by \studylabel{(i)} taking $\sup_{n\ge1}(-a_n)=-\inf_{n\ge1}a_n$ with $a_n:=\sup_{m\ge n}(-x_m)$.
The same reversal of order exchanges arbitrary upper bounds with lower bounds even when the bounds are infinite. Thus both applications of the sup/inf identity remain valid for the extended-real tail values.
%
% !TEX root = ../main.tex
\subsection*{Def limsup of collection of sets}%
\manualtoc{subsection}{\texorpdfstring{Def limsup of collection of sets}{Def limsup of collection of sets}}{studiabile-nav-0111}%

Let $(E_n)_{n\in\mathbb N}\subseteq\mathcal F$ and define
\[
\begin{aligned}
\limsup_{n\to+\infty}E_n
&:=\{\omega\in\Omega:\omega\in E_n\text{ for infinitely many }n\text{'s}\}\\&=\{\omega\in\Omega:\text{for any }n\in\mathbb N\text{ there is }m\ge n\text{ such that }\omega\in E_m\} =\bigcap_{n\ge1}\bigcup_{m\ge n}E_m.
\end{aligned}
\]
For fixed $n$, $\bigcup_{m\ge n}E_m$ contains exactly those $\omega$ that belong to at least one $E_m$ with $m\ge n$. Intersecting over all $n$ requires such occurrences arbitrarily far out in the sequence.
Notice also that $\limsup_{n\to+\infty}E_n\in\mathcal F$ because it is countable intersection of countable unions of sets in $\mathcal F$.
%
% !TEX root = ../main.tex
\subsection*{HW 23 --- Measurability of the limsup event}%
\manualtoc{subsection}{\texorpdfstring{HW 23 --- Measurability of the limsup event}{HW 23 --- Measurability of the limsup event}}{studiabile-nav-0112}%

Prove that $\limsup_{n\to+\infty}E_n\in\mathcal F$.
\par\noindent\studyheading{Proof.}\par\nobreak
Let $(\Omega,\mathcal F,\mathbb P)$ be a probability space and let $(E_n)_{n\in\mathbb N}\subseteq\mathcal F$.
By definition $\limsup_{n\to+\infty}E_n=\bigcap_{n\ge1}\bigcup_{m\ge n}E_m$.
Since $\mathcal F$ is a $\sigma$-algebra is closed under countable unions therefore $\bigcup_{m\ge n}E_m\in\mathcal F$ $\forall n\in\mathbb N$.
Moreover by HW~3, a $\sigma$-algebra is also closed under countable intersections hence
$\bigcap_{n\ge1}(\bigcup_{m\ge n}E_m)\in\mathcal F$.
%
% !TEX root = ../main.tex
\subsection*{Lemma 2.2 (Reverse Fatou's lemma for sets)}%
\manualtoc{subsection}{\texorpdfstring{Lemma 2.2 (Reverse Fatou's lemma for sets)}{Lemma 2.2 (Reverse Fatou's lemma for sets)}}{studiabile-nav-0114}%

Let $(\Omega,\mathcal F,\mathbb P)$.
Given $(E_n)_{n\in\mathbb N}\subseteq\mathcal F$ we have
\[
\mathbb P\left(\limsup_{n\to+\infty}E_n\right)\ge\limsup_{n\to+\infty}\mathbb P(E_n).
\]
\par\noindent\studyheading{Proof.}\par\nobreak
Define $G_m:=\bigcup_{n\ge m}E_n$.
Notice that $G_{m+1}\subseteq G_m$ $\forall m\in\mathbb N$.
Indeed by definition $G_m=\bigcup_{n\ge m}E_n=E_m\cup(\bigcup_{n\ge m+1}E_n)=E_m\cup G_{m+1}\supseteq G_{m+1}$.
Therefore $(G_m)_{m\in\mathbb N}$ is a non-increasing sequence of sets and $G_m\downarrow G$ where $G=\bigcap_{m=1}^\infty G_m=\bigcap_{m=1}^\infty\bigcup_{n\ge m}E_n=\limsup_{n\to+\infty}E_n$.
Since $\mathbb P$ is a probability measure we have $\mathbb P(G_m)\le1<\infty$ $\forall m\in\mathbb N$.
Hence by Lemma 1.6 (continuity from above) $\mathbb P(G)=\lim_{m\to+\infty}\mathbb P(G_m)$.
Moreover by definition of $G_m$ we have $E_n\subseteq G_m$ $\forall n\ge m$.
By the monotonicity of probability measures (HW~12), $\mathbb P(E_n)\le\mathbb P(G_m)$ $\forall n\ge m$.
Taking the limit as $m\to+\infty$ we have
$\lim_{m\to+\infty}\sup_{n\ge m}\mathbb P(E_n)\le\lim_{m\to+\infty}\mathbb P(G_m)$.

Recalling the definition of limsup for numerical sequences and using Lemma 1.6 we conclude that
\[
\begin{aligned}
\mathbb P\left(\limsup_{n\to+\infty}E_n\right)
&=\mathbb P(G)=\lim_{m\to+\infty}\mathbb P(G_m) \ge\lim_{m\to+\infty}\sup_{n\ge m}\mathbb P(E_n)
=\limsup_{n\to+\infty}\mathbb P(E_n).
\end{aligned}
\]
Hence $\mathbb P(\limsup_{n\to+\infty}E_n)\ge\limsup_{n\to+\infty}\mathbb P(E_n)$.
%
% !TEX root = ../main.tex
\subsection*{Lemma 2.3 (First Borel--Cantelli's lemma)}%
\manualtoc{subsection}{\texorpdfstring{Lemma 2.3 (First Borel--Cantelli's lemma)}{Lemma 2.3 (First Borel--Cantelli's lemma)}}{studiabile-nav-0116}%

Let $(\Omega,\mathcal F,\mathbb P)$ be a probability space and let $(E_n)_{n\in\mathbb N}\subseteq\mathcal F$ be such that $\sum_{n=1}^\infty\mathbb P(E_n)<\infty$.
Then $\mathbb P(\limsup_{n\to+\infty}E_n)=0$ (i.e. almost every $\omega$ belongs to only finitely many $E_n$'s).

\par\noindent\studyheading{Proof.}\par\nobreak
Define $G_m:=\bigcup_{n\ge m}E_n$.
Notice that $G_{m+1}\subseteq G_m$ $\forall m\in\mathbb N$.
Indeed by definition $G_m=E_m\cup G_{m+1}$.
Therefore $(G_m)$ is a non-increasing sequence of sets and $G_m\downarrow G$ where $G=\bigcap_{m=1}^\infty G_m=\limsup_{n\to+\infty}E_n$.
Since $\mathbb P(G_m)\le1<\infty$ by Lemma 1.6 (continuity from above) we obtain $\mathbb P(G)=\lim_{m\to+\infty}\mathbb P(G_m)$.
Now for each fixed $m\in\mathbb N$ we define $G_m^k:=\bigcup_{n=m}^{m+k}E_n$.
Notice that $G_m^k\subseteq G_m^{k+1}$ since $G_m^{k+1}=G_m^k\cup E_{m+k+1}$.
Therefore $(G_m^k)_{k\in\mathbb N}$ is a non-decreasing and $G_m^k\uparrow G_m=\bigcup_{n\ge m}E_n$.
Hence by Lemma 1.5 (continuity from below) $\mathbb P(G_m)=\lim_{k\to+\infty}\mathbb P(G_m^k)$.
Moreover by finite sub-additivity of probability measures (Lemma 1.4)
$\mathbb P(G_m^k)=\mathbb P(\bigcup_{n=m}^{m+k}E_n)\le\sum_{n=m}^{m+k}\mathbb P(E_n)$.
Taking the limit as $k\to+\infty$ we obtain
\[
\mathbb P(G_m)=\lim_{k\to+\infty}\mathbb P(G_m^k)
\le\lim_{k\to+\infty}\sum_{n=m}^{m+k}\mathbb P(E_n)=\sum_{n=m}^\infty\mathbb P(E_n).
\]
Combining this inequality with the previous result we have
\[
\mathbb P\left(\limsup_{n\to+\infty}E_n\right)
=\mathbb P(G)=\lim_{m\to+\infty}\mathbb P(G_m)
\le\lim_{m\to+\infty}\sum_{n=m}^\infty\mathbb P(E_n)=0,
\]
where the last equality follows from the ass. $\sum_{n=1}^\infty\mathbb P(E_n)<\infty$.
Finally, since probabilities are non-negative $0\le\mathbb P(\limsup_{n\to+\infty}E_n)\le0$.
Therefore $\mathbb P(\limsup_{n\to+\infty}E_n)=0$.
%
% !TEX root = ../main.tex
\subsection*{Def liminf of collection of sets}%
\manualtoc{subsection}{\texorpdfstring{Def liminf of collection of sets}{Def liminf of collection of sets}}{studiabile-nav-0118}%

Let $(E_n)_{n\in\mathbb N}\subseteq\mathcal F$ and define
\[
\begin{aligned}
\liminf_{n\to+\infty}E_n
&:=\{\omega\in\Omega:\omega\in E_n\text{ for all but finitely many }n\text{'s}\}\\&=\{\omega\in\Omega:\exists N_\omega\in\mathbb N\text{ such that }\omega\in E_n\ \forall n\ge N_\omega\} =\bigcup_{n\ge1}\bigcap_{m\ge n}E_m.
\end{aligned}
\]
Notice that $\bigcap_{m\ge n}E_m$ contains those $\omega$'s that belong to all sets $E_m$ for $m\ge n$.
Taking the union allows to include all those $\omega$'s that belong to all but finitely many sets $E_n$.
Notice that $\liminf_{n\to+\infty}E_n\in\mathcal F$ because it is countable union of countable intersections of sets in $\mathcal F$.
%
% !TEX root = ../main.tex
\subsection*{HW 24 --- Measurability of the liminf event}%
\manualtoc{subsection}{\texorpdfstring{HW 24 --- Measurability of the liminf event}{HW 24 --- Measurability of the liminf event}}{studiabile-nav-0119}%

Prove that $\liminf_{n\to+\infty}E_n\in\mathcal F$.

\par\noindent\studyheading{Proof.}\par\nobreak
Define $Y_n:=\bigcap_{m\ge n}E_m=E_n\cap(\bigcap_{m\ge n+1}E_m)=E_n\cap Y_{n+1}$.
Therefore $Y_n\subseteq Y_{n+1}$ $\forall n\in\mathbb N$, hence $(Y_n)_{n\in\mathbb N}$ is a non-decreasing sequence of sets.
Moreover since $E_m\in\mathcal F$ $\forall m\in\mathbb N$ and $\mathcal F$ is a $\sigma$-algebra, by HW~3 it follows that $Y_n=\bigcap_{m\ge n}E_m\in\mathcal F$ $\forall n\in\mathbb N$.
Finally since $\mathcal F$ is closed under countable unions, we obtain
\[
\liminf_{n\to+\infty}E_n=\bigcup_{n=1}^\infty Y_n
=\bigcup_{n=1}^\infty\left(\bigcap_{m\ge n}E_m\right)\in\mathcal F.
\]
%
% !TEX root = ../main.tex
\subsection*{HW 25 --- Complement relation of event limits}%
\manualtoc{subsection}{\texorpdfstring{HW 25 --- Complement relation of event limits}{HW 25 --- Complement relation of event limits}}{studiabile-nav-0121}%

Prove that $\liminf_{n\to+\infty}E_n=(\limsup_{n\to+\infty}E_n^c)^c$.

\par\noindent\studyheading{Proof.}\par\nobreak
We obtain
\[
\begin{aligned}
\left(\limsup_{n\to+\infty}E_n^c\right)^c
&=\left(\bigcap_{n=1}^\infty\bigcup_{m\ge n}E_m^c\right)^c
=\bigcup_{n=1}^\infty\left(\bigcup_{m\ge n}E_m^c\right)^c\\
&=\bigcup_{n=1}^\infty\bigcap_{m\ge n}(E_m^c)^c
=\bigcup_{n=1}^\infty\bigcap_{m\ge n}E_m
=\liminf_{n\to+\infty}E_n,
\end{aligned}
\]
where the second and third equalities follow De Morgan's laws, the fourth from double complement property and the last from the definition of liminf.
%
% !TEX root = ../main.tex
\subsection*{Lemma 2.4 (Fatou's lemma for sets)}%
\manualtoc{subsection}{\texorpdfstring{Lemma 2.4 (Fatou's lemma for sets)}{Lemma 2.4 (Fatou's lemma for sets)}}{studiabile-nav-0123}%

Let $(\Omega,\mathcal F,\mathbb P)$ be a probability space and let $(E_n)_{n\in\mathbb N}\subseteq\mathcal F$. Prove that
\[
\mathbb P\left(\liminf_{n\to+\infty}E_n\right)\le\liminf_{n\to+\infty}\mathbb P(E_n).
\]
%
% !TEX root = ../main.tex
\subsection*{HW 26 --- Fatou inequality by complements}%
\manualtoc{subsection}{\texorpdfstring{HW 26 --- Fatou inequality by complements}{HW 26 --- Fatou inequality by complements}}{studiabile-nav-0124}%

Let $(\Omega,\mathcal F,\mathbb P)$ be a probability space and let $(E_n)_{n\in\mathbb N}\subseteq\mathcal F$. Prove Lemma~2.4, namely that
\[
\mathbb P\left(\liminf_{n\to+\infty}E_n\right)\le\liminf_{n\to+\infty}\mathbb P(E_n).
\]


% BEGIN PFF RESTORATION HW26_HINT
\studyheading{\inlinetoc{subsubsection}{\texorpdfstring{Hint:}{Hint:}}{studiabile-nav-0125}Hint:} follow the steps in the proof of Lemma~2.2.
% END PFF RESTORATION HW26_HINT
\par\noindent\studyheading{Proof.}\par\nobreak
We obtain
\[
\begin{aligned}
\mathbb P\left(\liminf_{n\to+\infty}E_n\right)
&=\mathbb P\left(\left(\limsup_{n\to+\infty}E_n^c\right)^c\right) =1-\mathbb P\left(\limsup_{n\to+\infty}E_n^c\right)
\le1-\limsup_{n\to+\infty}\mathbb P(E_n^c),
\end{aligned}
\]
The first equality is HW~25, the second is the complement identity from HW~18, and the inequality is Lemma~2.2 applied to $(E_n^c)$, after multiplying by $-1$ and adding $1$. Hence
\[
\begin{aligned}
\mathbb P\left(\liminf_{n\to+\infty}E_n\right)
&\le1-\limsup_{n\to+\infty}\mathbb P(E_n^c)
=1-\limsup_{n\to+\infty}(1-\mathbb P(E_n))\\&=1-\left(1+\limsup_{n\to+\infty}(-\mathbb P(E_n))\right) =-\limsup_{n\to+\infty}(-\mathbb P(E_n))
=\liminf_{n\to+\infty}\mathbb P(E_n),
\end{aligned}
\]
The first equality again uses HW~18; the next two use the constant-shift rule for limsup and cancel $1$, and the last is HW~22.
%
% !TEX root = ../main.tex
\subsection*{Definition of indicator function}%
\manualtoc{subsection}{\texorpdfstring{Definition of indicator function}{Definition of indicator function}}{studiabile-nav-0127}%

The indicator function of a set $E\in\mathcal F$ is defined as
\[
\mathbf{1}_E(\omega)=\begin{cases}1&\text{if }\omega\in E,\\0&\text{if }\omega\notin E.\end{cases}
\]
For example, if we take $\Omega=[0,1]$ (i.e. $\omega$'s are numbers in $[0,1]$) and $\mathcal F=\mathcal B([0,1])$, then $\mathbf{1}_{[\frac12,1]}(\omega)$ is a function which is zero when $\omega\in[0,\frac12)$ and it is equal to one when $\omega\in[\frac12,1]$.
%
% !TEX root = ../main.tex
\subsection*{HW 27 --- Indicators and limsup}%
\manualtoc{subsection}{\texorpdfstring{HW 27 --- Indicators and limsup}{HW 27 --- Indicators and limsup}}{studiabile-nav-0128}%


% BEGIN PFF RESTORATION HW27_SOURCE_ATTRIBUTION
(Ex.~2.9 from the book.)
% END PFF RESTORATION HW27_SOURCE_ATTRIBUTION
Given $(E_n)_{n\in\mathbb N}\subseteq\mathcal F$. Prove that, for every $\omega\in\Omega$
\[
\mathbf{1}_{\limsup_{n\to+\infty}E_n}(\omega)=\limsup_{n\to+\infty}\mathbf{1}_{E_n}(\omega).
\]
\studyheading{\inlinetoc{subsubsection}{\texorpdfstring{Hint:}{Hint:}}{studiabile-nav-0129}Hint:} follow the intermediate steps set out below.
\studylabel{(i)} prove that $\mathbf{1}_{E_i\cap E_j}=\mathbf{1}_{E_i}\cdot\mathbf{1}_{E_j}=\min\{\mathbf{1}_{E_i},\mathbf{1}_{E_j}\}$;
\studylabel{(ii)} prove that $\mathbf{1}_{E_i\cup E_j}=\mathbf{1}_{E_i}+\mathbf{1}_{E_j}-\mathbf{1}_{E_i\cap E_j}=\max\{\mathbf{1}_{E_i},\mathbf{1}_{E_j}\}$;
\studylabel{(iii)} use \studylabel{(i)} and \studylabel{(ii)} to conclude.

\par\noindent\studyheading{Proof.}\par\nobreak
\emph{\studylabel{(i)} We prove the intersection identity by its two possible values.}
\[
\mathbf{1}_{A\cap B}(\omega)=\begin{cases}
\begin{aligned}[t]
0\quad&\text{if }\omega\notin A\cap B\Longleftrightarrow\omega\notin A\text{ or }\omega\notin B\\
&\Longleftrightarrow\mathbf{1}_A(\omega)\mathbf{1}_B(\omega)=0=\min\{\mathbf{1}_A(\omega),\mathbf{1}_B(\omega)\},
\end{aligned}\\[3pt]
\begin{aligned}[t]
1\quad&\text{if }\omega\in A\cap B\Longleftrightarrow\omega\in A\text{ and }\omega\in B\\
&\Longleftrightarrow\mathbf{1}_A(\omega)\mathbf{1}_B(\omega)=1=\min\{\mathbf{1}_A(\omega),\mathbf{1}_B(\omega)\}.
\end{aligned}
\end{cases}
\]

\emph{\studylabel{(ii)} We prove the union identity by its two possible values.}
\[
\mathbf{1}_{A\cup B}(\omega)=\begin{cases}
\begin{aligned}[t]
0\quad&\text{if }\omega\notin A\text{ and }\omega\notin B\\
&\Longleftrightarrow\mathbf{1}_A(\omega)+\mathbf{1}_B(\omega)-\mathbf{1}_{A\cap B}(\omega)=0=\max\{\mathbf{1}_A(\omega),\mathbf{1}_B(\omega)\},
\end{aligned}\\[3pt]
\begin{aligned}[t]
1\quad&\text{if }\omega\in A\text{ or }\omega\in B\\
&\Longleftrightarrow\mathbf{1}_A(\omega)+\mathbf{1}_B(\omega)-\mathbf{1}_{A\cap B}(\omega)=1=\max\{\mathbf{1}_A(\omega),\mathbf{1}_B(\omega)\}.
\end{aligned}
\end{cases}
\]
The same pointwise reasoning applies to countably many indicators. An intersection indicator equals one exactly when all the indicators equal one; a union indicator equals one exactly when at least one equals one. Since every value is either zero or one, infimum and supremum are attained. Thus
\[
\mathbf{1}_{\bigcap_{n=1}^\infty A_n}(\omega)=\inf_{n\ge1}\mathbf{1}_{A_n}(\omega)=\min_{n\ge1}\mathbf{1}_{A_n}(\omega),
\qquad
\mathbf{1}_{\bigcup_{n=1}^\infty A_n}(\omega)=\sup_{n\ge1}\mathbf{1}_{A_n}(\omega)=\max_{n\ge1}\mathbf{1}_{A_n}(\omega).
\]
For simplicity define $F_n:=\bigcup_{m\ge n}E_m$ and let $\omega\in\Omega$, we have
\[
\begin{aligned}
\mathbf{1}_{\limsup_{n\to+\infty}E_n}(\omega)
&=\mathbf{1}_{\bigcap_{n\ge1}F_n}(\omega)=\inf_{n\ge1}\mathbf{1}_{F_n}(\omega)\\
&=\inf_{n\ge1}\mathbf{1}_{\bigcup_{m\ge n}E_m}(\omega)
=\inf_{n\ge1}\sup_{m\ge n}\mathbf{1}_{E_m}(\omega)
=\limsup_{n\to+\infty}\mathbf{1}_{E_n}(\omega).
\end{aligned}
\]
In the chain, the first equality substitutes the limsup event, the second uses the intersection identity, the third substitutes $F_n$, the fourth uses the union identity, and the last is the numerical definition of limsup. Each step holds for the same arbitrary $\omega$.
%
\clearpage
%
\section*{3. Random Variables}
\manualtoc{section}{\texorpdfstring{3. Random Variables}{3. Random Variables}}{studiabile-nav-0136}
% !TEX root = ../main.tex
\subsection*{Random variables and pre-images}%
\manualtoc{subsection}{\texorpdfstring{Random variables and pre-images}{Random variables and pre-images}}{studiabile-nav-0137}%

Random variables are measurable functions.
%
% !TEX root = ../main.tex
\subsection*{Def --- pre-image of a set}%
\manualtoc{subsection}{\texorpdfstring{Def --- pre-image of a set}{Def --- pre-image of a set}}{studiabile-nav-0138}%

Let $(S,\Sigma)$ be a measurable space and let $h:S\to\mathbb R$ be a function.
We define the pre-image of a set $A\subseteq\mathbb R$ as $h^{-1}(A)=\{s\in S:h(s)\in A\}$.
%
% !TEX root = ../main.tex
\subsection*{Def 3.1 (Measurable functions)}%
\manualtoc{subsection}{\texorpdfstring{Def 3.1 (Measurable functions)}{Def 3.1 (Measurable functions)}}{studiabile-nav-0139}%

A function $h:S\to\mathbb R$ is $\Sigma$-measurable (also denoted $h\in\mathrm m\Sigma$) if $h^{-1}(A)\in\Sigma$ $\forall A\in\mathcal B(\mathbb R)$.
Equivalently we say that $h\in\mathrm m\Sigma$ if $h^{-1}:\mathcal B(\mathbb R)\to\Sigma$.

\studyheading{\inlinetoc{subsubsection}{\texorpdfstring{Notation:}{Notation:}}{studiabile-nav-0140}Notation:} we use $(\mathrm m\Sigma)^+$ for the class of non-negative $\Sigma$-measurable functions.
We use $\mathrm b\Sigma$ for the class of bounded $\Sigma$-measurable functions.
\studyheading{\inlinetoc{subsubsection}{\texorpdfstring{Borel-measurable functions:}{Borel-measurable functions:}}{studiabile-nav-0141}Borel-measurable functions:} If $S$ is a topological space and $\Sigma=\mathcal B(S)$, then $\mathrm m\Sigma$ is the class of Borel-measurable functions.
For example if $S=\mathbb R$, $\Sigma=\mathcal B(\mathbb R)$ and $h:\mathbb R\to\mathbb R$ is such $h\in\mathrm m\Sigma$, then $h$ is Borel-measurable.
%
% !TEX root = ../main.tex
\subsection*{Properties of pre-image mapping}%
\manualtoc{subsection}{\texorpdfstring{Properties of pre-image mapping}{Properties of pre-image mapping}}{studiabile-nav-0142}%

Let $h:S\to\mathbb R$ and let $(A_\alpha)_{\alpha\in\mathcal I}$ be a family of subsets of $\mathbb R$ with $\mathcal I$ possibly uncountable (i.e. $\mathcal I=[0,1]$). Then
\[
\begin{aligned}
\text{\studylabel{(i)}}\quad&h^{-1}\left(\bigcup_{\alpha\in\mathcal I}A_\alpha\right)=\bigcup_{\alpha\in\mathcal I}h^{-1}(A_\alpha),\\
\text{\studylabel{(ii)}}\quad&h^{-1}(A^c)=(h^{-1}(A))^c,\\
\text{\studylabel{(iii)}}\quad&h^{-1}\left(\bigcap_{\alpha\in\mathcal I}A_\alpha\right)=\bigcap_{\alpha\in\mathcal I}h^{-1}(A_\alpha).
\end{aligned}
\]
\par\noindent\studyheading{Proof \studylabel{(i)}.}\par\nobreak
\emph{We first prove that $h^{-1}(\bigcup_{\alpha\in\mathcal I}A_\alpha)\subseteq\bigcup_{\alpha\in\mathcal I}h^{-1}(A_\alpha)$.}
Let $s_0\in h^{-1}(\bigcup_{\alpha\in\mathcal I}A_\alpha)$.
By definition of pre-image we have $h(s_0)\in\bigcup_{\alpha\in\mathcal I}A_\alpha$.
By definition of union there exists $\alpha_0\in\mathcal I$ such that $h(s_0)\in A_{\alpha_0}$.
Therefore $s_0\in h^{-1}(A_{\alpha_0})\subseteq\bigcup_{\alpha\in\mathcal I}h^{-1}(A_\alpha)$, where the membership $(\in)$ follows from the definition of pre-image and the inclusion follows the fact that $h^{-1}(A_{\alpha_0})$ is one of the sets in the union.
Since $s_0$ was arbitrary it follows that $h^{-1}(\bigcup_{\alpha\in\mathcal I}A_\alpha)\subseteq\bigcup_{\alpha\in\mathcal I}h^{-1}(A_\alpha)$.
For the reverse inclusion, let $s_0\in\bigcup_{\alpha\in\mathcal I}h^{-1}(A_\alpha)$.
By definition of union, there exists $\alpha_0\in\mathcal I$ such that $s_0\in h^{-1}(A_{\alpha_0})$.
By definition of pre-image we obtain $h(s_0)\in A_{\alpha_0}$.
Therefore $h(s_0)\in A_{\alpha_0}\subseteq\bigcup_{\alpha\in\mathcal I}A_\alpha$

where the inclusion follows from the definition of union.
Hence by definition of pre-image $s_0\in h^{-1}(\bigcup_{\alpha\in\mathcal I}A_\alpha)$.
Since $s_0$ is arbitrary we obtain $\bigcup_{\alpha\in\mathcal I}h^{-1}(A_\alpha)\subseteq h^{-1}(\bigcup_{\alpha\in\mathcal I}A_\alpha)$.
Combining the two inclusions we conclude the proof.
%
% !TEX root = ../main.tex
\subsection*{HW 28 --- Pre-images of complements and intersections}%
\manualtoc{subsection}{\texorpdfstring{HW 28 --- Pre-images of complements and intersections}{HW 28 --- Pre-images of complements and intersections}}{studiabile-nav-0144}%

For $h:S\to\mathbb R$, prove properties \studylabel{(ii)} and \studylabel{(iii)}:
\[
h^{-1}(A^c)=(h^{-1}(A))^c,\qquad
h^{-1}\left(\bigcap_{\alpha\in\mathcal I}A_\alpha\right)=\bigcap_{\alpha\in\mathcal I}h^{-1}(A_\alpha).
\]
Complements on the left are taken in $\mathbb R$ and those on the right in $S$.
\par\noindent\studyheading{Proof.}\par\nobreak
\studylabel{(ii)} Let $A\subseteq\mathbb R$, \emph{we first prove that $h^{-1}(A^c)\subseteq(h^{-1}(A))^c$.}
Let $s_0\in h^{-1}(A^c)$.
By definition of pre-image we have $h(s_0)\in A^c\Rightarrow h(s_0)\notin A$.
Therefore $s_0\notin h^{-1}(A)\Rightarrow s_0\in(h^{-1}(A))^c$ where the first statement follows from the definition of pre-image and the second from the definition of complement.
Since $s_0$ is arbitrary it follows that $h^{-1}(A^c)\subseteq(h^{-1}(A))^c$.
For the reverse inclusion let $s_0\in(h^{-1}(A))^c\Rightarrow s_0\notin h^{-1}(A)$.
By the definition of pre-image this implies that $h(s_0)\notin A$.
Therefore $h(s_0)\in A^c\Rightarrow s_0\in h^{-1}(A^c)$ where the first membership follows from the definition of complement and the second from the definition of pre-image.
Since $s_0$ was arbitrary we obtain $(h^{-1}(A))^c\subseteq h^{-1}(A^c)$.
Combining the two inclusionss we conclude the proof.

\studylabel{(iii)}
Let $(A_\alpha)_{\alpha\in\mathcal I}$ be a family of subsets of $\mathbb R$ where $\mathcal I$ is possibly uncountable.
\emph{We first prove that $h^{-1}(\bigcap_{\alpha\in\mathcal I}A_\alpha)\subseteq\bigcap_{\alpha\in\mathcal I}h^{-1}(A_\alpha)$.}
Let $s_0\in h^{-1}(\bigcap_{\alpha\in\mathcal I}A_\alpha)$.
By definition of pre-image we have $h(s_0)\in\bigcap_{\alpha\in\mathcal I}A_\alpha$.
By definition of intersection this means that $h(s_0)\in A_\alpha$ $\forall\alpha\in\mathcal I$.
Therefore $s_0\in h^{-1}(A_\alpha)$ $\forall\alpha\in\mathcal I$ where the membership follows the def of pre-image.
Hence $s_0\in\bigcap_{\alpha\in\mathcal I}h^{-1}(A_\alpha)$ by def of intersection.
Since $s_0$ is arbitrary it follows that $h^{-1}(\bigcap_{\alpha\in\mathcal I}A_\alpha)\subseteq\bigcap_{\alpha\in\mathcal I}h^{-1}(A_\alpha)$.
For the reverse inclusion, let $s_0\in\bigcap_{\alpha\in\mathcal I}h^{-1}(A_\alpha)$.
By definition of intersection we have $s_0\in h^{-1}(A_\alpha)$ $\forall\alpha\in\mathcal I$.
By definition of pre-image it follows that $h(s_0)\in A_\alpha$ $\forall\alpha\in\mathcal I$.
Therefore by def of intersection $h(s_0)\in\bigcap_{\alpha\in\mathcal I}A_\alpha$ hence $s_0\in h^{-1}(\bigcap_{\alpha\in\mathcal I}A_\alpha)$.
Since $s_0$ is arbitrary we obtain $\bigcap_{\alpha\in\mathcal I}h^{-1}(A_\alpha)\subseteq h^{-1}(\bigcap_{\alpha\in\mathcal I}A_\alpha)$.
Combining the two inclusions we conclude the proof.
%
% !TEX root = ../main.tex
\subsection*{Lemma 3.1 --- The pre-image of Borel sets form a $\sigma$-algebra}%
\manualtoc{subsection}{\texorpdfstring{Lemma 3.1 --- The pre-image of Borel sets form a $\sigma$-algebra}{Lemma 3.1 --- The pre-image of Borel sets form a sigma-algebra}}{studiabile-nav-0148}%

Let $h:S\to\mathbb R$, we want to prove that
$h^{-1}(\mathcal B(\mathbb R)):=\{h^{-1}(A):A\in\mathcal B(\mathbb R)\}$ is a $\sigma$-algebra on $S$.
\par\noindent\studyheading{Proof.}\par\nobreak
By definition of $\sigma$-algebra we must verify the following three prop.

\studylabel{(i)} $\emptyset\in h^{-1}(\mathcal B(\mathbb R))$.
Notice that $\emptyset,\mathbb R\in\mathcal B(\mathbb R)$ since $\mathcal B(\mathbb R)$ is a $\sigma$-algebra on $\mathbb R$.
Moreover by definition of pre-image $h^{-1}(\emptyset)=\emptyset$ and $h^{-1}(\mathbb R)=S$.
Therefore $\emptyset,S\in h^{-1}(\mathcal B(\mathbb R))$.

\studylabel{(ii)} Closure under complements.
Let $H\in h^{-1}(\mathcal B(\mathbb R))$.
By definition of this class, there exists $A\in\mathcal B(\mathbb R)$ such that $H=h^{-1}(A)$.
Since $\mathcal B(\mathbb R)$ is a $\sigma$-algebra, it is closed under complements.
Therefore $A^c\in\mathcal B(\mathbb R)$.
Now consider the complement of $H$:
\[
H^c=(h^{-1}(A))^c=h^{-1}(A^c)\in h^{-1}(\mathcal B(\mathbb R)),
\]
where the first equality follows from the definition of $H$, the second from prop \studylabel{(ii)} HW~28 and the last membership follows from $A^c\in\mathcal B(\mathbb R)$.
Therefore $H\in h^{-1}(\mathcal B(\mathbb R))\Rightarrow H^c\in h^{-1}(\mathcal B(\mathbb R))$.

\studylabel{(iii)} Closure under countable unions.
Let $(H_n)_{n\in\mathbb N}\subseteq h^{-1}(\mathcal B(\mathbb R))$.
By definition of this class $\forall n\in\mathbb N$ $\exists A_n\in\mathcal B(\mathbb R)$ such that $H_n=h^{-1}(A_n)$.
Since $\mathcal B(\mathbb R)$ is a $\sigma$-alg. it is closed under countable unions therefore $\bigcup_{n=1}^\infty A_n\in\mathcal B(\mathbb R)$.
Now consider the union of the sets $H_n$:
\[
\bigcup_{n=1}^\infty H_n=\bigcup_{n=1}^\infty h^{-1}(A_n)
=h^{-1}\left(\bigcup_{n=1}^\infty A_n\right)\in h^{-1}(\mathcal B(\mathbb R)),
\]
where the first equality follows from the def of $H_n$, the second from prop \studylabel{(i)} of the pre-image mapping and the last membership follows from $\bigcup_{n=1}^\infty A_n\in\mathcal B(\mathbb R)$.
Therefore $(H_n)_{n\in\mathbb N}\subseteq h^{-1}(\mathcal B(\mathbb R))\Rightarrow\bigcup_{n=1}^\infty H_n\in h^{-1}(\mathcal B(\mathbb R))$.
Since all three prop are satisfied $h^{-1}(\mathcal B(\mathbb R))$ is a $\sigma$-alg on $S$.

\studyheading{\inlinetoc{subsubsection}{\texorpdfstring{Rmk:}{Rmk:}}{studiabile-nav-0153}Rmk:} It should be clear that $h$ is measurable respect to $h^{-1}(\mathcal B(\mathbb R))$.
Moreover for any $h:S\to\mathbb R$ we can construct a $\sigma$-alg on $S$ with respect to which $h$ is measurable. The smallest of such $\sigma$-algebras is exactly $h^{-1}(\mathcal B(\mathbb R))$.

\studyheading{\inlinetoc{subsubsection}{\texorpdfstring{Problem:}{Problem:}}{studiabile-nav-0154}Problem:} According to Def 3.1 in order to check whether a function is $\Sigma$-measurable or not one need to verify the pre-images of ALL Borel sets are in $\Sigma$. Checking every Borel set directly is generally impractical; the next lemma provides a generating-class test.
%
% !TEX root = ../main.tex
\subsection*{Lemma 3.2. Measurability through a generating class}%
\manualtoc{subsection}{\texorpdfstring{Lemma 3.2. Measurability through a generating class}{Lemma 3.2. Measurability through a generating class}}{studiabile-nav-0155}%

Let $(S,\Sigma)$ be a measurable space, $h:S\to\mathbb R$ and let $\mathcal C\subseteq\mathcal B(\mathbb R)$ be such that $\sigma(\mathcal C)=\mathcal B(\mathbb R)$.
If $h^{-1}(A)\in\Sigma$ $\forall A\in\mathcal C$ then $h$ is $\Sigma$-measurable i.e. $h^{-1}(A)\in\Sigma$ $\forall A\in\mathcal B(\mathbb R)$.

\par\noindent\studyheading{Proof.}\par\nobreak
Define $\mathcal E:=\{A\in\mathcal B(\mathbb R):h^{-1}(A)\in\Sigma\}$.
By definition of $\mathcal E$ and by our assumption we have $\mathcal C\subseteq\mathcal E\subseteq\mathcal B(\mathbb R)$.
\emph{We want to prove that $\mathcal E$ is a $\sigma$-algebra on $\mathbb R$.}
By definition of $\sigma$-alg we must verify the three prop.

\studylabel{(i)} $\emptyset$ and $\mathbb R$ belong to $\mathcal E$.
Notice that by definition of pre-image $h^{-1}(\emptyset)=\emptyset$ and $h^{-1}(\mathbb R)=S$.
Since $\Sigma$ is a $\sigma$-algebra on $S$ we have $\emptyset,S\in\Sigma$.
Therefore $h^{-1}(\emptyset)\in\Sigma$ and $h^{-1}(\mathbb R)\in\Sigma$ hence by definition of $\mathcal E$, $\emptyset,\mathbb R\in\mathcal E$.

\studylabel{(ii)} Closure under complements.
Let $A\in\mathcal E$, by def of $\mathcal E$ we have $A\in\mathcal B(\mathbb R)$ and $h^{-1}(A)\in\Sigma$.
Since both $\mathcal B(\mathbb R)$ and $\Sigma$ are $\sigma$-algebras they are closed under complements.
Therefore $A^c\in\mathcal B(\mathbb R)$ and $(h^{-1}(A))^c\in\Sigma$.
Now by prop \studylabel{(ii)} of HW~28 $h^{-1}(A^c)=(h^{-1}(A))^c\in\Sigma$ where the equality follows from the property of pre-images of complements and the last membership follows from the closure of $\Sigma$ under complements.
Since $A^c\in\mathcal B(\mathbb R)$ and $h^{-1}(A^c)\in\Sigma$ by definition of $\mathcal E$, $A^c\in\mathcal E$.
Hence $\mathcal E$ is closed under complements.

\studylabel{(iii)} Closure under countable unions.
Let $(A_n)_{n\in\mathbb N}\subseteq\mathcal E$, by definition of $\mathcal E$, $A_n\in\mathcal B(\mathbb R)$ and $h^{-1}(A_n)\in\Sigma$ $\forall n\in\mathbb N$.
Since $\mathcal B(\mathbb R)$ is a $\sigma$-algebra, it is closed under countable unions.
Therefore $\bigcup_{n=1}^\infty A_n\in\mathcal B(\mathbb R)$.
Moreover by property \studylabel{(i)} of the pre-image mapping
\[
h^{-1}\left(\bigcup_{n=1}^\infty A_n\right)=\bigcup_{n=1}^\infty h^{-1}(A_n)\in\Sigma
\]
where the equality follows from property \studylabel{(i)} of pre-images and the last membership follows from the closure of $\Sigma$ under countable unions, since $h^{-1}(A_n)\in\Sigma$ for every $n\in\mathbb N$.
Hence by definition of $\mathcal E$, $\bigcup_{n=1}^\infty A_n\in\mathcal E$, thus $\mathcal E$ is closed under countable unions.
Since all three prop are satisfied $\mathcal E$ is a $\sigma$-alg on $\mathbb R$.
Moreover since $\mathcal C\subseteq\mathcal E$ and $\mathcal E$ is a $\sigma$-algebra, by the definition of the generated $\sigma$-algebra $\sigma(\mathcal C)\subseteq\mathcal E$.
Therefore $\mathcal B(\mathbb R)=\sigma(\mathcal C)\subseteq\mathcal E\subseteq\mathcal B(\mathbb R)$ hence $\mathcal E=\mathcal B(\mathbb R)$.
By def of $\mathcal E$ this implies $h^{-1}(A)\in\Sigma$ $\forall A\in\mathcal B(\mathbb R)$ therefore $h\in\mathrm m\Sigma$.
%
% !TEX root = ../main.tex
\subsection*{Remarks}%
\manualtoc{subsection}{\texorpdfstring{Remarks}{Remarks}}{studiabile-nav-0160}%

There are some important consequences of the above lemma:

\par\noindent\studyheading{Proof.}\par\nobreak
\studylabel{(i)} If $S$ is a topological space and $h:S\to\mathbb R$ is continuous then $h$ is Borel measurable (i.e. $h^{-1}(A)\in\mathcal B(S)$ $\forall A\in\mathcal B(\mathbb R)$).
Indeed, $h^{-1}(A)$ is an open set in $S$ whenever $A$ is an open set in $\mathbb R$, due to continuity.
Take the generating class of open sets in $\mathbb R$. Their pre-images are open in $S$, hence in $\mathcal B(S)$, and they generate $\mathcal B(\mathbb R)$. Lemma~3.2 therefore applies with domain sigma-algebra $\mathcal B(S)$.

\studylabel{(ii)} Given a measurable space $(S,\Sigma)$ a function $h:S\to\mathbb R$ is $\Sigma$-measurable if and only if $\{h\le c\}:=\{s\in S:h(s)\in(-\infty,c]\}=h^{-1}((-\infty,c])\in\Sigma$ for all $c\in\mathbb R$.
Indeed, if $h\in\mathrm m\Sigma$, then $h^{-1}((-\infty,c])\in\Sigma$ for all $c\in\mathbb R$ because $(-\infty,c]\in\mathcal B(\mathbb R)$.
If instead $h^{-1}((-\infty,c])\in\Sigma$ for all $c\in\mathbb R$, then we apply the lemma upon recalling the $\pi$-system $\pi(\mathbb R)$ and the fact that $\sigma(\pi(\mathbb R))=\mathcal B(\mathbb R)$.
This is a clever way to test measurability.
%
% !TEX root = ../main.tex
\subsection*{HW 29 --- A Borel indicator}%
\manualtoc{subsection}{\texorpdfstring{HW 29 --- A Borel indicator}{HW 29 --- A Borel indicator}}{studiabile-nav-0164}%

Let $S=\mathbb R$, $\Sigma=\mathcal B(\mathbb R)$ and $h(s)=\mathbf{1}_{[0,1]}(s)$. Prove that $h$ is Borel-measurable. Hint: first you need to show that $[0,1]\in\mathcal B(\mathbb R)$.
\par\noindent\studyheading{Proof.}\par\nobreak
First, $[0,1]=\mathbb R\setminus(( -\infty,0)\cup(1,+\infty))\in\mathcal B(\mathbb R)$, since the set being complemented is open. Recall that by consequence of \studylabel{(ii)} of Lemma 3.2 a function $h:\mathbb R\to\mathbb R$ is Borel measurable if and only if $h^{-1}((-\infty,c])\in\mathcal B(\mathbb R)$ $\forall c\in\mathbb R$.
By definition of indicator function $h(s)=\begin{cases}0&s\notin[0,1],\\1&s\in[0,1].\end{cases}$
Therefore by definition of pre-image
\[
h^{-1}((-\infty,c])=\{s\in\mathbb R:h(s)\le c\}
=\begin{cases}\emptyset&c<0,\\(-\infty,0)\cup(1,+\infty)&0\le c<1,\\\mathbb R&c\ge1.\end{cases}
\]
Indeed if $c<0$ neither $0$ nor $1$ satisfies $h(s)\le c$ so the pre-image is empty.
If $0\le c<1$ only $h(s)=0$ satisfies the inequality which occurs when $s\notin[0,1]$.
Finally if $c\ge1$, both value satisfies the inequality so the pre-image is $\mathbb R$.
Now since $\mathcal B(\mathbb R)$ is a $\sigma$-algebra we have $\emptyset,\mathbb R\in\mathcal B(\mathbb R)$.
Moreover since $(-\infty,0)\cup(1,+\infty)$ are open sets, their union is also open.
Therefore by definition of $\sigma$-algebra $(-\infty,0)\cup(1,+\infty)\in\mathcal B(\mathbb R)$.
Hence in all three cases $h^{-1}((-\infty,c])\in\mathcal B(\mathbb R)$ $\forall c\in\mathbb R$, therefore by consequence \studylabel{(ii)} of Lemma 3.2 $h\in\mathrm m\mathcal B(\mathbb R)$.
%
% !TEX root = ../main.tex
\subsection*{HW 30 --- Equivalent threshold tests}%
\manualtoc{subsection}{\texorpdfstring{HW 30 --- Equivalent threshold tests}{HW 30 --- Equivalent threshold tests}}{studiabile-nav-0166}%

Prove that $h:S\to\mathbb R$ is $\Sigma$-measurable if and only if each of the following holds:
\[
\begin{aligned}
\text{\studylabel{(i)}}\quad&\{h<c\}:=h^{-1}((-\infty,c))\in\Sigma\quad\forall c\in\mathbb R,\\
\text{\studylabel{(ii)}}\quad&\{h>c\}:=h^{-1}((c,+\infty))\in\Sigma\quad\forall c\in\mathbb R,\\
\text{\studylabel{(iii)}}\quad&\{h\ge c\}:=h^{-1}([c,+\infty))\in\Sigma\quad\forall c\in\mathbb R.
\end{aligned}
\]
\par\noindent\studyheading{Proof.}\par\nobreak
\noindent\studyheading{($\Rightarrow$)} We assume that $h$ is $\Sigma$-measurable then by Definition 3.1 $\forall A\in\mathcal B(\mathbb R),\ h^{-1}(A)\in\Sigma$.
Since $(-\infty,c)$ and $(c,+\infty)$ are open sets they are Borel and $[c,+\infty)=(-\infty,c)^c$ hence is Borel so their pre-image belong to $\Sigma$.

\noindent\studyheading{($\Leftarrow$)}

\studylabel{(i)} Assume that $\{h<c\}\in\Sigma$ $\forall c\in\mathbb R$.
Notice that $(-\infty,c]=\bigcap_{n=1}^\infty(-\infty,c+\frac1n)$.
Therefore
\[
\begin{aligned}
\{h\le c\}&=h^{-1}((-\infty,c])
=h^{-1}\left(\bigcap_{n=1}^\infty(-\infty,c+\tfrac1n)\right)\\
&=\bigcap_{n=1}^\infty h^{-1}((-\infty,c+\tfrac1n))
=\bigcap_{n=1}^\infty\{h<c+\tfrac1n\}\in\Sigma,
\end{aligned}
\]
where the first equality follows from the definition of pre-image, the second from the identity above, the third from property \studylabel{(iii)} of HW~28 and the fourth from the definition of $\{h<c\}$.
The last membership follows from the closure under countable intersections since $\{h<c+\frac1n\}\in\Sigma$ $\forall n\in\mathbb N$ by assumption.
Hence $\{h\le c\}\in\Sigma$ $\forall c\in\mathbb R$ and by Lemma 3.2 $h$ is $\Sigma$-measurable.

\studylabel{(ii)} Assume that $\{h>c\}\in\Sigma$ $\forall c\in\mathbb R$.
Notice that $(-\infty,c]=(c,+\infty)^c$.
Therefore
\[
\begin{aligned}
\{h\le c\}&=h^{-1}((-\infty,c])=h^{-1}((c,+\infty)^c) =(h^{-1}((c,+\infty)))^c=\{h>c\}^c\in\Sigma,
\end{aligned}
\]
where the first equality follows from the definition of pre-image, the second from the complement identity, the third from prop \studylabel{(ii)} of HW~28, the fourth from the definition of $\{h>c\}$.
The last membership follows closure of $\Sigma$ under complements since $\{h>c\}\in\Sigma$ by assumption.
Hence by Lemma 3.2 $h$ is $\Sigma$-measurable.

\studylabel{(iii)} Assume that $\{h\ge c\}\in\Sigma$ $\forall c\in\mathbb R$.
Notice that $(-\infty,c)=[c,+\infty)^c$.
Therefore
\[
\begin{aligned}
\{h<c\}&=h^{-1}((-\infty,c))=h^{-1}([c,+\infty)^c) =(h^{-1}([c,+\infty)))^c=\{h\ge c\}^c\in\Sigma,
\end{aligned}
\]
where the first equality follows from the definition of pre-image, the second from the complement identity, the third from property \studylabel{(ii)} of HW~28 and the fourth from the definition of $\{h\ge c\}$.
The last membership follows from the closure of $\Sigma$ under complements since $\{h\ge c\}\in\Sigma$ by assumption.
Hence condition \studylabel{(i)} holds and by the first part of the proof $h$ is $\Sigma$-measurable.
Therefore each of the three condition is equivalent to $\Sigma$-measurability.
%
% !TEX root = ../main.tex
\subsection*{Lemma 3.3. Composition of functions}%
\manualtoc{subsection}{\texorpdfstring{Lemma 3.3. Composition of functions}{Lemma 3.3. Composition of functions}}{studiabile-nav-0174}%

Let $h:S\to\mathbb R$ and $f:\mathbb R\to\mathbb R$ be such that $h\in\mathrm m\Sigma$ and $f\in\mathrm m\mathcal B$.
Then $g:=f\circ h:S\to\mathbb R$ (i.e. $g(s)=f(h(s))$, $s\in S$) is a $\Sigma$-measurable functions.
\par\noindent\studyheading{Proof.}\par\nobreak
For any $A\in\mathcal B(\mathbb R)$ we have
\[
\begin{aligned}
g^{-1}(A)&=\{s\in S:f(h(s))\in A\}
=\{s\in S:h(s)\in f^{-1}(A)\} =h^{-1}(f^{-1}(A))\in\Sigma,
\end{aligned}
\]
The first equality uses $g(s)=f(h(s))$, the second the definition of the pre-image under $f$, and the third that under $h$. Since $f^{-1}(A)$ is Borel by measurability of $f$, measurability of $h$ gives the final membership.
%
% !TEX root = ../main.tex
\subsection*{Lemma 3.4}%
\manualtoc{subsection}{\texorpdfstring{Lemma 3.4}{Lemma 3.4}}{studiabile-nav-0176}%

If $h,h_1,h_2\in\mathrm m\Sigma$ then \studylabel{(i)} $h_1+h_2\in\mathrm m\Sigma$, \studylabel{(ii)} $h_1\cdot h_2\in\mathrm m\Sigma$, \studylabel{(iii)} $\lambda h\in\mathrm m\Sigma$ for any $\lambda\in\mathbb R$.
\par\noindent\studyheading{Proof.}\par\nobreak
\studylabel{(i)} Recall by HW~30 it is sufficient to show that $\{h_1+h_2>c\}\in\Sigma$ $\forall c\in\mathbb R$.
Notice that if $h_1(s)+h_2(s)>c$ then $c-h_2(s)<h_1(s)$.
Since $\mathbb Q$ is dense in $\mathbb R$ there exists $q\in\mathbb Q$ such that $h_1(s)>q>c-h_2(s)$.
Therefore
\[
\begin{aligned}
\{h_1+h_2>c\}
&=\bigcup_{q\in\mathbb Q}\{s\in S:h_1(s)>q>c-h_2(s)\} =\bigcup_{q\in\mathbb Q}(\{h_1>q\}\cap\{h_2>c-q\})\in\Sigma.
\end{aligned}
\]
Since $h_1,h_2$ are $\Sigma$-measurable by HW~30 we have $\{h_1>q\}\in\Sigma$ and $\{h_2>c-q\}\in\Sigma$.
The second equality rewrites $q>c-h_2(s)$ as $h_2(s)>c-q$. Finite-intersection closure and then countable-union closure of $\Sigma$, using countability of $\mathbb Q$, give the final membership.
The rational-union identity in \studylabel{(i)} also has the reverse implication: if $h_1(s)>q$ and $h_2(s)>c-q$, adding gives $h_1(s)+h_2(s)>c$. Thus no point is added by that union. Since $c\in\mathbb R$ was arbitrary, HW~30 gives $h_1+h_2\in\mathrm m\Sigma$.

\studylabel{(iii)} By HW~30 it is sufficient to prove that $\{\lambda h>c\}\in\Sigma$ $\forall c\in\mathbb R$. We distinguish three cases.
Case 1. $\lambda=0$. Notice that $\{\lambda h>c\}=\{s\in S:0>c\}=\begin{cases}\emptyset&c\ge0,\\S&c<0.\end{cases}$
Since $\emptyset,S\in\Sigma$ it follows that $\{\lambda h>c\}\in\Sigma$.
Case 2. $\lambda>0$. Notice that $\{\lambda h>c\}=\{s\in S:\lambda h(s)>c\}=\{s\in S:h(s)>\frac c\lambda\}\in\Sigma$ since $h\in\mathrm m\Sigma$.
Case 3. $\lambda<0$. Notice that $\{\lambda h>c\}=\{s\in S:h(s)<\frac c\lambda\}\in\Sigma$ by HW~30.
Hence in all three cases $\{\lambda h>c\}\in\Sigma$ therefore by HW~30 $\lambda h\in\mathrm m\Sigma$.

\studylabel{(ii)} Let $h_1,h_2\in\mathrm m\Sigma$. Notice that the function $f(x)=x^2$ is continuous and therefore is Borel-measurable by consequence \studylabel{(i)} of Lemma 3.2.
Hence by Lemma 3.3 $h_1^2,h_2^2\in\mathrm m\Sigma$ and, applying the same composition lemma to the measurable sum from \studylabel{(i)}, $(h_1+h_2)^2\in\mathrm m\Sigma$.
Now notice that $h_1h_2=\frac12(h_1+h_2)^2-\frac12(h_1^2+h_2^2)$, since sum, difference and scalar multiples are $\Sigma$-measurable also $h_1h_2\in\mathrm m\Sigma$.
%
% !TEX root = ../main.tex
\subsection*{Lemma 3.5. Measurability of infimum, supremum, liminf and limsup}%
\manualtoc{subsection}{\texorpdfstring{Lemma 3.5. Measurability of infimum, supremum, liminf and limsup}{Lemma 3.5. Measurability of infimum, supremum, liminf and limsup}}{studiabile-nav-0181}%

Let $(S,\Sigma)$ be a measurable space and let $(h_n)_{n\in\mathbb N}\subseteq\mathrm m\Sigma$. The infimum, supremum and limits below may take values in $\overline{\mathbb R}=[-\infty,+\infty]$. For a map $u:S\to\overline{\mathbb R}$, measurability means that $\{u\le c\}\in\Sigma$ for every $c\in\mathbb R$; the equivalent threshold tests from HW~30 apply in this extended setting as well. Memberships in $\mathrm m\Sigma$ within this lemma are understood in this extended sense. Whenever one of these functions is finite at every point, it belongs to the real-valued class of Definition~3.1. Then
\studylabel{(i)} $\inf_{n\in\mathbb N}h_n\in\mathrm m\Sigma$, $\sup_{n\in\mathbb N}h_n\in\mathrm m\Sigma$;
\studylabel{(ii)} $\liminf_{n\to+\infty}h_n\in\mathrm m\Sigma$, $\limsup_{n\to+\infty}h_n\in\mathrm m\Sigma$.

\par\noindent\studyheading{Proof.}\par\nobreak
\studylabel{(i)} Let $h(s):=\inf_{n\in\mathbb N}h_n(s)$.
Recall that by HW~30 to prove that $h$ is $\Sigma$-measurable it is sufficient to show that $\{h\ge c\}\in\Sigma$ $\forall c\in\mathbb R$.
As we will see in HW~31, $\{h\ge c\}=\bigcap_{n=1}^\infty\{h_n\ge c\}$. Each set in the intersection is measurable by HW~30, so closure under countable intersections gives $\{h\ge c\}\in\Sigma$.
By HW~22 $\sup_{n\in\mathbb N}h_n=-\inf_{n\in\mathbb N}(-h_n)$, since $h_n\in\mathrm m\Sigma$ by Lemma 3.4 \studylabel{(iii)}, $-h_n\in\mathrm m\Sigma$ and by the first part of the proof $\inf_{n\in\mathbb N}(-h_n)\in\mathrm m\Sigma$. Negation also preserves extended-real measurability, since $\{-u>c\}=\{u<-c\}$ for every real threshold $c$. Consequently $-\inf_{n\in\mathbb N}(-h_n)=\sup_{n\in\mathbb N}h_n$ is measurable.

\studylabel{(ii)} By definition of liminf $\liminf_{n\to+\infty}h_n(s)=\sup_{n\ge1}\inf_{m\ge n}h_m(s)$.
For every $n\in\mathbb N$ define $\ell_n(s):=\inf_{m\ge n}h_m(s)$.
Notice that $\ell_n$ is the infimum of a countable family of $\Sigma$-measurable functions therefore by prop \studylabel{(i)} $\ell_n\in\mathrm m\Sigma$ $\forall n\in\mathbb N$.
Therefore also $\sup_{n\ge1}\ell_n(s)\in\mathrm m\Sigma$ thanks to \studylabel{(i)} hence $\liminf h_n(s)\in\mathrm m\Sigma$.
By HW~22 $\limsup_{n\to+\infty}h_n(s)=-\liminf_{n\to+\infty}(-h_n(s))$ and the same threshold test for negation shows that $\limsup_{n\to+\infty}h_n$ is measurable.
%
% !TEX root = ../main.tex
\subsection*{HW 31 --- Superlevel sets of an infimum}%
\manualtoc{subsection}{\texorpdfstring{HW 31 --- Superlevel sets of an infimum}{HW 31 --- Superlevel sets of an infimum}}{studiabile-nav-0187}%

For $h=\inf_n h_n$, prove that $\{h\ge c\}=\bigcap_{n=1}^\infty\{h_n\ge c\}$.
\par\noindent\studyheading{Proof.}\par\nobreak
\emph{We first prove that $\{h\ge c\}\subseteq\bigcap_{n=1}^\infty\{h_n\ge c\}$.}
Let $s_0\in\{h\ge c\}$.
By definition of this set we have $h(s_0)=\inf_{n\in\mathbb N}h_n(s_0)\ge c$.
By definition of infimum $\inf_{n\in\mathbb N}h_n(s_0)\le h_n(s_0)$ $\forall n\in\mathbb N$.
Therefore $h_n(s_0)\ge\inf_{n\in\mathbb N}h_n(s_0)\ge c$ $\forall n\in\mathbb N$ hence $s_0\in\{h_n\ge c\}$ which implies $s_0\in\bigcap_{n=1}^\infty\{h_n\ge c\}$.
Since $s_0$ was arbitrary $\bigcap_{n=1}^\infty\{h_n\ge c\}\supseteq\{h\ge c\}$.
For the reverse inclusion let $s_0\in\bigcap_{n=1}^\infty\{h_n\ge c\}$.
By definition of intersection we have $s_0\in\{h_n\ge c\}$ $\forall n\in\mathbb N$.
Therefore by definition of these sets $h_n(s_0)\ge c$ $\forall n\in\mathbb N$.
Notice that $c$ is a lower bound for the family $\{h_n(s_0),n\in\mathbb N\}$.
Since the infimum is the greatest lower bound we obtain $h(s_0)=\inf_{n\in\mathbb N}h_n(s_0)\ge c$ hence $s_0\in\{h\ge c\}$.
Since $s_0$ was arbitrary $\bigcap_{n=1}^\infty\{h_n\ge c\}\subseteq\{h\ge c\}$.
Combining the two inclusions we conclude the proof.
%
% !TEX root = ../main.tex
\subsection*{Corollary 3.6 --- Measurability of the Convergence Set}%
\manualtoc{subsection}{\texorpdfstring{Corollary 3.6 --- Measurability of the Convergence Set}{Corollary 3.6 --- Measurability of the Convergence Set}}{studiabile-nav-0189}%

Let $(h_n)_{n\in \mathbb{N}}\subset m\Sigma$ . Then
$\{s \in \mathcal {S}: \lim _ {n \to \infty} h _ {n} (s) \text{exists} \} \in \Sigma .$
 
\begin{proof}[Expanded proof]
Let
\[
\underline h =\liminf_{n\to\infty}h_n,
\; \overline h =\limsup_{n\to\infty}h_n.
\]
By Lemma 3.5 these are extended-real measurable, and
$\underline h\le\overline h$ pointwise. We follow the lecturer's
idea of testing their difference at zero, first restricting to where
the subtraction is defined.

\noindent\studyheading{Part 1 --- Both liminf and limsup are finite real numbers.}
Let $C=\{s\in\mathcal S:\underline h(s),\overline h(s)\in\R\}$.
This set is measurable: since $\underline h\le\overline h$,
$C=\bigcup_{k\ge1}(\{\underline h\ge-k\}\cap\{\overline h\le k\})\in\Sigma$.
On $C$ both functions are real-valued, so we can use the lecturer's subtraction argument.
Define $g(s):=\underline h(s)-\overline h(s)$ for $s\in C$.
By Lemma~3.5 the two functions are measurable, and Lemma~3.4\studylabel{(i)},\studylabel{(iii)}
makes their difference $g$ measurable on $C$.

\emph{We show that finite real convergence is equivalent to $g=0$.}
A real sequence converges to a finite limit if and only if its liminf and limsup
coincide in $\R$. Since both values are finite on $C$, this is equivalent to
their difference being zero. Therefore
\[
E_{\mathrm{fin}}=\{s\in C:\underline h(s)-\overline h(s)=0\}
=\{s\in C:g(s)=0\}=g^{-1}(\{0\})\in\Sigma.
\]
The first equality is the finite liminf/limsup convergence criterion,
the second substitutes the definition of $g$, and the third is the definition
of pre-image. Since $\{0\}\in\mathcal B(\R)$ and $g$ is measurable on the
measurable set $C$, its pre-image belongs to $\Sigma$.

\noindent\studyheading{Part 2 --- Both liminf and limsup equal $+\infty$.}
\emph{We show that convergence to $+\infty$ is a measurable event.}
The inequality $\underline h\le\overline h$ implies that
$\underline h(s)=+\infty$ forces $\overline h(s)=+\infty$.
Consequently,
\[
E_+:=\{s:\underline h(s)=\overline h(s)=+\infty\}
=\{\underline h=+\infty\}
=\bigcap_{k\ge1}\{\underline h>k\}\in\Sigma.
\]
The first equality uses the inequality between liminf and limsup.
For the second, an extended real number equals $+\infty$ exactly when it exceeds
every positive integer. Each level set is measurable by Lemma~3.5;
closure under countable intersections gives the final membership.
If $\underline h(s)=+\infty$, the identity
$\underline h(s)=\sup_N\inf_{n\ge N}h_n(s)$ shows that, for every real $M$,
some tail infimum exceeds $M$. Hence $h_n(s)>M$ throughout that tail,
which is convergence to $+\infty$.
Conversely, convergence to $+\infty$ puts $h_n(s)>M+1$ eventually for every $M$.
The corresponding tail infimum is at least $M+1$, so $\underline h(s)>M$ for every $M$,
and therefore $\underline h(s)=+\infty$.

\noindent\studyheading{Part 3 --- Both liminf and limsup equal $-\infty$.}
\emph{We show that convergence to $-\infty$ is a measurable event.}
Again $\underline h\le\overline h$, so $\overline h(s)=-\infty$ forces
$\underline h(s)=-\infty$. Thus
\[
E_-:=\{s:\underline h(s)=\overline h(s)=-\infty\}
=\{\overline h=-\infty\}
=\bigcap_{k\ge1}\{\overline h<-k\}\in\Sigma.
\]
The first equality uses the same order inequality. The second characterizes $-\infty$
by all negative integer thresholds. Each level set is measurable by Lemma~3.5,
so its countable intersection belongs to $\Sigma$.
If $\overline h(s)=-\infty$, the identity
$\overline h(s)=\inf_N\sup_{n\ge N}h_n(s)$ gives, for every real $M$,
a tail supremum below $M$. Every term of that tail is then below $M$,
which is convergence to $-\infty$.
Conversely, that convergence gives $h_n(s)<M-1$ eventually for every $M$.
The tail supremum is at most $M-1$, so $\overline h(s)<M$ for every $M$,
and therefore $\overline h(s)=-\infty$.

Combining the three parts, the event of convergence in $\overline{\R}$ is
\[
\{s:\lim_{n\to\infty}h_n(s)\text{ exists in }\overline{\R}\}
=E_{\mathrm{fin}}\cup E_+\cup E_-\in\Sigma.
\]
For convergence to a finite real number, the event is $E_{\mathrm{fin}}$ alone.
No difference of two infinite values is used.
\end{proof}
%
% !TEX root = ../main.tex
\subsection*{Def: Random variable}%
\manualtoc{subsection}{\texorpdfstring{Def: Random variable}{Def: Random variable}}{studiabile-nav-0194}%

Given a probability space $(\Omega,\mathcal F,\mathbb P)$ a r.v. $X$ is a $\mathcal F$-measurable function i.e. $X:\Omega\to\mathbb R$ and $X^{-1}:\mathcal B(\mathbb R)\to\mathcal F$.
%
% !TEX root = ../main.tex
\subsection*{Ex 1b proof}%
\manualtoc{subsection}{\texorpdfstring{Example 1(b) --- Frequency event}{Example 1(b) --- Frequency event}}{studiabile-nav-0195}%

Consider the infinite coin tossing where $\Omega=\{H,T\}^{\mathbb N}$ and
$\mathcal F=\sigma(\{\omega\in\Omega:\omega_n=W\},\ n\in\mathbb N,\ W\in\{H,T\})$.
For every $n\in\mathbb N$ define $X_n(\omega)=\begin{cases}1&\text{if }\omega_n=H,\\0&\text{if }\omega_n=T,\end{cases}$ and $S_n=X_1+\cdots+X_n$.
Prove that $F:=\{\omega\in\Omega:\lim_{n\to+\infty}\frac{S_n(\omega)}n=\frac12\}\in\mathcal F$.
\par\noindent\studyheading{Proof.}\par\nobreak
\emph{We first prove that $X_n$ is $\mathcal F$-measurable $\forall n\in\mathbb N$.}
By HW~30 it is sufficient to show that $\{X_n\le c\}\in\mathcal F$ $\forall c\in\mathbb R$.
By definition of $X_n$ we have
\[
\{X_n\le c\}=\begin{cases}\emptyset&c<0,\\\{\omega\in\Omega:\omega_n=T\}&0\le c<1,\\\Omega&c\ge1.\end{cases}
\]
Notice that $\emptyset,\Omega\in\mathcal F$ and $\{\omega\in\Omega:\omega_n=T\}\in\mathcal F$ by def of $\mathcal F$ in example 1b.
Therefore $\{X_n\le c\}\in\mathcal F$ $\forall c\in\mathbb R$ hence $X_n\in\mathrm m\mathcal F$ $\forall n\in\mathbb N$.
Now by Lemma 3.4 \studylabel{(i)} $S_n:=\sum_{k=1}^n X_k\in\mathrm m\mathcal F$ $\forall n\in\mathbb N$.
Moreover since $\frac1n\in\mathbb R$, by Lemma 3.4 \studylabel{(iii)} $\frac1n S_n\in\mathrm m\mathcal F$ $\forall n\in\mathbb N$.
\emph{We prove now that $F\in\mathcal F$.}
Recall that a sequence of real numbers converges to $a\in\mathbb R$ if and only if its limsup and liminf are both equal to $a$.
Therefore
\[
\begin{aligned}
F&=\{\omega\in\Omega:\lim_{n\to+\infty}\tfrac{S_n(\omega)}n=\tfrac12\} =\{\omega\in\Omega:\limsup_{n\to+\infty}\tfrac{S_n(\omega)}n=\tfrac12\}
\cap\{\omega\in\Omega:\liminf_{n\to+\infty}\tfrac{S_n(\omega)}n=\tfrac12\}.
\end{aligned}
\]
Define $U(\omega):=\limsup_{n\to+\infty}\frac{S_n(\omega)}n$ and $L(\omega):=\liminf_{n\to+\infty}\frac{S_n(\omega)}n$.
Since $\frac{S_n}n\in\mathrm m\mathcal F$ $\forall n\in\mathbb N$ by Lemma 3.5 \studylabel{(ii)} $U,L\in\mathrm m\mathcal F$.
Notice that both functions are real valued since $0\le\frac{S_n(\omega)}n\le1$ $\forall\omega$.
Moreover $\{\frac12\}\in\mathcal B(\mathbb R)$ since it is a singleton.
In particular
\[
\begin{aligned}
A&:=\{\omega\in\Omega:\limsup_{n\to+\infty}\tfrac{S_n(\omega)}n=\tfrac12\}
=U^{-1}(\{\tfrac12\})\in\mathcal F,\\
B&:=\{\omega\in\Omega:\liminf_{n\to+\infty}\tfrac{S_n(\omega)}n=\tfrac12\}
=L^{-1}(\{\tfrac12\})\in\mathcal F.
\end{aligned}
\]
Hence recalling $F=A\cap B$ since $\mathcal F$ is closed under intersection we conclude that $F\in\mathcal F$.
As a consequence, since $\mathbb P:\mathcal F\to[0,1]$ is a probability measure and $F\in\mathcal F$, the probability $\mathbb P(F)$ is well-defined.
%
% !TEX root = ../main.tex
\subsection*{Def Law of a random variable}%
\manualtoc{subsection}{\texorpdfstring{Def Law of a random variable}{Def Law of a random variable}}{studiabile-nav-0197}%

On a probability space $(\Omega,\mathcal F,\mathbb P)$ let $X$ be a real-valued random variable.
The law of $X$ is a probability measure on $\mathcal B(\mathbb R)$ i.e. a mapping $\mathcal L_X:\mathcal B(\mathbb R)\to[0,1]$ defined as $\mathcal L_X:=\mathbb P\circ X^{-1}$.
That is, for any $A\in\mathcal B(\mathbb R)$ $\mathcal L_X(A):=\mathbb P(X\in A)$.
The diagram below should help your intuition:
$\mathcal B(\mathbb R)\xrightarrow{X^{-1}}\mathcal F\xrightarrow{\mathbb P}[0,1]$.

\par\noindent\studyheading{Proof.}\par\nobreak
Notice that $\mathcal L_X$ is well defined. Indeed since $X$ is $\mathcal F$-measurable
by definition 3.1, $X^{-1}(A)\in\mathcal F$ $\forall A\in\mathcal B(\mathbb R)$.
Therefore $\mathcal L_X(A)=\mathbb P(X^{-1}(A))\in[0,1]$.
\emph{We now verify the properties of a probability measure.}

\emph{\studylabel{(i)} We first prove that $\mathcal L_X(\emptyset)=0$.}
We obtain
\[
\mathcal L_X(\emptyset)=\mathbb P(X^{-1}(\emptyset))=\mathbb P(\emptyset)=0
\]
where the first equality follows from the definition of the law $X$,
the second from the fact that pre-image of the empty set is empty,
the last from properties of the probability measure $\mathbb P$.

\emph{\studylabel{(ii)} We now prove that $\mathcal L_X(\mathbb R)=1$.}
We obtain
\[
\mathcal L_X(\mathbb R)=\mathbb P(X^{-1}(\mathbb R))=\mathbb P(\Omega)=1
\]
where the first equality is the definition of $\mathcal L_X$, the second follows from the fact that $X:\Omega\to\mathbb R$ so $X(\omega)\in\mathbb R$
$\forall\omega\in\Omega$ and the last equality follows from the fact that $\mathbb P$
is a probability measure.

\emph{\studylabel{(iii)} Finally we prove countable additivity.}
Let $(A_n)_{n\in\mathbb N}\subset\mathcal B(\mathbb R)$ be a sequence of pairwise disjoint sets
namely $A_i\cap A_j=\emptyset$ $\forall i\ne j$.
\emph{We want to prove that
\[
\mathcal L_X\left(\bigcup_{n=1}^{\infty}A_n\right)=\sum_{n=1}^{\infty}\mathcal L_X(A_n).
\]}
As we will see in HW~32, $X^{-1}(\bigcup_{n=1}^{\infty}A_n)=\bigcup_{n=1}^{\infty}X^{-1}(A_n)$. We use that identity without proving it here.
First notice that the pre-images $X^{-1}(A_n)$ are also pairwise disjoint.
Indeed for $i\ne j$ by property \studylabel{(iii)} of HW~28
\[
X^{-1}(A_i)\cap X^{-1}(A_j)=X^{-1}(A_i\cap A_j)=X^{-1}(\emptyset)=\emptyset
\]
Moreover since $X$ is $\mathcal F$-measurable $X^{-1}(A_n)\in\mathcal F$ $\forall n\in\mathbb N$.
Therefore we can apply the countable additivity of $\mathbb P$ to the sequence
$(X^{-1}(A_n))_{n\in\mathbb N}$. We obtain
\[
\begin{aligned}
\mathcal L_X\left(\bigcup_{n=1}^{\infty}A_n\right)
&=\mathbb P\left(X^{-1}\left(\bigcup_{n=1}^{\infty}A_n\right)\right)
=\mathbb P\left(\bigcup_{n=1}^{\infty}X^{-1}(A_n)\right) =\sum_{n=1}^{\infty}\mathbb P(X^{-1}(A_n))=\sum_{n=1}^{\infty}\mathcal L_X(A_n)
\end{aligned}
\]
where the first equality follows from the definition of $\mathcal L_X$,
the second from HW~32, the third from the countable additivity of $\mathbb P$
since the pre-images are pairwise disjoint, and the last from the definition of $\mathcal L_X$.
Hence $\mathcal L_X(\bigcup_{n=1}^{\infty}A_n)=\sum_{n=1}^{\infty}\mathcal L_X(A_n)$.
Therefore $\mathcal L_X$ is countably additive.
Since $\mathcal L_X(\emptyset)=0$, $\mathcal L_X(\mathbb R)=1$ and $\mathcal L_X$ is countably additive
on $\mathcal B(\mathbb R)$ we conclude that $\mathcal L_X$ is a probability measure on $\mathcal B(\mathbb R)$.
%
% !TEX root = ../main.tex
\subsection*{HW 32 --- Pre-image of a countable union}%
\manualtoc{subsection}{\texorpdfstring{HW 32 --- Pre-image of a countable union}{HW 32 --- Pre-image of a countable union}}{studiabile-nav-0202}%

Prove that if $(A_n)_{n\in\mathbb N}$ is a family of pairwise disjoint subsets of $\mathbb R$ then
\[
\{\omega\in\Omega:X(\omega)\in\bigcup_{n=1}^\infty A_n\}
=\bigcup_{n=1}^\infty\{\omega\in\Omega:X(\omega)\in A_n\}.
\]
\par\noindent\studyheading{Proof.}\par\nobreak
We first prove that
$\{\omega\in\Omega:X(\omega)\in\bigcup_{n=1}^\infty A_n\}\subseteq\bigcup_{n=1}^\infty\{\omega\in\Omega:X(\omega)\in A_n\}$.
Let $\omega_0\in\{\omega\in\Omega:X(\omega)\in\bigcup_{n=1}^\infty A_n\}$.
By definition of this set we have $X(\omega_0)\in\bigcup_{n=1}^\infty A_n$.
By definition of union there exists $n_0\in\mathbb N$ such that $X(\omega_0)\in A_{n_0}$.
$n_0$ is unique since $A_n$ are pairwise disjoint.
Therefore $\omega_0\in\{\omega\in\Omega:X(\omega)\in A_{n_0}\}\subseteq\bigcup_{n=1}^\infty\{\omega\in\Omega:X(\omega)\in A_n\}$
where the first membership follows from the definition of the set and the inclusion follows from the fact that $\{\omega\in\Omega:X(\omega)\in A_{n_0}\}$ is one of the sets of the union.
Since $\omega_0$ was arbitrary we obtain
$\{\omega\in\Omega:X(\omega)\in\bigcup_{n=1}^\infty A_n\}\subseteq\bigcup_{n=1}^\infty\{\omega\in\Omega:X(\omega)\in A_n\}$.
For the reverse inclusion we let $\omega_0\in\bigcup_{n=1}^\infty\{\omega\in\Omega:X(\omega)\in A_n\}$.

By definition of union there exists $n_0\in\mathbb N$ such that $\omega_0\in\{\omega\in\Omega:X(\omega)\in A_{n_0}\}$.
By definition of this set we have $X(\omega_0)\in A_{n_0}$.
Therefore $X(\omega_0)\in A_{n_0}\subseteq\bigcup_{n=1}^\infty A_n$.
Hence $\omega_0\in\{\omega\in\Omega:X(\omega)\in\bigcup_{n=1}^\infty A_n\}$.
Since $\omega_0$ was arbitrary it follows that $\bigcup_{n=1}^\infty\{\omega\in\Omega:X(\omega)\in A_n\}\subseteq\{\omega\in\Omega:X(\omega)\in\bigcup_{n=1}^\infty A_n\}$.
Combining the two inclusions we conclude the proof.
%
% !TEX root = ../main.tex
\subsection*{Why is it sufficient to specify the law on a $\pi$-system}%
\manualtoc{subsection}{\texorpdfstring{Why is it sufficient to specify the law on a $\pi$-system}{Why is it sufficient to specify the law on a pi-system}}{studiabile-nav-0204}%

Since $\mathcal L_X$ is a probability measure on $\mathcal B(\mathbb R)$ we may ask whether it is necessary to specify $\mathcal L_X(A)$ $\forall$ Borel set $A$ in order to uniquely determine the law of $X$.
The answer is no. By Lemma 1.2 it is sufficient to specify the values of $\mathcal L_X$ on a $\pi$-system generating $\mathcal B(\mathbb R)$.
Recall that by Prop 1.1 $\pi(\mathbb R)=\{(-\infty,c]:c\in\mathbb R\}$ is a $\pi$-system and $\sigma(\pi(\mathbb R))=\mathcal B(\mathbb R)$.
Moreover by Lemma 1.2 if two finite measures on a $\sigma$-algebra generated by a $\pi$-system have the same total mass and coincide on that $\pi$-system, then they coincide on the entire $\sigma$-algebra.
In particular let $\mu_1$ and $\mu_2$ be two probability measures on $\mathcal B(\mathbb R)$ such that $\mu_1((-\infty,c])=\mu_2((-\infty,c])$ $\forall c\in\mathbb R$.
Since both measures are probability measures $\mu_1(\mathbb R)=\mu_2(\mathbb R)=1<\infty$.
Therefore all the assumption of Lemma 1.2 are satisfied and we conclude that $\mu_1(A)=\mu_2(A)$ $\forall A\in\mathcal B(\mathbb R)$.
Hence a probability measure on $\mathcal B(\mathbb R)$ is uniquely determined by its values on the intervals $(-\infty,c]$ for every $c\in\mathbb R$.
In particular the law of a random variable can be uniquely determined without specifying its value on every Borel set.
It is sufficient to know $\mathcal L_X((-\infty,c])$ $\forall c\in\mathbb R$.
This motivates the definition of the cumulative distribution function.
%
% !TEX root = ../main.tex
\subsection*{Definition CDF}%
\manualtoc{subsection}{\texorpdfstring{Definition CDF}{Definition CDF}}{studiabile-nav-0205}%

Let $X:\Omega\to\mathbb R$ be a random variable.
The CDF of $X$ is the function $F_X:\mathbb R\to[0,1]$ defined by
$F_X(c):=\mathcal L_X((-\infty,c])=\mathbb P(X\le c)$ $\forall c\in\mathbb R$.
Indeed by definition of the law of $X$
\[
\begin{aligned}
F_X(c)&=\mathcal L_X((-\infty,c])=\mathbb P(X^{-1}((-\infty,c])) =\mathbb P(\{\omega\in\Omega:X(\omega)\le c\}),
\end{aligned}
\]
where the second equality follows from the definition of $\mathcal L_X$ and the third from the definition of pre-image.

Since the intervals $(-\infty,c]$ form a $\pi$-system generating $\mathcal B(\mathbb R)$ by Lemma 1.2 the CDF uniquely determines the law $\mathcal L_X$.
In particular $F_X(c)=F_Y(c)$ $\forall c\in\mathbb R\Longleftrightarrow\mathcal L_X=\mathcal L_Y$.
Thus two random variables have the same law exactly when they have the same CDF.
%
% !TEX root = ../main.tex
\subsection*{Properties of cdf}%
\manualtoc{subsection}{\texorpdfstring{Properties of cdf}{Properties of cdf}}{studiabile-nav-0206}%

The cdf enjoys the following properties:
\studylabel{(a)} Range and monotonicity $F_X:\mathbb R\to[0,1]$, $F_X(y)\le F_X(z)$ $\forall y\le z$.

\studylabel{(b)} Limit at infinity. If $\mathbb P(|X|<\infty)=1$ then $\lim_{z\to+\infty}F_X(z)=1$, $\lim_{z\to-\infty}F_X(z)=0$.

\studylabel{(c)} Right-continuity $\forall z\in\mathbb R$, $\lim_{\varepsilon\downarrow0}F_X(z+\varepsilon)=F_X(z)$.

\par\noindent\studyheading{Proof.}\par\nobreak
\studylabel{(a)} By definition of the CDF $F_X(z)=\mathbb P(X\le z)$.
Since $\mathbb P$ is a probability measure we have $0\le\mathbb P(X\le z)\le1$, $\forall z\in\mathbb R$.
Therefore $0\le F_X(z)\le1$ $\forall z\in\mathbb R$ hence $F_X:\mathbb R\to[0,1]$.
For monotonicity, let $y,z\in\mathbb R$ such that $y\le z$. Notice that $\{X\le y\}\subseteq\{X\le z\}$.
Indeed if $\omega\in\{X\le y\}$ then $X(\omega)\le y\le z$ therefore $\omega\in\{X\le z\}$.
Since $X$ is $\mathcal F$-measurable both events belong to $\mathcal F$ and by monotonicity of probability HW~12 we obtain $\mathbb P(X\le y)\le\mathbb P(X\le z)$.
Hence by definition of cdf $F_X(y)\le F_X(z)$.

\studylabel{(b)} For the positive tail,
$\forall n\in\mathbb N$ define $A_n:=\{\omega\in\Omega:X(\omega)\le n\}$.
Since $X$ is $\mathcal F$-measurable by definition 3.1 $A_n\in\mathcal F$ $\forall n\in\mathbb N$.
Notice that $A_n\subseteq A_{n+1}$ $\forall n\in\mathbb N$, because $X(\omega)\le n$ implies $X(\omega)\le n+1$.
Moreover $\bigcup_{n=1}^\infty A_n=\{\omega\in\Omega:X(\omega)<+\infty\}$.
Since $\mathbb P(|X|<\infty)=1$, we have $\mathbb P(\bigcup_{n=1}^\infty A_n)=1$.
Therefore by Lemma 1.5 (continuity of probability from below)
\[
\lim_{n\to+\infty}F_X(n)=\lim_{n\to+\infty}\mathbb P(A_n)
=\mathbb P\left(\bigcup_{n=1}^\infty A_n\right)=1,
\]
where the first equality follows from the definition of the CDF, the second from Lemma 1.5 and third from assumption.
Since $F_X$ is non-decreasing by prop (a) the convergence $F_X(n)\to1$ implies that $\lim_{z\to+\infty}F_X(z)=1$.
For the negative tail,
$\forall n\in\mathbb N$ define $B_n:=\{\omega\in\Omega:X(\omega)\le-n\}$.

Again since $X$ is measurable, $B_n\in\mathcal F$ $\forall n\in\mathbb N$.
Notice that $B_n\supseteq B_{n+1}$ $\forall n\in\mathbb N$ because $X(\omega)\le-(n+1)$ implies $X(\omega)\le-n$.
Moreover $\bigcap_{n=1}^\infty B_n=\{\omega\in\Omega:X(\omega)=-\infty\}$.
By assumption $\mathbb P(|X|<\infty)=1$, $\mathbb P(\bigcap_{n=1}^\infty B_n)=0$.
Since the events $B_n$ form a non-increasing sequence and $\mathbb P(B_1)\le1<\infty$, we may apply Lemma 1.6 (continuity from above).
Therefore
\[
\lim_{n\to+\infty}F_X(-n)=\lim_{n\to+\infty}\mathbb P(B_n)
=\mathbb P\left(\bigcap_{n=1}^\infty B_n\right)=0,
\]
where the first equality follows from the definition of CDF, the second from Lemma 1.6 and third from ass.
Since $F_X$ is non-decreasing by prop (a) the convergence $F_X(-n)\to0$ implies that $\lim_{z\to-\infty}F_X(z)=0$.

\studylabel{(c)}
Fix an arbitrary $z$. $\forall n\in\mathbb N$ define $A_n:=\{\omega\in\Omega:X(\omega)\le z+\frac1n\}$.
Since $X$ is $\mathcal F$-measurable, $A_n\in\mathcal F$ $\forall n\in\mathbb N$.
Notice that $A_n\supseteq A_{n+1}$ $\forall n\in\mathbb N$ because $z+\frac1{n+1}\le z+\frac1n$.
\emph{We now show that $\bigcap_{n=1}^\infty A_n=\{X\le z\}$ by showing both inclusions.}
\emph{First we show that $\bigcap_{n=1}^\infty A_n\subseteq\{X\le z\}$.} Let $\omega_0\in\bigcap_{n=1}^\infty A_n$, by definition of intersection $\omega_0\in A_n$ $\forall n\in\mathbb N$.
Therefore by definition of $A_n$, $X(\omega_0)\le z+\frac1n$ $\forall n\in\mathbb N$.
Since $\frac1n\to0$ we obtain $X(\omega_0)\le z$.
Hence $\omega_0\in\{X\le z\}$.
Since $\omega_0$ was arbitrary we conclude that $\bigcap_{n=1}^\infty A_n\subseteq\{X\le z\}$.
\emph{For the reverse inclusion, we show that $\{X\le z\}\subseteq\bigcap_{n=1}^\infty A_n$.} Let $\omega_0\in\{X\le z\}$.
By definition of this set $X(\omega_0)\le z$.
Therefore $X(\omega_0)\le z<z+\frac1n$ $\forall n\in\mathbb N$.
Hence $\omega_0\in A_n$ $\forall n\in\mathbb N$.
By definition of intersection $\omega_0\in\bigcap_{n=1}^\infty A_n$.
Since $\omega_0$ was arbitrary we obtain $\{X\le z\}\subseteq\bigcap_{n=1}^\infty A_n$.
Combining the two inclusionss we conclude that $\bigcap_{n=1}^\infty A_n=\{X\le z\}$.
Now since $(A_n)_{n\in\mathbb N}$ is a decreasing sequence of events and $\mathbb P(A_1)\le1<\infty$, by Lemma 1.6 we obtain
\[
\begin{aligned}
\lim_{n\to+\infty}F_X(z+\tfrac1n)
&=\lim_{n\to+\infty}\mathbb P(A_n)=\mathbb P\left(\bigcap_{n=1}^\infty A_n\right) =\mathbb P(X\le z)=F_X(z),
\end{aligned}
\]
where the first equality follows from the definition of the CDF, the second from Lemma 1.6, the third from the equality of sets proved above and the last the definition of $F_X$.
Finally, to obtain the limit for arbitrary $\varepsilon\downarrow0$, let $\delta>0$.
From the preceding convergence, there exists $N\in\mathbb N$ such that $F_X(z+\frac1N)<F_X(z)+\delta$.
For every $0<\varepsilon\le\frac1N$, by monotonicity of $F_X$,
$F_X(z)\le F_X(z+\varepsilon)\le F_X(z+\frac1N)<F_X(z)+\delta$.
Since $\delta>0$ was arbitrary we conclude that
\[
\lim_{\varepsilon\downarrow0}F_X(z+\varepsilon)=F_X(z).
\]
Since $z\in\mathbb R$ was arbitrary, $F_X$ is right-continuous.
%
% !TEX root = ../main.tex
\subsection*{HW 33 --- CDF tails with mass at infinity}%
\manualtoc{subsection}{\texorpdfstring{HW 33 --- CDF tails with mass at infinity}{HW 33 --- CDF tails with mass at infinity}}{studiabile-nav-0214}%

Let $X$ be an extended-real random variable. Prove if $\mathbb P(X=+\infty)=p$ and $\mathbb P(X=-\infty)=q$ then $\lim_{z\to+\infty}F_X(z)=1-p$ and $\lim_{z\to-\infty}F_X(z)=q$.
\par\noindent\studyheading{Proof.}\par\nobreak
Recall that by definition of CDF $F_X(z)=\mathbb P(X\le z)$, $z\in\mathbb R$.
For the positive tail,
$\forall n\in\mathbb N$ define $A_n:=\{X\le n\}$.
Notice that $A_n\subseteq A_{n+1}$ $\forall n\in\mathbb N$ since $X(\omega)\le n\Rightarrow X(\omega)\le n+1$.
Therefore $(A_n)_{n\in\mathbb N}$ is an increasing sequence of events.
Moreover $\bigcup_{n=1}^\infty A_n=\{X<\infty\}$.
Indeed, $X(\omega)\le n$ for at least one $n\in\mathbb N$ if and only if $X(\omega)<+\infty$.
Hence by Lemma 1.5 continuity of probability from below we obtain
\[
\begin{aligned}
\lim_{n\to+\infty}F_X(n)&=\lim_{n\to+\infty}\mathbb P(A_n)
=\mathbb P\left(\bigcup_{n=1}^\infty A_n\right) =\mathbb P(X<+\infty)=1-\mathbb P(X=+\infty)=1-p,
\end{aligned}
\]
where the first equality follows from the definition of the CDF, the second from Lemma 1.5, the third from the equality of sets above and the fourth from the complement rule for probability.
Since $F_X$ is non-decreasing by prop (a) of the CDF, and $F_X(n)\le F_X(z)\le1-p$ $\forall z\ge n$ we conclude that $\lim_{z\to+\infty}F_X(z)=1-p$.
For the negative tail,
$\forall n\in\mathbb N$ define $B_n:=\{X\le-n\}$.
Notice that $B_n\supseteq B_{n+1}$ $\forall n\in\mathbb N$ since $X(\omega)\le-(n+1)\Rightarrow X(\omega)\le-n$.
Therefore $(B_n)_{n\in\mathbb N}$ is a decreasing sequence of events.
Moreover $\bigcap_{n=1}^\infty B_n=\{X=-\infty\}$.
Indeed $X(\omega)\le-n$ $\forall n\in\mathbb N\Longleftrightarrow X(\omega)=-\infty$.
Hence, since $\mathbb P(B_1)\le1<\infty$ by Lemma 1.6 (continuity of probability from above) we obtain
\[
\lim_{n\to+\infty}F_X(-n)=\lim_{n\to+\infty}\mathbb P(B_n)
=\mathbb P\left(\bigcap_{n=1}^\infty B_n\right)=\mathbb P(X=-\infty)=q,
\]
where the first equality follows from the definition of the CDF, the second from Lemma 1.6, the third from the equality of sets above and the last from the assumption.
Since $F_X$ is non-decreasing by prop (a) and $q\le F_X(z)\le F_X(-n)$ $\forall z\le-n$ we conclude that $\lim_{z\to-\infty}F_X(z)=q$.
%
% !TEX root = ../main.tex
\subsection*{HW 34 --- Left limits and atoms}%
\manualtoc{subsection}{\texorpdfstring{HW 34 --- Left limits and atoms}{HW 34 --- Left limits and atoms}}{studiabile-nav-0216}%

Prove that $\forall z\in\mathbb R$, $\{X\le z-\varepsilon\}\uparrow\{X<z\}$ as $\varepsilon\downarrow0$.
Then conclude that $\lim_{\varepsilon\downarrow0}F_X(z-\varepsilon)=F_X(z)-\mathcal L_X(\{z\})$.

\par\noindent\studyheading{Proof.}\par\nobreak
Fix an arbitrary $z\in\mathbb R$. For every $n\in\mathbb N$ define $A_n:=\{\omega\in\Omega:X(\omega)\le z-\frac1n\}$.
Since $X$ is $\mathcal F$-measurable we have $(A_n)_{n\in\mathbb N}\subseteq\mathcal F$.
Notice that $A_n\subseteq A_{n+1}$ $\forall n\in\mathbb N$.
Indeed let $\omega_0\in A_n$.
By definition of $A_n$ $X(\omega_0)\le z-\frac1n$, since $z-\frac1n<z-\frac1{n+1}$ we obtain $X(\omega_0)\le z-\frac1n<z-\frac1{n+1}$.
Therefore $\omega_0\in A_{n+1}$.
Since $\omega_0$ was arbitrary the inclusion follows.
Hence $(A_n)_{n\in\mathbb N}$ is an increasing sequence of events, \emph{we now prove that $\bigcup_{n=1}^\infty A_n=\{X<z\}$.}
We establish the equality by proving both inclusionss.
First inclusion. Let $\omega_0\in\bigcup_{n=1}^\infty A_n$.
By definition of union, there exists $n_0\in\mathbb N$ such that $\omega_0\in A_{n_0}$.
By definition of $A_{n_0}$, $X(\omega_0)\le z-\frac1{n_0}<z$.
Therefore $X(\omega_0)<z$ which implies that $\omega_0\in\{X<z\}$.
Since $\omega_0$ was arbitrary we conclude that $\bigcup_{n=1}^\infty A_n\subseteq\{X<z\}$.
Reverse inclusion. Let $\omega_0\in\{X<z\}$.
By definition of this set $X(\omega_0)<z$.
Therefore $z-X(\omega_0)>0$.
Since $\frac1n\to0$ as $n\to+\infty$ there exists $n_0\in\mathbb N$ such that $\frac1{n_0}<z-X(\omega_0)$.
Rearranging this inequality we obtain $X(\omega_0)<z-\frac1{n_0}$ which implies $X(\omega_0)\le z-\frac1{n_0}$ hence $\omega_0\in A_{n_0}$ and consequently $\omega_0\in\bigcup_{n=1}^\infty A_n$.
Since $\omega_0$ was arbitrary we conclude that $\{X<z\}\subseteq\bigcup_{n=1}^\infty A_n$.
Combining the two inclusionss, we obtain $\bigcup_{n=1}^\infty A_n=\{X<z\}$.
Since $(A_n)_{n\in\mathbb N}$ is an increasing sequence of events we have proved that $\{X\le z-\varepsilon\}\uparrow\{X<z\}$ as $\varepsilon\downarrow0$.
Since $(A_n)_{n\in\mathbb N}$ is an increasing sequence of events, by Lemma 1.5 we obtain
\[
\lim_{n\to+\infty}F_X(z-\tfrac1n)=\lim_{n\to+\infty}\mathbb P(A_n)
=\mathbb P\left(\bigcup_{n=1}^\infty A_n\right)=\mathbb P(X<z)
\]
where the first equality follows from the definition of the CDF and of $A_n$, the second from Lemma 1.5 and the third from the equality of sets proved above.
Moreover since $F_X$ is non-decreasing for every $0<\varepsilon\le\frac1n$,
$F_X(z-\frac1n)\le F_X(z-\varepsilon)\le\mathbb P(X<z)$.
Since the left-hand side converges to $\mathbb P(X<z)$ as $n\to+\infty$ we conclude that $\lim_{\varepsilon\downarrow0}F_X(z-\varepsilon)=\mathbb P(X<z)$.
Finally we notice that $\{X\le z\}=\{X<z\}\cup\{X=z\}$

where the two events are disjoint on the right-hand side.
Therefore by the additivity of the probability measure $\mathbb P$
$\mathbb P(X\le z)=\mathbb P(X<z)+\mathbb P(X=z)$.
Rearranging we obtain $\mathbb P(X<z)=\mathbb P(X\le z)-\mathbb P(X=z)$.
Hence
\[
\begin{aligned}
\lim_{\varepsilon\downarrow0}F_X(z-\varepsilon)
&=\mathbb P(X<z)=\mathbb P(X\le z)-\mathbb P(X=z) =F_X(z)-\mathcal L_X(\{z\}),
\end{aligned}
\]
where the first equality follows from the continuity of probability from below established above, the second from the additivity of $\mathbb P$, the third from the definition of the CDF and of the law of $X$.
Since $z\in\mathbb R$ was arbitrary we conclude that $\lim_{\varepsilon\downarrow0}F_X(z-\varepsilon)=F_X(z)-\mathcal L_X(\{z\})$.
%
% !TEX root = ../main.tex
\subsection*{Definition --- Probability Density Function (pdf)}%
\manualtoc{subsection}{\texorpdfstring{Definition --- Probability Density Function (pdf)}{Definition --- Probability Density Function (pdf)}}{studiabile-nav-0218}%

Let $(\Omega,\mathcal F,\mathbb P)$ be a probability space and let $X:\Omega\to\mathbb R$ be a r.v.
Recall that cdf $F_X(z):=\mathbb P(X\le z)$, $z\in\mathbb R$.
A r.v. $X$ admits a probability density function (pdf) if there exists a non-negative measurable function $f_X:\mathbb R\to[0,+\infty)$ such that
\[
\mathbb P(X\in A)=\int_A f_X(z)\,dz\quad\forall A\in\mathcal B(\mathbb R).
\]
In particular taking $A=(-\infty,z]$ we obtain $F_X(z)=\int_{-\infty}^z f_X(t)\,dt$.
If the CDF is absolutely continuous then it is differentiable almost everywhere and its density is given by $f_X(z)=F_X'(z)$ for almost every $z\in\mathbb R$.
%
% !TEX root = ../main.tex
\subsection*{Properties of pdf}%
\manualtoc{subsection}{\texorpdfstring{Properties of pdf}{Properties of pdf}}{studiabile-nav-0219}%

\studylabel{(a)} Non-negativity $f_X(z)\ge0$ for almost every $z\in\mathbb R$.
\studylabel{(b)} Normalization $\int_{-\infty}^{+\infty}f_X(z)\,dz=1$.
\par\noindent\studyheading{Proof.}\par\nobreak
\studylabel{(a)} Recall that by prop (a) of the CDF, $F_X$ is non-decreasing.
Therefore $\forall z\in\mathbb R$ and $\forall h>0$, $F_X(z+h)\ge F_X(z)$.
Subtracting $F_X(z)$ and dividing by $h>0$ we obtain $\frac{F_X(z+h)-F_X(z)}h\ge0$.
At every point where $F_X$ is differentiable letting $h\downarrow0$ gives
$F_X'(z)=\lim_{h\downarrow0}\frac{F_X(z+h)-F_X(z)}h\ge0$.
The integral representation $F_X(z)=\int_{-\infty}^{z}f_X(t)\,dt$ from the definition of a density makes $F_X$ absolutely continuous, with $F_X'=f_X$ almost everywhere by the fundamental theorem for the Lebesgue integral. Consequently, $f_X(z)\ge0$ almost everywhere, consistently with the non-negativity already required in the definition.
\studylabel{(b)} Recall that by def of probability density function
$\mathbb P(X\in A)=\int_A f_X(z)\,dz$ $\forall A\in\mathcal B(\mathbb R)$.
In particular taking

$A=\mathbb R$ we obtain $\int_{-\infty}^{+\infty}f_X(z)\,dz=\mathbb P(X\in\mathbb R)=\mathbb P(\Omega)=1$
where the first equality follows from the definition of the pdf, the second from the fact that $X:\Omega\to\mathbb R$ and the third from the normalization of the probability measure $\mathbb P$.

%
% !TEX root = ../main.tex
\subsection*{Theorem -- Existence of a r.v. with a prescribed CDF}%
\manualtoc{subsection}{\texorpdfstring{Theorem -- Existence of a r.v. with a prescribed CDF}{Theorem -- Existence of a r.v. with a prescribed CDF}}{studiabile-nav-0225}%

Let $F:\mathbb R\to[0,1]$ be a function satisfying the following properties:
\begin{itemize}
\item[(a)] Range and monotonicity $F:\mathbb R\to[0,1]$, $F(y)\le F(z)$ $\forall y\le z$
\item[(b)] Limits of infinity $\lim_{z\to+\infty}F(z)=1$ $\lim_{z\to-\infty}F(z)=0$
\item[(c)] Right continuity $\lim_{\varepsilon\downarrow0}F(z+\varepsilon)=F(z)$ $\forall z\in\mathbb R$
\end{itemize}
Then there exists a probability space $(\Omega,\mathcal F,\mathbb P)$ and a r.v. $X:\Omega\to\mathbb R$
such that $F_X(c)=F(c)$ $\forall c\in\mathbb R$

\subsubsection*{Proof}%
%

First we construct the probability space. Let $\Omega=\mathbb R$ and $\mathcal F=\mathcal B(\mathbb R)$.
We now introduce a measure $\mathcal L$ by defining $\forall a,b\in\mathbb R$ with $a<b$
$\mathcal L((a,b]):=F(b)-F(a)$.
Notice that, since $F$ is non-decreasing by property (a) we have $F(a)\le F(b)$ whenever $a<b$.
Therefore $\mathcal L((a,b])=F(b)-F(a)\ge0$.
The monotonicity and right-continuity of $F$ allow this assignment on intervals to be extended
to a measure $\mathcal L$ on $\mathcal B(\mathbb R)$ by the Lebesgue--Stieltjes measure construction. This construction is taken as a known result here; it is the same existence step used in the handwritten proof.
\emph{We now verify that $\mathcal L$ is a probability measure.}
Since $\mathcal L$ is already a measure on $\mathcal B(\mathbb R)$ it remains to prove that $\mathcal L(\mathbb R)=1$.
Notice that $(-n,n]\subseteq(-(n+1),n+1]$ $\forall n\in\mathbb N$ and $\bigcup_{n=1}^{\infty}(-n,n]=\mathbb R$.
Therefore since $((-n,n])_{n\in\mathbb N}$ is an increasing sequence of Borel sets,
by continuity of measures from below we obtain
\[
\begin{aligned}
\mathcal L(\mathbb R)&=\mathcal L\left(\bigcup_{n=1}^{\infty}(-n,n]\right)
=\lim_{n\to\infty}\mathcal L((-n,n]) =\lim_{n\to\infty}(F(n)-F(-n))=1-0=1
\end{aligned}
\]
Hence $\mathcal L$ is a probability measure on $\mathcal B(\mathbb R)$.
We can therefore consider the probability space $(\Omega,\mathcal F,\mathbb P)=(\mathbb R,\mathcal B(\mathbb R),\mathcal L)$.

Second we construct the random variable. Define $X:\Omega\to\mathbb R$ by $X(\omega):=\omega$ $\forall\omega\in\Omega$.
Notice that $X$ is identity function on $\mathbb R$.
\emph{We now verify that $X$ is a r.v. namely that is $\mathcal F$-measurable.}
Let $A\in\mathcal B(\mathbb R)$.
By definition of pre-image
\[
X^{-1}(A)=\{\omega\in\Omega:X(\omega)\in A\}=\{\omega\in\mathbb R:\omega\in A\}=A
\]
where the second equality follows from $\Omega=\mathbb R$ and $X(\omega)=\omega$.
Since $A\in\mathcal B(\mathbb R)=\mathcal F$ we conclude that $X^{-1}(A)\in\mathcal F$.
Since $A\in\mathcal B(\mathbb R)$ was arbitrary by Definition 3.1, $X$ is $\mathcal F$-measurable and therefore a random variable.

\emph{Finally we prove that the CDF of $X$ coincides with $F$.}
Fix an arbitrary $c\in\mathbb R$.
First notice that $(-n,c]\uparrow(-\infty,c]$ as $n\to\infty$
where we consider sufficiently large $n$ so that $-n<c$.
Therefore by continuity of measure from below
\[
\mathcal L((-\infty,c])=\lim_{n\to\infty}\mathcal L((-n,c])
=\lim_{n\to\infty}(F(c)-F(-n))=F(c)
\]
where the first equality follows from continuity of measure from below,
the second from the definition of $\mathcal L$ and the third from property (b)
since $\lim_{n\to\infty}F(-n)=0$.
Now by definition of the CDF
\[
\begin{aligned}
F_X(c)&=\mathcal L_X((-\infty,c])=\mathbb P(X\le c)
=\mathbb P(\{\omega\in\Omega:X(\omega)\le c\})\\
&=\mathbb P(\{\omega\in\mathbb R:\omega\le c\})
=\mathbb P((-\infty,c])=\mathcal L((-\infty,c])=F(c)
\end{aligned}
\]
where the first and second equalities follow from the definition of CDF and the law of $X$,
the fourth from $X(\omega)=\omega$ and $\Omega=\mathbb R$,
the sixth from $\mathbb P=\mathcal L$ and the last from calculation above.
Since $c\in\mathbb R$ was arbitrary we conclude that $F_X(c)=F(c)$ $\forall c\in\mathbb R$.
Thus we have constructed a r.v. $X$ whose CDF is precisely $F$.
%
% !TEX root = ../main.tex
\subsection*{Construction --- Inverse Transform Method}%
\manualtoc{subsection}{\texorpdfstring{Construction --- Inverse Transform Method}{Construction --- Inverse Transform Method}}{studiabile-nav-0227}%


Another way to construct a random variable with a prescribed CDF $F$ is to start from a uniformly distributed random variable $U\sim\operatorname{Unif}(0,1)$.

Now, define
\[
X:=F^{-1}(U),
\]
where $F^{-1}$ denotes the inverse function when it exists, or the generalized inverse for a general CDF.

The idea is that applying $F^{-1}$ to a uniform random variable transforms its distribution into the prescribed distribution $F$. 

Thus, the inverse transform method allows us to construct a random variable with any prescribed CDF $F$.
%
% !TEX root = ../main.tex
\subsection*{Lemma 3.7 --- The $\sigma$-Algebra Generated by a Random Variable}%
\manualtoc{subsection}{\texorpdfstring{Lemma 3.7 --- The $\sigma$-Algebra Generated by a Random Variable}{Lemma 3.7 --- The sigma-Algebra Generated by a Random Variable}}{studiabile-nav-0228}%


\subsubsection*{Statement}%
%


Let $Y:\Omega\to\mathbb R$ be a random variable, and consider the class
\[
\begin{aligned}
\pi(Y)&:=\bigl\{\{\omega\in\Omega:Y(\omega)\le x\}:x\in\mathbb R\bigr\} =Y^{-1}(\pi(\mathbb R)),
\end{aligned}
\]
where
\[
\pi(\mathbb R):=\{(-\infty,x]:x\in\mathbb R\}.
\]
Then
\[
\sigma(Y)=\sigma(\pi(Y)).
\]

\subsubsection*{Expanded Proof of Lemma 3.7}%
%


As we will see in HW 35, $\pi(Y)$ is a $\pi$-system on $\Omega$.

We prove the two inclusions.

\emph{First, we prove that
\[
\sigma(Y)\subseteq\sigma(\pi(Y)).
\]}
By definition of the pre-image of a set, for every $x\in\mathbb R$,
\[
\begin{aligned}
Y^{-1}((-\infty,x])
&=\{\omega\in\Omega:Y(\omega)\in(-\infty,x]\} =\{\omega\in\Omega:Y(\omega)\le x\} =\{Y\le x\},
\end{aligned}
\]
where the second equality follows from the fact that $Y(\omega)\in(-\infty,x]$ if and only if $Y(\omega)\le x$.

Since $\{Y\le x\}$ is precisely one of the sets defining $\pi(Y)$, we obtain
\[
Y^{-1}((-\infty,x])=\{Y\le x\}\in\pi(Y).
\]
Since $\sigma(\pi(Y))$ is the smallest $\sigma$-algebra containing $\pi(Y)$, we have
\[
\pi(Y)\subseteq\sigma(\pi(Y)).
\]
Consequently,
\[
Y^{-1}((-\infty,x])\in\sigma(\pi(Y))
\qquad\forall x\in\mathbb R.
\]
In other words, the pre-image under $Y$ of every set belonging to $\pi(\mathbb R)$ belongs to $\sigma(\pi(Y))$.

Moreover, recall that, by Proposition 1.1,
\[
\sigma(\pi(\mathbb R))=\mathcal B(\mathbb R).
\]
Thus, all the assumptions of Lemma 3.2 are satisfied.

Therefore, by Lemma 3.2, $Y$ is measurable with respect to $\sigma(\pi(Y))$, namely
\[
Y\in m\sigma(\pi(Y)).
\]
Now, recall that $\sigma(Y)$ is the smallest $\sigma$-algebra on $\Omega$ with respect to which $Y$ is measurable.

Since $\sigma(\pi(Y))$ is a $\sigma$-algebra with respect to which $Y$ is measurable, by the minimality of $\sigma(Y)$, we conclude that
\[
\sigma(Y)\subseteq\sigma(\pi(Y)).
\]
\emph{We now prove the reverse inclusion, namely
\[
\sigma(\pi(Y))\subseteq\sigma(Y).
\]}
Recall that
\[
\sigma(Y)=Y^{-1}(\mathcal B(\mathbb R)).
\]
\emph{We first prove that
\[
\pi(Y)\subseteq\sigma(Y).
\]}
Let $A\in\pi(Y)$.

By definition of $\pi(Y)$, there exists $x_0\in\mathbb R$ such that
\[
A=\{Y\le x_0\}.
\]
As shown above, by definition of pre-image,
\[
A=\{Y\le x_0\}=Y^{-1}((-\infty,x_0]).
\]
Notice that $(-\infty,x_0]$ is a Borel set, since it is closed in $\mathbb R$. Therefore,
\[
(-\infty,x_0]\in\mathcal B(\mathbb R).
\]
Hence, by definition of $\sigma(Y)$,
\[
A=Y^{-1}((-\infty,x_0])\in\sigma(Y).
\]
Since $A\in\pi(Y)$ was arbitrary, we conclude that
\[
\pi(Y)\subseteq\sigma(Y).
\]
Now, since $\sigma(Y)$ is a $\sigma$-algebra containing $\pi(Y)$, and $\sigma(\pi(Y))$ is the smallest $\sigma$-algebra containing $\pi(Y)$, we obtain
\[
\sigma(\pi(Y))\subseteq\sigma(Y).
\]
Combining the two inclusions, we conclude that
\[
\sigma(Y)=\sigma(\pi(Y)).
\]
%
% !TEX root = ../main.tex
\subsection*{HW 35 --- Prove That $\pi(Y)$ Is a $\pi$-System on $\Omega$}%
\manualtoc{subsection}{\texorpdfstring{HW 35 --- Prove That $\pi(Y)$ Is a $\pi$-System on $\Omega$}{HW 35 --- Prove That pi(Y) Is a pi-System on Omega}}{studiabile-nav-0231}%


\subsubsection*{Statement}%
%


Let $Y:\Omega\to\mathbb R$ be a random variable, and consider
\[
\pi(Y):=\{\{Y\le x\}:x\in\mathbb R\}.
\]
Prove that $\pi(Y)$ is a $\pi$-system on $\Omega$.

\subsubsection*{Proof}%
%


Recall that a $\pi$-system is a non-empty collection of sets closed under finite intersections.

It is trivial that $\pi(Y)$ is non-empty.

\emph{We now prove that $\pi(Y)$ is closed under finite intersections.}

Let $A,B\in\pi(Y)$. By definition, there exist $x,y\in\mathbb R$ such that
\[
A=\{Y\le x\},\qquad B=\{Y\le y\}.
\]
Recall that, as shown in HW 6,
\[
(-\infty,x]\cap(-\infty,y]
=(-\infty,\min\{x,y\}].
\]
Therefore,
\[
\begin{aligned}
A\cap B
&=\{Y\le x\}\cap\{Y\le y\} =Y^{-1}((-\infty,x])\cap Y^{-1}((-\infty,y])\\&=Y^{-1}((-\infty,x]\cap(-\infty,y]) =Y^{-1}((-\infty,\min\{x,y\}]) =\{Y\le\min\{x,y\}\},
\end{aligned}
\]
where the third equality follows from HW 28-\studylabel{(iii)}, and the fourth from HW 6.

Since $\min\{x,y\}\in\mathbb R$, by definition of $\pi(Y)$,
\[
\{Y\le\min\{x,y\}\}\in\pi(Y).
\]
Hence,
\[
A\cap B\in\pi(Y).
\]
Since $A,B\in\pi(Y)$ were arbitrary, $\pi(Y)$ is closed under finite intersections.

Therefore,
\[
\pi(Y)\text{ is a }\pi\text{-system on }\Omega.
\]
The statement is proved.
%
% !TEX root = ../main.tex
\subsection*{Properties --- Measurable Functions of Random Variables}%
\manualtoc{subsection}{\texorpdfstring{Properties --- Measurable Functions of Random Variables}{Properties --- Measurable Functions of Random Variables}}{studiabile-nav-0234}%


Here are some useful facts that we state without proof.

\noindent\studyheading{(a)} Let $Y:\Omega\to\mathbb R$ and $Z:\Omega\to\mathbb R$. We have
\[
Z\in m\sigma(Y)
\iff
\exists f\in m\mathcal B(\mathbb R)
\text{ such that }Z=f(Y),
\]
where $f:\mathbb R\to\mathbb R$.

\noindent\studyheading{(b)} Let $Y:\Omega\to\mathbb R^n$, i.e.,
\[
Y(\omega)=
(Y_1(\omega),Y_2(\omega),\ldots,Y_n(\omega)),
\]
and let $Z:\Omega\to\mathbb R$. We have
\[
Z\in m\sigma(Y)
\iff
\exists f\in m\mathcal B(\mathbb R^n)
\text{ such that }Z=f(Y),
\]
where $f:\mathbb R^n\to\mathbb R$.

\noindent\studyheading{(c)} Let $(Y_\gamma)_{\gamma\in I}$ be a collection of functions $Y_\gamma:\Omega\to\mathbb R$, and let $Z:\Omega\to\mathbb R$. We have
\[
\begin{aligned}
Z\in m\sigma(Y_\gamma,\gamma\in I)
\iff {}&
\text{there exist a countable collection}\\
&(\gamma_n)_{n\in\mathbb N}
\text{ and a function}\\
&f\in m\mathcal B(\mathbb R^{\mathbb N})\\
&\text{such that }
Z=f((Y_{\gamma_n})_{n\in\mathbb N}),
\end{aligned}
\]
where $f:\mathbb R^{\mathbb N}\to\mathbb R$.

\subsubsection*{Financial Interpretation}%
\manualtoc{subsubsection}{\texorpdfstring{Financial Interpretation}{Financial Interpretation}}{studiabile-nav-0238}%


Item (b) can be interpreted as follows.

Let $(Y_1,\ldots,Y_n)$ be $n$ assets traded on the market, and let $Z$ be the payoff of a derivative contract written on the vector $(Y_1,\ldots,Y_n)$.

Then there must exist a function $f$ such that
\[
Z=f(Y_1,\ldots,Y_n),
\]
which describes the random payoff $Z$ of the derivative contract.
%
\clearpage
%
\section*{4. Independence}
\manualtoc{section}{\texorpdfstring{4. Independence}{4. Independence}}{studiabile-nav-0239}
\subsection*{Definition 4.1 --- Independent Sigma-Algebras}%
\manualtoc{subsection}{\texorpdfstring{Definition 4.1 --- Independent Sigma-Algebras}{Definition 4.1 --- Independent Sigma-Algebras}}{studiabile-nav-0240}%

Sub-sigma-algebras $\mathcal G$ and $\mathcal H$ of $\mathcal F$ are independent when
\[
\mathbb P(G\cap H)=\mathbb P(G)\mathbb P(H)
\qquad\forall G\in\mathcal G,\ H\in\mathcal H.
\]
For a family of sub-sigma-algebras, the corresponding factorisation is required for every finite selection of distinct members and one event from each selected member.

\subsection*{Definition 4.2 --- Independent Random Variables}%
\manualtoc{subsection}{\texorpdfstring{Definition 4.2 --- Independent Random Variables}{Definition 4.2 --- Independent Random Variables}}{studiabile-nav-0241}%

Random variables are independent when their generated sigma-algebras are independent. In particular, the threshold criteria below test this property using the classes $\pi(X)$ introduced in Lemma~3.7.

\subsection*{Definition 4.3 --- Independent Events}%
\manualtoc{subsection}{\texorpdfstring{Definition 4.3 --- Independent Events}{Definition 4.3 --- Independent Events}}{studiabile-nav-0242}%

Events are independent when the sigma-algebras generated by their indicator functions are independent. HW~36 describes each of these sigma-algebras explicitly.
%
% !TEX root = ../main.tex
\subsection*{HW 36 --- $\sigma$-Algebra Generated by an Indicator Function}%
\manualtoc{subsection}{\texorpdfstring{HW 36 --- $\sigma$-Algebra Generated by an Indicator Function}{HW 36 --- sigma-Algebra Generated by an Indicator Function}}{studiabile-nav-0243}%


\subsubsection*{Statement}%
%


Let $E\in\mathcal F$. Prove that
\[
\sigma(\mathbf 1_E)=\{\emptyset,\Omega,E,E^c\}.
\]

\subsubsection*{Proof}%
%


Recall that, by Lemma 3.7,
\[
\sigma(\mathbf 1_E)=\sigma(\pi(\mathbf 1_E)),
\]
where
\[
\pi(\mathbf 1_E)=\{\{\mathbf 1_E\le x\}:x\in\mathbb R\}.
\]

Since the indicator function takes only the values $0$ and $1$, we have
\[
\{\mathbf 1_E\le x\}
=
\begin{cases}
\emptyset,&x<0,\\
E^c,&0\le x<1,\\
\Omega,&x\ge1.
\end{cases}
\]

Therefore,
\[
\pi(\mathbf 1_E)=\{\emptyset,E^c,\Omega\}.
\]

Since every $\sigma$-algebra is closed under complements, we must also include
\[
E=(E^c)^c.
\]

Notice that the collection
\[
\mathcal G:=\{\emptyset,\Omega,E,E^c\}
\]
is already a $\sigma$-algebra on $\Omega$, since it contains $\Omega$, is closed under complements and countable unions, and consequently also under intersections.

Therefore, no additional sets are needed, and we obtain
\[
\sigma(\pi(\mathbf 1_E))
=\{\emptyset,\Omega,E,E^c\}.
\]

Finally, by Lemma 3.7,
\[
\sigma(\mathbf 1_E)
=\sigma(\pi(\mathbf 1_E))
=\{\emptyset,\Omega,E,E^c\}.
\]

Therefore, the statement is proved. 

\subsubsection*{Connection with Independence of Events}%
\manualtoc{subsubsection}{\texorpdfstring{Connection with Independence of Events}{Connection with Independence of Events}}{studiabile-nav-0246}%


This tells us how to calculate all elements in $\sigma(\mathbf 1_E)$. In particular, given two independent events $E,F\in\mathcal F$, according to the definition above we have
\[
\begin{aligned}
\mathbb P(E\cap F)
&=\mathbb P(E)\mathbb P(F),\\[4pt]
\mathbb P(E^c\cap F)
&=\mathbb P(E^c)\mathbb P(F),\\[4pt]
\mathbb P(E^c\cap F^c)
&=\mathbb P(E^c)\mathbb P(F^c),\\[4pt]
\mathbb P(E\cap F^c)
&=\mathbb P(E)\mathbb P(F^c).
\end{aligned}
\]

Plus analogous relations involving $\emptyset$ and $\Omega$. So here we recognise the ``old'' notion of independence!

Next comes a lemma that will further allow us to deepen our understanding of independence for random variables.
%
% !TEX root = ../main.tex
\subsection*{Lemma 4.1 --- Independence of $\sigma$-Algebras Generated by $\pi$-Systems}%
\manualtoc{subsection}{\texorpdfstring{Lemma 4.1 --- Independence of $\sigma$-Algebras Generated by $\pi$-Systems}{Lemma 4.1 --- Independence of sigma-Algebras Generated by pi-Systems}}{studiabile-nav-0247}%


\subsubsection*{Statement}%
%


Let $(\Omega,\mathcal F,\mathbb P)$ be a probability space. Let $\mathcal G$ and $\mathcal H$ be sub-$\sigma$-algebras of $\mathcal F$, and let $\mathcal I$ and $\mathcal J$ be $\pi$-systems such that
\[
\sigma(\mathcal I)=\mathcal G,
\qquad
\sigma(\mathcal J)=\mathcal H.
\]

Then the $\sigma$-algebras $\mathcal G$ and $\mathcal H$ are independent if and only if
\[
\mathbb P(I\cap J)=\mathbb P(I)\mathbb P(J)
\qquad
\forall I\in\mathcal I,\ J\in\mathcal J.
\]

\subsubsection*{Expanded Proof of Lemma 4.1}%
%


\studyheading{$(\Rightarrow)$} Assume that $\mathcal G$ and $\mathcal H$ are independent.

By Definition 4.1,
\[
\mathbb P(G\cap H)=\mathbb P(G)\mathbb P(H)
\qquad
\forall G\in\mathcal G,\ H\in\mathcal H.
\]

Since
\[
\mathcal I\subseteq\sigma(\mathcal I)=\mathcal G,
\qquad
\mathcal J\subseteq\sigma(\mathcal J)=\mathcal H,
\]

the equality holds, in particular, for every $I\in\mathcal I$ and $J\in\mathcal J$. Therefore,
\[
\mathbb P(I\cap J)=\mathbb P(I)\mathbb P(J)
\qquad\forall I\in\mathcal I,\ J\in\mathcal J.
\]

\studyheading{$(\Leftarrow)$} Assume the factorisation in the statement.

First, fix an arbitrary $I\in\mathcal I$, and consider the two set functions defined on $(\Omega,\mathcal H)$ by
\[
\mu_1(H):=\mathbb P(I\cap H),
\qquad
\mu_2(H):=\mathbb P(I)\mathbb P(H),
\quad H\in\mathcal H.
\]

As we will see in HW 37, these two set functions are measures on $(\Omega,\mathcal H)$.

Notice that they have the same finite total mass. Indeed,
\[
\begin{aligned}
\mu_1(\Omega)
=\mathbb P(I\cap\Omega) =\mathbb P(I) =\mathbb P(I)\mathbb P(\Omega) =\mu_2(\Omega)<\infty,
\end{aligned}
\]
where the second equality follows from $I\cap\Omega=I$, and the third from $\mathbb P(\Omega)=1$.

Moreover, we prove that the two measures coincide on the $\pi$-system $\mathcal J$.

Let $J\in\mathcal J$. By our initial assumption,
\[
\begin{aligned}
\mu_1(J)
=\mathbb P(I\cap J) =\mathbb P(I)\mathbb P(J) =\mu_2(J),
\end{aligned}
\]
where the second equality follows precisely from the assumption of this implication.

Since $J\in\mathcal J$ was arbitrary, we obtain
\[
\mu_1(J)=\mu_2(J)
\qquad\forall J\in\mathcal J.
\]

Thus, the two measures have the same finite total mass and coincide on the $\pi$-system $\mathcal J$.

Since $\sigma(\mathcal J)=\mathcal H$, all the assumptions of Lemma 1.2 are satisfied. Therefore, by Lemma 1.2, the two measures coincide on the entire $\sigma$-algebra $\mathcal H$:
\[
\mu_1(H)=\mu_2(H)
\qquad\forall H\in\mathcal H.
\]

Substituting their definitions, and recalling that $I\in\mathcal I$ was arbitrary, we conclude that
\[
\mathbb P(I\cap H)=\mathbb P(I)\mathbb P(H)
\qquad
\forall I\in\mathcal I,\ H\in\mathcal H.
\]

We now fix an arbitrary $H\in\mathcal H$ and consider the two measures on $(\Omega,\mathcal G)$ defined by
\[
\nu_1(G):=\mathbb P(G\cap H),
\qquad
\nu_2(G):=\mathbb P(G)\mathbb P(H),
\quad G\in\mathcal G.
\]

By the same argument as above, these two measures have the same finite total mass:
\[
\nu_1(\Omega)=\mathbb P(H)=\nu_2(\Omega)<\infty.
\]

Moreover, by the equality established in the first part of the proof, they coincide on the $\pi$-system $\mathcal I$:
\[
\nu_1(I)=\nu_2(I)
\qquad\forall I\in\mathcal I.
\]

Therefore, since $\sigma(\mathcal I)=\mathcal G$, by Lemma 1.2 the two measures also coincide on the entire $\sigma$-algebra $\mathcal G$:
\[
\nu_1(G)=\nu_2(G)
\qquad\forall G\in\mathcal G.
\]

Substituting their definitions, and recalling that $H\in\mathcal H$ was arbitrary, we conclude that
\[
\mathbb P(G\cap H)=\mathbb P(G)\mathbb P(H)
\qquad
\forall G\in\mathcal G,\ H\in\mathcal H.
\]

Thus, by Definition 4.1, the $\sigma$-algebras $\mathcal G$ and $\mathcal H$ are independent.

%
% !TEX root = ../main.tex
\subsection*{HW 37 --- Prove That Two Set Functions Define Measures}%
\manualtoc{subsection}{\texorpdfstring{HW 37 --- Prove That Two Set Functions Define Measures}{HW 37 --- Prove That Two Set Functions Define Measures}}{studiabile-nav-0252}%


\subsubsection*{Statement}%
%


Let $(\Omega,\mathcal F,\mathbb P)$ be a probability space, let $\mathcal H\subseteq\mathcal F$ be a $\sigma$-algebra, and let $I\in\mathcal F$.

Prove that the two set functions
\[
H\longmapsto\mathbb P(I\cap H),
\qquad
H\longmapsto\mathbb P(I)\mathbb P(H)
\]
define two measures on $(\Omega,\mathcal H)$.

\subsubsection*{Expanded Proof of HW 37}%
%


Write $\mu_1$ and $\mu_2$ for the first and second set functions in the statement, respectively.

\emph{First, we prove that $\mu_1$ is a measure.}

Let $H\in\mathcal H$. Since $\mathcal H\subseteq\mathcal F$, we have $H\in\mathcal F$. Moreover, since $I\in\mathcal F$ and $\mathcal F$ is closed under intersections,
\[
I\cap H\in\mathcal F.
\]

Therefore, $\mu_1(H)=\mathbb P(I\cap H)$ is well-defined and non-negative:
\[
\mu_1(H)\ge0
\qquad\forall H\in\mathcal H.
\]

Furthermore,
\[
\begin{aligned}
\mu_1(\emptyset)
=\mathbb P(I\cap\emptyset) =\mathbb P(\emptyset) =0.
\end{aligned}
\]

\emph{We now prove countable additivity.}

Let $(H_n)_{n\in\mathbb N}\subseteq\mathcal H$ be a sequence of pairwise disjoint sets, namely
\[
H_i\cap H_j=\emptyset
\qquad\forall i\neq j.
\]

Notice that the sets $I\cap H_n$ are also pairwise disjoint. Indeed, for every $i\neq j$,
\[
\begin{aligned}
(I\cap H_i)\cap(I\cap H_j)
&=I\cap(H_i\cap H_j) =I\cap\emptyset =\emptyset.
\end{aligned}
\]

Moreover, since $I\in\mathcal F$ and $H_n\in\mathcal H\subseteq\mathcal F$, we have $I\cap H_n\in\mathcal F$ for every $n\in\mathbb N$.

Therefore,
\[
\begin{aligned}
\mu_1\left(\bigcup_{n=1}^{\infty}H_n\right)
&=\mathbb P\left(I\cap\bigcup_{n=1}^{\infty}H_n\right) =\mathbb P\left(\bigcup_{n=1}^{\infty}(I\cap H_n)\right) =\sum_{n=1}^{\infty}\mathbb P(I\cap H_n) =\sum_{n=1}^{\infty}\mu_1(H_n),
\end{aligned}
\]
where the second equality follows from the distributive property of intersections over unions, and the third from the countable additivity of $\mathbb P$, since the sets $I\cap H_n$ are measurable and pairwise disjoint.

Thus, $\mu_1$ is a measure on $(\Omega,\mathcal H)$.

\emph{We now prove that $\mu_2$ is also a measure.}

Since $\mathbb P(I)\ge0$ and $\mathbb P$ is a probability measure, we immediately have
\[
\mu_2(H)\ge0
\qquad\forall H\in\mathcal H,
\]
and
\[
\mu_2(\emptyset)
=\mathbb P(I)\mathbb P(\emptyset)
=0.
\]

Moreover, let $(H_n)_{n\in\mathbb N}\subseteq\mathcal H$ be a sequence of pairwise disjoint sets. By the countable additivity of $\mathbb P$,
\[
\begin{aligned}
\mu_2\left(\bigcup_{n=1}^{\infty}H_n\right)
&=\mathbb P(I)\mathbb P\left(\bigcup_{n=1}^{\infty}H_n\right) =\mathbb P(I)\sum_{n=1}^{\infty}\mathbb P(H_n) =\sum_{n=1}^{\infty}\mathbb P(I)\mathbb P(H_n) =\sum_{n=1}^{\infty}\mu_2(H_n).
\end{aligned}
\]

Therefore, $\mu_2$ is also a measure on $(\Omega,\mathcal H)$.

We conclude that both set functions
\[
\begin{aligned}
\mu_1(H)&=\mathbb P(I\cap H),\\
\mu_2(H)&=\mathbb P(I)\mathbb P(H),
\end{aligned}
\qquad H\in\mathcal H
\]
define measures on $(\Omega,\mathcal H)$.

%
% !TEX root = ../main.tex
\subsection*{Corollary --- Characterisation of Independence via Threshold Events}%
\manualtoc{subsection}{\texorpdfstring{Corollary --- Characterisation of Independence via Threshold Events}{Corollary --- Characterisation of Independence via Threshold Events}}{studiabile-nav-0255}%


Let $X,Y:\Omega\to\mathbb R$ be two random variables.

Then $X$ and $Y$ are independent if and only if
\[
\mathbb P(X\le x,\ Y\le y)
=
\mathbb P(X\le x)\,\mathbb P(Y\le y)
\qquad
\forall x,y\in\mathbb R.
\]

Indeed, this is a direct application of Lemma 4.1 with
\[
\mathcal I:=\pi(X)=\bigl\{\{\omega\in\Omega:X(\omega)\le x\}:x\in\mathbb R\bigr\},
\]

and
\[
\mathcal J:=\pi(Y)=\bigl\{\{\omega\in\Omega:Y(\omega)\le y\}:y\in\mathbb R\bigr\}.
\]

Recall that, by Lemma 3.7,
\[
\sigma(\pi(X))=\sigma(X),\qquad \sigma(\pi(Y))=\sigma(Y).
\]

Therefore, the independence of the random variables $X$ and $Y$ is equivalent to the factorisation of the probabilities of the events $\{X\le x\}$ and $\{Y\le y\}$ for every $x,y\in\mathbb R$.

\subsection*{Generalisation --- Independence of Finitely Many Random Variables}%
\manualtoc{subsection}{\texorpdfstring{Generalisation --- Independence of Finitely Many Random Variables}{Generalisation --- Independence of Finitely Many Random Variables}}{studiabile-nav-0256}%


More generally, random variables $X_1,\ldots,X_n$ are independent if and only if, for every $x_1,\ldots,x_n\in\mathbb R$,
\[
\mathbb P(X_1\le x_1,\ldots,X_n\le x_n)
=
\prod_{i=1}^{n}\mathbb P(X_i\le x_i).
\]

This is the analogous finite-dimensional characterisation of independence, obtained by applying the same idea to the $\pi$-systems generated by the threshold events of the variables $X_1,\ldots,X_n$.
%
% !TEX root = ../main.tex
\subsection*{HW 38 --- Independence of Three $\sigma$-Algebras Generated by $\pi$-Systems}%
\manualtoc{subsection}{\texorpdfstring{HW 38 --- Independence of Three $\sigma$-Algebras Generated by $\pi$-Systems}{HW 38 --- Independence of Three sigma-Algebras Generated by pi-Systems}}{studiabile-nav-0257}%


% BEGIN PFF RESTORATION HW38_SOURCE_ATTRIBUTION
(Williams, E4.1, p.~225.)
% END PFF RESTORATION HW38_SOURCE_ATTRIBUTION

\subsubsection*{Statement}%
%


Let $(\Omega,\mathcal F,\mathbb P)$ be a probability space, and let $\mathcal I_1,\mathcal I_2,\mathcal I_3$ be three $\pi$-systems such that
\[
\mathcal I_k\subseteq\mathcal F,
\qquad
\Omega\in\mathcal I_k,
\qquad k=1,2,3.
\]

Prove that $\sigma(\mathcal I_1)$, $\sigma(\mathcal I_2)$ and $\sigma(\mathcal I_3)$ are independent if and only if
\[
\mathbb P(I_1\cap I_2\cap I_3)
=
\prod_{k=1}^{3}\mathbb P(I_k)
\]
for every $I_k\in\mathcal I_k$, $k=1,2,3$.

Why is it necessary to assume that $\Omega\in\mathcal I_k$, $k=1,2,3$?

\subsubsection*{Part I --- Expanded Proof of the Equivalence}%
%


Let
\[
\mathcal G_k:=\sigma(\mathcal I_k),
\qquad k=1,2,3.
\]

\studyheading{($\Rightarrow$)} Assume that $\mathcal G_1,\mathcal G_2,\mathcal G_3$ are independent.

By Definition 4.1, for every $G_k\in\mathcal G_k$, $k=1,2,3$,
\[
\mathbb P(G_1\cap G_2\cap G_3)
=
\prod_{k=1}^{3}\mathbb P(G_k).
\]

Since
\[
\mathcal I_k\subseteq\sigma(\mathcal I_k)=\mathcal G_k,
\qquad k=1,2,3,
\]

the equality holds, in particular, for every $I_k\in\mathcal I_k$. Therefore,
\[
\mathbb P(I_1\cap I_2\cap I_3)
=
\prod_{k=1}^{3}\mathbb P(I_k).
\]

\studyheading{($\Leftarrow$)} Assume the factorisation in the statement.

First, without loss of generality, we consider $\mathcal G_1$ and $\mathcal G_2$, since the three indices play symmetric roles. In particular, since $\Omega\in\mathcal I_3$, we can choose $I_3=\Omega$ in our assumption.

Therefore, for every $I_1\in\mathcal I_1$ and $I_2\in\mathcal I_2$,
\[
\begin{aligned}
\mathbb P(I_1\cap I_2)
&=\mathbb P(I_1\cap I_2\cap\Omega) =\mathbb P(I_1)\mathbb P(I_2)\mathbb P(\Omega) =\mathbb P(I_1)\mathbb P(I_2),
\end{aligned}
\]
where the second equality follows from our assumption, and the third from $\mathbb P(\Omega)=1$.

Since $\mathcal I_1$ and $\mathcal I_2$ are $\pi$-systems generating $\mathcal G_1$ and $\mathcal G_2$, respectively, by Lemma 4.1 we conclude that
\[
\mathcal G_1 \text{ and } \mathcal G_2
\text{ are independent}.
\]

We now introduce the collection
\[
\mathcal I_{1,2}
:=
\{I_1\cap I_2:
I_1\in\mathcal I_1,\ I_2\in\mathcal I_2\}.
\]

\emph{We first prove that $\mathcal I_{1,2}$ is a $\pi$-system.}

Let $A,B\in\mathcal I_{1,2}$. By definition, there exist $I_1,J_1\in\mathcal I_1$ and $I_2,J_2\in\mathcal I_2$ such that
\[
A=I_1\cap I_2,
\qquad
B=J_1\cap J_2.
\]

Therefore,
\[
\begin{aligned}
A\cap B
&=(I_1\cap I_2)\cap(J_1\cap J_2) =(I_1\cap J_1)\cap(I_2\cap J_2),
\end{aligned}
\]
where the second equality follows from the associative and commutative properties of intersections.

Since $\mathcal I_1$ and $\mathcal I_2$ are $\pi$-systems,
\[
I_1\cap J_1\in\mathcal I_1,
\qquad
I_2\cap J_2\in\mathcal I_2.
\]

Hence,
\[
A\cap B\in\mathcal I_{1,2}.
\]

Moreover, $\Omega=\Omega\cap\Omega\in\mathcal I_{1,2}$, since $\Omega\in\mathcal I_1$ and $\Omega\in\mathcal I_2$.

Thus, $\mathcal I_{1,2}$ is a non-empty collection closed under intersections, and therefore a $\pi$-system.

Notice also that, since $\Omega\in\mathcal I_1$ and $\Omega\in\mathcal I_2$,
\[
I_1=I_1\cap\Omega\in\mathcal I_{1,2}
\qquad\forall I_1\in\mathcal I_1,
\]

and
\[
I_2=\Omega\cap I_2\in\mathcal I_{1,2}
\qquad\forall I_2\in\mathcal I_2.
\]

Therefore,
\[
\mathcal I_1\subseteq\mathcal I_{1,2},
\qquad
\mathcal I_2\subseteq\mathcal I_{1,2}.
\]

Since $\sigma(\mathcal I_{1,2})$ is a $\sigma$-algebra containing both $\mathcal I_1$ and $\mathcal I_2$, by the minimality of generated $\sigma$-algebras we obtain
\[
\mathcal G_1\subseteq\sigma(\mathcal I_{1,2}),
\qquad
\mathcal G_2\subseteq\sigma(\mathcal I_{1,2}).
\]

\emph{We now prove that $\sigma(\mathcal I_{1,2})$ and $\mathcal G_3$ are independent.}

Let $A\in\mathcal I_{1,2}$ and $I_3\in\mathcal I_3$. By definition, there exist $I_1\in\mathcal I_1$ and $I_2\in\mathcal I_2$ such that
\[
A=I_1\cap I_2.
\]

Therefore,
\[
\begin{aligned}
\mathbb P(A\cap I_3)
&=\mathbb P(I_1\cap I_2\cap I_3) =\mathbb P(I_1)\mathbb P(I_2)\mathbb P(I_3) =\mathbb P(I_1\cap I_2)\mathbb P(I_3) =\mathbb P(A)\mathbb P(I_3),
\end{aligned}
\]
where the second equality follows from our initial assumption, and the third from the independence of $\mathcal G_1$ and $\mathcal G_2$, established above.

Since $A\in\mathcal I_{1,2}$ and $I_3\in\mathcal I_3$ were arbitrary,
\[
\begin{gathered}
\mathbb P(A\cap I_3)=\mathbb P(A)\mathbb P(I_3),\\
\forall A\in\mathcal I_{1,2},\ I_3\in\mathcal I_3.
\end{gathered}
\]

Since $\mathcal I_{1,2}$ and $\mathcal I_3$ are $\pi$-systems, by Lemma 4.1 we conclude that
\[
\sigma(\mathcal I_{1,2})
\text{ and }\mathcal G_3
\text{ are independent}.
\]

Finally, let
\[
G_k\in\mathcal G_k,
\qquad k=1,2,3.
\]

Since $\mathcal G_1,\mathcal G_2\subseteq\sigma(\mathcal I_{1,2})$, and $\sigma(\mathcal I_{1,2})$ is closed under intersections, we have
\[
G_1\cap G_2\in\sigma(\mathcal I_{1,2}).
\]

Therefore,
\[
\begin{aligned}
\mathbb P(G_1\cap G_2\cap G_3)
&=\mathbb P((G_1\cap G_2)\cap G_3) =\mathbb P(G_1\cap G_2)\mathbb P(G_3) =\mathbb P(G_1)\mathbb P(G_2)\mathbb P(G_3),
\end{aligned}
\]
where the second equality follows from the independence of $\sigma(\mathcal I_{1,2})$ and $\mathcal G_3$, and the third from the independence of $\mathcal G_1$ and $\mathcal G_2$.

Since $G_k\in\mathcal G_k$, $k=1,2,3$, were arbitrary, we conclude that
\[
\begin{gathered}
\mathbb P(G_1\cap G_2\cap G_3)
=
\prod_{k=1}^{3}\mathbb P(G_k),\\
\forall G_k\in\mathcal G_k.
\end{gathered}
\]

Thus, by Definition 4.1, the $\sigma$-algebras $\mathcal G_1,\mathcal G_2,\mathcal G_3$ are independent.

Therefore, the equivalence is proved.
%
% !TEX root = ../main.tex
\subsection*{HW 38 --- Part II: Why Is the Assumption $\Omega\in\mathcal I_k$ Necessary?}%
\manualtoc{subsection}{\texorpdfstring{HW 38 --- Part II: Why Is the Assumption $\Omega\in\mathcal I_k$ Necessary?}{HW 38 --- Part II: Why Is the Assumption Omega in I k Necessary?}}{studiabile-nav-0262}%


In Part I, the assumption $\Omega\in\mathcal I_3$ allowed us to choose $I_3=\Omega$ and establish the independence of $\mathcal G_1$ and $\mathcal G_2$.

However, pairwise independence alone does not imply joint independence. We first illustrate this fact and then construct a counterexample showing why the assumption $\Omega\in\mathcal I_k$ is needed.

\subsubsection*{Counterexample 1 --- Pairwise Independence Does Not Imply Joint Independence}%
\manualtoc{subsubsection}{\texorpdfstring{Counterexample 1 --- Pairwise Independence Does Not Imply Joint Independence}{Counterexample 1 --- Pairwise Independence Does Not Imply Joint Independence}}{studiabile-nav-0263}%


Let $X_1$ and $X_2$ be two independent and identically distributed random variables such that
\[
\mathbb P(X_i=1)
=
\mathbb P(X_i=-1)
=
\frac12,
\qquad i=1,2.
\]

Define
\[
X_3:=X_1X_2.
\]

\emph{We first prove that $X_1,X_2,X_3$ are pairwise independent.}

By assumption, $X_1$ and $X_2$ are independent. Therefore, it remains to consider the pairs $(X_1,X_3)$ and $(X_2,X_3)$.

Notice that $X_3=1$ if and only if $X_1$ and $X_2$ have the same sign. Thus,
\[
\{X_3=1\}
=
\bigl(\{X_1=1\}\cap\{X_2=1\}\bigr)
\cup
\bigl(\{X_1=-1\}\cap\{X_2=-1\}\bigr).
\]

Since these two events are disjoint, by additivity and the independence of $X_1$ and $X_2$,
\[
\begin{aligned}
\mathbb P(X_3=1)
&=\mathbb P(X_1=1,X_2=1) \quad+\mathbb P(X_1=-1,X_2=-1)\\&=\mathbb P(X_1=1)\mathbb P(X_2=1) \quad+\mathbb P(X_1=-1)\mathbb P(X_2=-1) =\frac14+\frac14 =\frac12.
\end{aligned}
\]

\emph{We now prove that $X_1$ and $X_3$ are independent.}

Since each $X_i$ takes only the values $-1$ and $1$, we have
\[
\sigma(X_i)
=
\{\emptyset,\Omega,\{X_i=1\},\{X_i=-1\}\},
\qquad i=1,2,3.
\]

Therefore, the collection
\[
\mathcal J_i:=\bigl\{\{X_i=1\}\bigr\}
\]

is a $\pi$-system generating $\sigma(X_i)$. Indeed, it is closed under intersections, and its generated $\sigma$-algebra also contains
\[
\{X_i=1\}^c=\{X_i=-1\}.
\]

Thus, by Lemma 4.1, it is sufficient to prove that
\[
\mathbb P(X_1=1,X_3=1)
=
\mathbb P(X_1=1)\mathbb P(X_3=1).
\]

Notice that, since $X_3=X_1X_2$,
\[
\{X_1=1,X_3=1\}
=
\{X_1=1,X_2=1\}.
\]
Indeed, when $X_1=1$, the equality $X_3=1$ holds if and only if $X_2=1$.

Therefore,
\[
\begin{aligned}
\mathbb P(X_1=1,X_3=1)
&=\mathbb P(X_1=1,X_2=1) =\mathbb P(X_1=1)\mathbb P(X_2=1) =\frac14 =\mathbb P(X_1=1)\mathbb P(X_3=1),
\end{aligned}
\]

where the second equality follows from independence, and the last from $\mathbb P(X_3=1)=1/2$, established above.

Hence, by Lemma 4.1, $\sigma(X_1)$ and $\sigma(X_3)$ are independent.

Similarly, exchanging $X_1$ and $X_2$, we obtain
\[
\begin{aligned}
\mathbb P(X_2=1,X_3=1)
=\mathbb P(X_1=1,X_2=1) =\frac14 =\mathbb P(X_2=1)\mathbb P(X_3=1).
\end{aligned}
\]

Therefore, again by Lemma 4.1, $\sigma(X_2)$ and $\sigma(X_3)$ are independent.

We conclude that
\[
X_1,X_2,X_3
\text{ are pairwise independent}.
\]

\emph{We now prove that they are not jointly independent.}

Since $X_3=X_1X_2$, we have
\[
\{X_1=1,X_2=1,X_3=1\}
=
\{X_1=1,X_2=1\}.
\]

Therefore,
\[
\begin{aligned}
\mathbb P(X_1=1,X_2=1,X_3=1)
=\mathbb P(X_1=1,X_2=1) =\frac14.
\end{aligned}
\]

On the other hand,
\[
\begin{aligned}
\prod_{i=1}^{3}\mathbb P(X_i=1)
=\frac12\cdot\frac12\cdot\frac12 =\frac18.
\end{aligned}
\]

Since $\frac14\neq\frac18$, we obtain
\[
\mathbb P(X_1=1,X_2=1,X_3=1)
\neq
\prod_{i=1}^{3}\mathbb P(X_i=1).
\]

Thus, by Definition 4.1, the three $\sigma$-algebras $\sigma(X_1)$, $\sigma(X_2)$ and $\sigma(X_3)$ are not jointly independent.

Therefore, this counterexample shows that pairwise independence does not imply joint independence.

In particular, it explains why, in Part I of HW 38, establishing the independence of $\mathcal G_1$ and $\mathcal G_2$ was not sufficient.

\subsubsection*{Counterexample 2 --- Necessity of $\Omega\in\mathcal I_k$}%
\manualtoc{subsubsection}{\texorpdfstring{Counterexample 2 --- Necessity of $\Omega\in\mathcal I_k$}{Counterexample 2 --- Necessity of Omega in I k}}{studiabile-nav-0264}%


\emph{We now prove that, without the assumption $\Omega\in\mathcal I_k$, the factorisation on the three $\pi$-systems does not necessarily imply the independence of the generated $\sigma$-algebras.}

Consider the probability space
\[
\Omega=\{0,1\},
\qquad
\mathcal F=\mathcal P(\Omega),
\]

with
\[
\mathbb P(\{0\})=\mathbb P(\{1\})=\frac12.
\]

Let $A:=\{1\}$, so that
\[
\mathbb P(A)=\frac12.
\]

Define
\[
\mathcal I_1=\mathcal I_2=\{\Omega,A\},
\qquad
\mathcal I_3=\{\emptyset\}.
\]

First, notice that these collections are $\pi$-systems.

Indeed, $\{\Omega,A\}$ is closed under intersections, since
\[
\Omega\cap A=A,
\qquad
A\cap A=A,
\qquad
\Omega\cap\Omega=\Omega.
\]

Moreover, $\{\emptyset\}$ is a non-empty collection closed under intersections, because
\[
\emptyset\cap\emptyset=\emptyset.
\]

However, although $\Omega\in\mathcal I_1$ and $\Omega\in\mathcal I_2$,
\[
\Omega\notin\mathcal I_3.
\]

\emph{We first prove that the factorisation in HW 38 holds on these three $\pi$-systems.}

Let $I_k\in\mathcal I_k$, $k=1,2,3$.

Since $\mathcal I_3=\{\emptyset\}$, necessarily $I_3=\emptyset$.

Therefore,
\[
\begin{aligned}
\mathbb P(I_1\cap I_2\cap I_3)
=\mathbb P(I_1\cap I_2\cap\emptyset) =\mathbb P(\emptyset) =0,
\end{aligned}
\]

where the second equality follows from the fact that the intersection of any set with the empty set is empty.

On the other hand,
\[
\mathbb P(I_1)\mathbb P(I_2)\mathbb P(I_3)
=
\mathbb P(I_1)\mathbb P(I_2)\cdot0
=
0.
\]

Consequently,
\[
\mathbb P(I_1\cap I_2\cap I_3)
=
\prod_{k=1}^{3}\mathbb P(I_k).
\]

for every $I_k\in\mathcal I_k$, $k=1,2,3$.

\emph{We now prove that the generated $\sigma$-algebras are not independent.}

Let
\[
\mathcal G_k:=\sigma(\mathcal I_k),
\qquad k=1,2,3.
\]

Since $A\in\mathcal I_1$ and $A\in\mathcal I_2$, we have
\[
A\in\mathcal G_1,
\qquad
A\in\mathcal G_2.
\]

Moreover, although $\Omega\notin\mathcal I_3$, we necessarily have
\[
\Omega\in\mathcal G_3,
\]

because every $\sigma$-algebra on $\Omega$ contains $\Omega$.

Therefore, we can choose
\[
G_1=A,\qquad G_2=A,\qquad G_3=\Omega.
\]

If the three $\sigma$-algebras were independent, by Definition 4.1 we would have
\[
\mathbb P(A\cap A\cap\Omega)
=
\mathbb P(A)\mathbb P(A)\mathbb P(\Omega).
\]

However,
\[
\begin{aligned}
\mathbb P(A\cap A\cap\Omega)
=\mathbb P(A) =\frac12,
\end{aligned}
\]

where the first equality follows from $A\cap A\cap\Omega=A$.

On the other hand,
\[
\begin{aligned}
\mathbb P(A)\mathbb P(A)\mathbb P(\Omega)
=\frac12\cdot\frac12\cdot1 =\frac14.
\end{aligned}
\]

Since $\frac12\neq\frac14$, we obtain
\[
\mathbb P(A\cap A\cap\Omega)
\neq
\mathbb P(A)\mathbb P(A)\mathbb P(\Omega).
\]

Thus, by Definition 4.1, the $\sigma$-algebras $\mathcal G_1,\mathcal G_2,\mathcal G_3$ are not independent.

Notice that this does not contradict the factorisation established above. Indeed, that factorisation holds only for events belonging to the original $\pi$-systems. Since $\Omega\notin\mathcal I_3$, we could not choose $I_3=\Omega$ there, whereas $\Omega$ necessarily belongs to the generated $\sigma$-algebra $\mathcal G_3$.

Therefore, the factorisation holds on the three $\pi$-systems but fails on their generated $\sigma$-algebras.

Hence, the assumption $\Omega\in\mathcal I_k$, $k=1,2,3$, is necessary to guarantee the equivalence in HW 38.

Therefore, the statement is proved.
%
% !TEX root = ../main.tex
\subsection*{Lemma 4.2 --- Second Borel--Cantelli's Lemma}%
\manualtoc{subsection}{\texorpdfstring{Lemma 4.2 --- Second Borel--Cantelli's Lemma}{Lemma 4.2 --- Second Borel--Cantelli's Lemma}}{studiabile-nav-0265}%


\subsubsection*{Statement}%
%


Let $(\Omega,\mathcal F,\mathbb P)$ be a probability space, and let $(E_n)_{n\in\mathbb N}\subseteq\mathcal F$ be a sequence of independent events such that
\[
\sum_{n=1}^{\infty}\mathbb P(E_n)=+\infty.
\]

Then
\[
\mathbb P\left(\limsup_{n\to\infty}E_n\right)=1.
\]

In other words, with probability one, infinitely many of the events $E_n$ occur.

\subsubsection*{Expanded Proof of Lemma 4.2}%
%


Set
\[
E:=\limsup_{n\to\infty}E_n.
\]

By definition of the limit superior of a sequence of sets and De Morgan's laws, we have
\[
\begin{aligned}
E^c
&=\left(
\bigcap_{n=1}^{\infty}
\bigcup_{m\ge n}E_m
\right)^c =\bigcup_{n=1}^{\infty}
\left(\bigcup_{m\ge n}E_m\right)^c =\bigcup_{n=1}^{\infty}
\bigcap_{m\ge n}E_m^c =\liminf_{n\to\infty}E_n^c,
\end{aligned}
\]
where the second and third equalities follow from De Morgan's laws, and the last from the definition of the limit inferior.

We want to prove that $\mathbb P(E^c)=0$, which will imply $\mathbb P(E)=1$.

Set
\[
p_n:=\mathbb P(E_n),
\qquad n\in\mathbb N.
\]

By assumption,
\[
\sum_{n=1}^{\infty}p_n=+\infty.
\]

Fix an arbitrary $n\in\mathbb N$. We first calculate
\[
\mathbb P\left(\bigcap_{m\ge n}E_m^c\right).
\]

For every $k\ge0$, define
\[
A_k:=\bigcap_{m=n}^{n+k}E_m^c.
\]

Notice that $(A_k)_{k\ge0}$ is a decreasing sequence of measurable sets. Indeed,
\[
A_{k+1}
=A_k\cap E_{n+k+1}^c
\subseteq A_k.
\]

Moreover,
\[
\bigcap_{k=0}^{\infty}A_k
=
\bigcap_{m\ge n}E_m^c.
\]

Since $\mathbb P$ is a probability measure, we also have
\[
\mathbb P(A_0)\le\mathbb P(\Omega)=1<\infty.
\]

Thus, all the assumptions of Lemma 1.6 are satisfied. Therefore, by continuity of measures from above,
\[
\mathbb P\left(\bigcap_{m\ge n}E_m^c\right)
=
\lim_{k\to\infty}
\mathbb P\left(\bigcap_{m=n}^{n+k}E_m^c\right).
\]

Now, since the events $(E_n)_{n\in\mathbb N}$ are independent, their complements are also independent. Indeed, by Definition 4.3 and HW 36,
\[
E_m^c\in\sigma(\mathbf 1_{E_m})
\qquad\forall m\in\mathbb N,
\]
and the $\sigma$-algebras generated by the indicator functions of the independent events are themselves independent.

Consequently, for every finite $k\ge0$,
\[
\begin{aligned}
\mathbb P\left(\bigcap_{m=n}^{n+k}E_m^c\right)
&=\prod_{m=n}^{n+k}\mathbb P(E_m^c) =\prod_{m=n}^{n+k}(1-p_m),
\end{aligned}
\]
where the first equality follows from the independence of the complementary events, and the second from
\[
\mathbb P(E_m^c)=1-\mathbb P(E_m)=1-p_m.
\]

Combining this equality with continuity from above, we obtain
\[
\mathbb P\left(\bigcap_{m\ge n}E_m^c\right)
=
\lim_{k\to\infty}
\prod_{m=n}^{n+k}(1-p_m).
\]

We now use the inequality
\[
1-x\le e^{-x},\qquad x\ge0.
\]

Since $p_m=\mathbb P(E_m)\in[0,1]$, applying this inequality with $x=p_m$ gives
\[
0\le1-p_m\le e^{-p_m}.
\]

Therefore, multiplying these inequalities for $m=n,\ldots,n+k$, we obtain
\[
\begin{aligned}
\prod_{m=n}^{n+k}(1-p_m)
&\le\prod_{m=n}^{n+k}e^{-p_m} =\exp\left(-\sum_{m=n}^{n+k}p_m\right),
\end{aligned}
\]
where the last equality follows from the properties of the exponential function.

Consequently,
\[
0\le
\mathbb P\left(\bigcap_{m\ge n}E_m^c\right)
\le
\lim_{k\to\infty}
\exp\left(-\sum_{m=n}^{n+k}p_m\right).
\]

Notice that, since
\[
\sum_{m=1}^{\infty}p_m=+\infty
\]
and each $p_m\in[0,1]$, removing finitely many terms does not change the divergence of the series. Hence, for every fixed $n\in\mathbb N$,
\[
\sum_{m=n}^{\infty}p_m=+\infty.
\]

Therefore,
\[
\lim_{k\to\infty}
\exp\left(-\sum_{m=n}^{n+k}p_m\right)=0,
\]
since $e^{-t}\to0$ as $t\to+\infty$.

Thus, by the inequality established above,
\[
0\le
\mathbb P\left(\bigcap_{m\ge n}E_m^c\right)
\le0.
\]

We conclude that
\[
\mathbb P\left(\bigcap_{m\ge n}E_m^c\right)=0
\qquad\forall n\in\mathbb N.
\]

Finally, recall that we established at the beginning of the proof that
\[
E^c=
\bigcup_{n=1}^{\infty}\bigcap_{m\ge n}E_m^c.
\]

Since each event in this countable union has probability zero, by Lemma 1.7 (a countable union of null sets is a null set), we obtain
\[
\mathbb P(E^c)=0.
\]

Therefore,
\[
\begin{aligned}
\mathbb P(E)
=1-\mathbb P(E^c) =1.
\end{aligned}
\]

Recalling that $E=\limsup_{n\to\infty}E_n$, we conclude that
\[
\mathbb P\left(\limsup_{n\to\infty}E_n\right)=1.
\]

Therefore, the statement is proved.
%
% !TEX root = ../main.tex
% Both entries may be inserted into the existing Probability for Finance notes.

\subsection*{Example 2 (Williams, Example 4.4) --- Application of the Borel--Cantelli Lemmas}%
\manualtoc{subsection}{\texorpdfstring{Example 2 (Williams, Example 4.4) --- Application of the Borel--Cantelli Lemmas}{Example 2 (Williams, Example 4.4) --- Application of the Borel--Cantelli Lemmas}}{studiabile-nav-0268}%


\subsubsection*{Statement}%
%


Let $(X_n)_{n\in\mathbb N}$ be a sequence of independent and identically distributed random variables such that
\[
X_n\sim\operatorname{Exp}(1),\qquad n\in\mathbb N.
\]
For $\alpha>0$, determine the probability of the event
\[
E_\alpha:=\{\omega\in\Omega:X_n(\omega)>\alpha\log n,\ \text{i.o.}\},
\]
where i.o. means infinitely often. Then study the almost sure behaviour of
\[
L(\omega):=\limsup_{n\to\infty}\frac{X_n(\omega)}{\log n}.
\]

\subsubsection*{Expanded Solution of Example 2}%
%


Fix $\alpha>0$, and define
\[
E_n^\alpha:=\{\omega\in\Omega:X_n(\omega)>\alpha\log n\},
\qquad n\in\mathbb N.
\]
The survival function of $X_n$ satisfies
\[
\mathbb P(X_n>t)=e^{-t},\qquad t\ge0.
\]
This follows from $X_n\sim\operatorname{Exp}(1)$. Consequently,
\[
\begin{aligned}
\mathbb P(E_n^\alpha)
=\mathbb P(X_n>\alpha\log n) =e^{-\alpha\log n} =n^{-\alpha}.
\end{aligned}
\]
The first equality uses the definition of $E_n^\alpha$, the second the exponential survival function (since $\alpha\log n\ge0$), and the last the identity $e^{a\log n}=n^a$.

Moreover,
\[
E_n^\alpha
=X_n^{-1}((\alpha\log n,\infty))
\in\sigma(X_n)\subseteq\mathcal F.
\]
The equality is the definition of a pre-image; the membership follows from the definition of $\sigma(X_n)$, and the final inclusion from the $\mathcal F$-measurability of $X_n$. Since the $X_n$ are independent, the events $(E_n^\alpha)_{n\in\mathbb N}$ are independent by Definition 4.2.

We now express $E_\alpha$ as a limit superior:
\[
\begin{aligned}
E_\alpha
&=\{\omega\in\Omega:X_n(\omega)>\alpha\log n
  \text{ for infinitely many }n\}\\&=\{\omega\in\Omega:\forall N\in\mathbb N,\ \exists m\ge N
  \text{ such that }X_m(\omega)>\alpha\log m\}\\&=\bigcap_{N=1}^{\infty}\bigcup_{m\ge N}
  \{\omega\in\Omega:X_m(\omega)>\alpha\log m\} =\bigcap_{N=1}^{\infty}\bigcup_{m\ge N}E_m^\alpha =\limsup_{n\to\infty}E_n^\alpha.
\end{aligned}
\]
The first equality states the definition of $E_\alpha$. The second expresses infinitely often: after every $N$ there must be some $m\ge N$ for which the inequality holds. The third translates ``for every $N$'' into an intersection and ``there exists $m\ge N$'' into a union. The fourth substitutes the definition of $E_m^\alpha$, and the last is the definition of the limit superior of sets.

We will use formula (4.3):
\[
\mathbb P(E_\alpha)=
\begin{cases}
0,&\alpha>1,\\
1,&\alpha\in(0,1].
\end{cases}
\]
As we will prove in HW~39, this follows from the two Borel--Cantelli lemmas and the convergence criterion for the $p$-series $\sum_{n=1}^{\infty}n^{-\alpha}$.

We now consider
\[
L(\omega):=\limsup_{n\to\infty}\frac{X_n(\omega)}{\log n},
\]
where the ratios are defined for $n\ge2$, because $\log 1=0$. For each $n\ge2$, the function $X_n/\log n$ is measurable, since $X_n$ is measurable and $\log n$ is a positive constant. Moreover,
\[
L=\inf_{N\ge2}\sup_{n\ge N}\frac{X_n}{\log n}.
\]
This is the definition of the numerical limit superior; since countable suprema and infima of measurable functions are measurable, $L$ is an extended-real-valued measurable random variable.

\textit{We first prove that $\mathbb P(L\ge1)=1$.}

Let $\omega_0\in E_1$. We have
\[
\frac{X_n(\omega_0)}{\log n}>1
\qquad\text{for infinitely many }n\ge2.
\]
Indeed, the definition of $E_1$ gives $X_n(\omega_0)>\log n$ infinitely often, and division by $\log n>0$ preserves the inequality. Consequently,
\[
L(\omega_0)\ge1.
\]
This follows from the definition of the limit superior: a sequence exceeding $1$ infinitely often cannot have limit superior below $1$. Since $\omega_0\in E_1$ was arbitrary,
\[
E_1\subseteq\{L\ge1\}.
\]
We therefore obtain
\[
1\ge\mathbb P(L\ge1)\ge\mathbb P(E_1)=1.
\]
The first bound holds for every probability, the second inequality follows from the inclusion just proved and monotonicity of probability, and the last equality uses formula (4.3) with $\alpha=1$. Hence $\mathbb P(L\ge1)=1$.

\textit{We now prove that $\mathbb P(L>1)=0$.}

First, notice that
\[
\{L>1\}
=\bigcup_{k=1}^{\infty}\left\{L\ge1+\frac2k\right\}
=\bigcup_{k=1}^{\infty}A_k,
\]
where, for each fixed $k\in\mathbb N$, we define
\[
A_k:=\left\{L\ge1+\frac2k\right\}.
\]
The first equality follows from $2/k\to0$ and the Archimedean property of the real numbers (including the case $L=+\infty$); the second follows from the definition of $A_k$.

Now fix $k\in\mathbb N$, and define
\[
B_k:=\left\{X_n>\left(1+\frac1k\right)\log n,\ \text{i.o.}\right\}.
\]
We want to show that
\[
A_k\subseteq B_k
\quad\Longleftrightarrow\quad
A_k\cap B_k^c=\emptyset.
\]
The equivalence holds because an inclusion $A_k\subseteq B_k$ means that no outcome can belong to $A_k$ without belonging to $B_k$.

Suppose, for a contradiction, that there exists
\[
\omega_0\in A_k\cap B_k^c.
\]
From the membership $\omega_0\in A_k$, we obtain
\[
L(\omega_0)\ge1+\frac2k.
\]
This is the definition of $A_k$. On the other hand,
\[
X_n(\omega_0)>\left(1+\frac1k\right)\log n
\qquad\text{only finitely often}.
\]
This follows from $\omega_0\in B_k^c$: the outcome does not belong to the event where the inequality holds infinitely often. Thus, for some $N\ge2$,
\[
X_n(\omega_0)\le\left(1+\frac1k\right)\log n
\qquad\forall n\ge N.
\]
Indeed, after the last index (if any) at which the strict inequality holds, its negation must hold for every subsequent index. Dividing yields
\[
\frac{X_n(\omega_0)}{\log n}\le1+\frac1k
\qquad\forall n\ge N.
\]
The inequality is preserved because $\log n>0$ for $n\ge2$. Therefore,
\[
L(\omega_0)\le1+\frac1k.
\]
This follows from the definition of limit superior, since every sufficiently late term is bounded by $1+1/k$, and contradicts the membership $\omega_0\in A_k$ established above. Thus,
\[
A_k\cap B_k^c=\emptyset,
\qquad A_k\subseteq B_k.
\]
The contradiction rules out any outcome in $A_k\cap B_k^c$, which is equivalent to the desired inclusion.

Consequently,
\[
\begin{aligned}
\mathbb P(A_k)
\le\mathbb P(B_k) =\mathbb P(E_{1+1/k}) =0.
\end{aligned}
\]
The inequality follows from $A_k\subseteq B_k$ and monotonicity of probability, the first equality from the definition of $E_\alpha$, and the last from formula (4.3), since $1+1/k>1$.

Since $k$ was arbitrary, we obtain
\[
\mathbb P(L>1)
=\mathbb P\left(\bigcup_{k=1}^{\infty}A_k\right)
=0.
\]
The first equality uses the identity for $\{L>1\}$, and the second follows from Lemma 1.7 because each $A_k$ has probability zero.

Finally,
\[
\begin{aligned}
\mathbb P(L=1)
=\mathbb P(L\ge1)-\mathbb P(L>1) =1-0 =1.
\end{aligned}
\]
The first equality follows from the disjoint decomposition $\{L\ge1\}=\{L=1\}\cup\{L>1\}$ and additivity of probability; the other equalities use the two probabilities already established. Recalling the definition of $L$, we conclude that
\[
\mathbb P\left(\limsup_{n\to\infty}\frac{X_n}{\log n}=1\right)=1.
\]
\subsection*{HW 39 --- Proof of Formula (4.3)}%
\manualtoc{subsection}{\texorpdfstring{HW 39 --- Formula (4.3)}{HW 39 --- Formula (4.3)}}{studiabile-nav-0271}%


\subsubsection*{Statement}%
%


Let $(X_n)_{n\in\mathbb N}$ be independent and identically distributed random variables with
\[
X_n\sim\operatorname{Exp}(1).
\]
For $\alpha>0$, let
\[
E_\alpha:=\{X_n>\alpha\log n,\ \text{i.o.}\}.
\]
Prove formula (4.3):
\[
\mathbb P(E_\alpha)=
\begin{cases}
0,&\alpha>1,\\
1,&\alpha\in(0,1].
\end{cases}
\]

\subsubsection*{Expanded Proof of HW 39}%
%


Fix $\alpha>0$, and consider
\[
E_n^\alpha:=\{X_n>\alpha\log n\},\qquad n\in\mathbb N.
\]
As established in Example 2, these events are measurable and independent, and
\[
\mathbb P(E_n^\alpha)=n^{-\alpha},
\qquad E_\alpha=\limsup_{n\to\infty}E_n^\alpha.
\]
The first identity was obtained from the exponential survival function; the second from the meaning of infinitely often and the definition of the limit superior. Hence,
\[
\sum_{n=1}^{\infty}\mathbb P(E_n^\alpha)
=\sum_{n=1}^{\infty}\frac1{n^\alpha}.
\]
The equality follows by substituting the probability of each $E_n^\alpha$.

If $\alpha>1$, the $p$-series criterion gives
\[
\sum_{n=1}^{\infty}\mathbb P(E_n^\alpha)<\infty.
\]
Consequently,
\[
\begin{aligned}
\mathbb P(E_\alpha)
=\mathbb P\left(\limsup_{n\to\infty}E_n^\alpha\right) =0.
\end{aligned}
\]
The first equality uses the representation of $E_\alpha$, and the second follows from Lemma 2.3 (First Borel--Cantelli), whose assumptions hold because the events are measurable and their probabilities have finite sum. Independence is not required here.

If $0<\alpha\le1$, the same $p$-series diverges, so
\[
\sum_{n=1}^{\infty}\mathbb P(E_n^\alpha)=+\infty.
\]
At $\alpha=1$ this is the harmonic series. Therefore,
\[
\begin{aligned}
\mathbb P(E_\alpha)
=\mathbb P\left(\limsup_{n\to\infty}E_n^\alpha\right) =1.
\end{aligned}
\]
The first equality uses the same representation of $E_\alpha$, and the second follows from Lemma 4.2 (Second Borel--Cantelli), because the events are independent and the sum of their probabilities diverges.

Both cases of (4.3) follow.
%
% !TEX root = ../main.tex
\subsection*{Fundamental Question for Modelling}%
\manualtoc{subsection}{\texorpdfstring{Fundamental Question for Modelling}{Fundamental Question for Modelling}}{studiabile-nav-0274}%


So far, we have considered sequences of random variables with arbitrary distributions. A natural question is whether we can actually construct a sequence of independent random variables $(X_n)_{n\in\mathbb N}$, each with a prescribed CDF $F_n$.

The answer is yes. Indeed, let $(U_n)_{n\in\mathbb N}$ be a sequence of i.i.d. random variables such that
\[
U_n\sim\operatorname{Unif}(0,1).
\]

Define
\[
X_n:=F_n^{-1}(U_n),
\qquad n\in\mathbb N,
\]
where $F_n^{-1}$ denotes the generalised inverse of the CDF $F_n$.

Then
\[
X_n\sim F_n,
\]
by the inverse transform method. Moreover, the random variables $(X_n)_{n\in\mathbb N}$ are independent, since each $X_n$ is a measurable function of $U_n$, and measurable transformations of independent random variables preserve independence.

\subsection*{Stochastic Processes}%
\manualtoc{subsection}{\texorpdfstring{Stochastic Processes}{Stochastic Processes}}{studiabile-nav-0275}%


A stochastic process on a probability space $(\Omega,\mathcal F,\mathbb P)$ is a family of random variables
\[
(Y_t)_{t\in\mathcal I},
\]
indexed by a parameter $t$, usually interpreted as time.

For example, the index set $\mathcal I$ may be $\mathbb N$, corresponding to discrete time, or $[0,\infty)$, corresponding to continuous time.

\subsection*{Existence of Stochastic Processes}%
\manualtoc{subsection}{\texorpdfstring{Existence of Stochastic Processes}{Existence of Stochastic Processes}}{studiabile-nav-0276}%


A natural question is whether a stochastic process can be constructed with a prescribed distributional structure.

Under suitable consistency conditions, the answer is yes, as established by the Kolmogorov Existence Theorem.

The technical details of this theorem are beyond the scope of the present discussion.
%
% !TEX root = ../main.tex
\subsection*{Example 3 --- Markov Chains}%
\manualtoc{subsection}{\texorpdfstring{Example 3 --- Markov Chains}{Example 3 --- Markov Chains}}{studiabile-nav-0277}%

\subsubsection*{Statement}%
%


Let $E$ be a finite or countable set, and let $P=(p_{ij})_{i,j\in E}$ be a square matrix with positive entries satisfying
\[
\sum_{j\in E}p_{ij}=1,
\qquad\forall i\in E.
\]

Such a matrix is called a stochastic matrix.

Let $\mu$ be a probability measure on $E$, and set
\[
\mu_i:=\mu(\{i\}),
\qquad i\in E.
\]

A time-homogeneous Markov chain $Z=(Z_n)_{n\in\mathbb N\cup\{0\}}$ on $E$, with initial distribution $\mu$ and transition matrix $P$, is a stochastic process satisfying
\[
\mathbb P(Z_0=i_0,Z_1=i_1,\ldots,Z_n=i_n)
=
\mu_{i_0}p_{i_0i_1}\cdots p_{i_{n-1}i_n}
\]
for every $n\in\mathbb N$ and every choice of states $i_0,\ldots,i_n\in E$.

We want to construct such a process using independent random variables.

\subsubsection*{Expanded Construction and Proof of Example 3}%
%


First, let $Z_0$ be an $E$-valued random variable with distribution $\mu$, so that
\[
\mathbb P(Z_0=i)=\mu_i,
\qquad\forall i\in E.
\]

Next, consider a family of $E$-valued random variables
\[
(Y_{i,n})_{i\in E,\ n\in\mathbb N}
\]
such that
\[
\mathbb P(Y_{i,n}=j)=p_{ij},
\qquad\forall i,j\in E,\ n\in\mathbb N.
\]

We assume that all the random variables $Y_{i,n}$ are mutually independent and also independent of $Z_0$.

Notice that this independence is an assumption of our initial construction: \textbf{the random variables $Y_{i,n}$ are chosen independently before defining the process $Z$}.

For each fixed state $i$ and time $n$, the random variable $Y_{i,n}$ takes values in $E$, according to the probabilities $p_{ij}$. Thus, the index $i$ identifies which random variable we are considering, whereas $j$ represents one of its possible values.

We now define the process recursively by setting
\[
Z_n(\omega):=Y_{Z_{n-1}(\omega),n}(\omega),
\qquad n\ge1.
\]

This means that the value of $Z_{n-1}(\omega)$ determines which random variable $Y_{i,n}$ is selected to obtain $Z_n(\omega)$.

In particular, if $Z_0(\omega)=i_0$, then
\[
Z_1(\omega)=Y_{i_0,1}(\omega).
\]

The equality follows from the recursive definition with $n=1$, after substituting $Z_0(\omega)=i_0$.

Notice that knowing $Z_0(\omega)=i_0$ identifies the random variable $Y_{i_0,1}$, \textbf{but does not determine its realised value}. We only know that
\[
\mathbb P(Y_{i_0,1}=j)=p_{i_0j},
\qquad j\in E.
\]

If we additionally consider an outcome $\omega$ for which $Z_1(\omega)=i_1$, then
\[
Z_2(\omega)=Y_{i_1,2}(\omega).
\]

This follows from the same recursive definition with $n=2$. Again, the condition $Z_1(\omega)=i_1$ identifies the variable used at the second step, without specifying its realised value.

Notice also that, although the variables $Y_{i,n}$ are independent, the resulting variables $Z_n$ are generally not independent, since the selection of the variable used at time $n$ depends on the previous state $Z_{n-1}$.

Nevertheless, the construction defines measurable random variables. Indeed, for every $j\in E$ and $n\ge1$,
\[
\{Z_n=j\}
=
\bigcup_{i\in E}
\left(
\{Z_{n-1}=i\}\cap\{Y_{i,n}=j\}
\right).
\]

The equality follows from the recursive definition: $Z_n=j$ precisely when the previous state equals some $i\in E$ and the corresponding variable $Y_{i,n}$ takes the value $j$. Since $E$ is finite or countable, the right-hand side is a countable union of measurable events. Starting from $Z_0$, this proves the measurability of each $Z_n$ by induction.

\emph{We now verify that the constructed process has the required joint distribution.}

We first consider the first two transitions.

Fix arbitrary states $i_0,i_1,i_2\in E$. We obtain
\[
\begin{aligned}
&\{Z_0=i_0,Z_1=i_1,Z_2=i_2\}\\
&\quad=
\{Z_0=i_0,Y_{i_0,1}=i_1,Z_2=i_2\}\\
&\quad=
\{Z_0=i_0,Y_{i_0,1}=i_1,Y_{i_1,2}=i_2\}.
\end{aligned}
\]

The first equality follows from the recursive definition: on the event $\{Z_0=i_0\}$, we have $Z_1=Y_{i_0,1}$. Therefore, requiring $Z_1=i_1$ is equivalent to requiring $Y_{i_0,1}=i_1$.

The second equality follows by applying the same reasoning to the next step. Once the first two conditions hold, we have $Z_1=i_1$, and therefore $Z_2=Y_{i_1,2}$. Hence, requiring $Z_2=i_2$ is equivalent to requiring $Y_{i_1,2}=i_2$.

\textbf{The important point is that $i_0,i_1,i_2$ are fixed states.} We are considering the event in which the process follows the particular path
\[
i_0\longrightarrow i_1\longrightarrow i_2.
\]

Consequently, the event on the right-hand side involves the three specific random variables
\[
Z_0,\qquad Y_{i_0,1},\qquad Y_{i_1,2},
\]
which are independent by our initial assumption. We are not using independence of $Z_0,Z_1,Z_2$.

Therefore,
\[
\begin{aligned}
&\mathbb P(Z_0=i_0,Z_1=i_1,Z_2=i_2)\\
&\quad=
\mathbb P(Z_0=i_0,Y_{i_0,1}=i_1,Y_{i_1,2}=i_2)\\
&\quad=
\mathbb P(Z_0=i_0)
\mathbb P(Y_{i_0,1}=i_1)
\mathbb P(Y_{i_1,2}=i_2)\\
&\quad=
\mu_{i_0}p_{i_0i_1}p_{i_1i_2}.
\end{aligned}
\]

The first equality follows from the equality of events established above. The second follows from the independence of $Z_0,Y_{i_0,1},Y_{i_1,2}$, and the last from their prescribed distributions.

We now extend the argument to an arbitrary number of transitions.

Fix $n\in\mathbb N$ and arbitrary states $i_0,\ldots,i_n\in E$.

Repeatedly applying the recursive definition gives
\[
\begin{aligned}
&\{Z_0=i_0,Z_1=i_1,\ldots,Z_n=i_n\}\\
&\quad=
\{Z_0=i_0\}
\cap
\bigcap_{r=1}^{n}
\{Y_{i_{r-1},r}=i_r\}.
\end{aligned}
\]

This follows from the same argument used for the first two transitions. For a fixed path, the condition $Z_{r-1}=i_{r-1}$ implies that the variable selected at time $r$ is $Y_{i_{r-1},r}$. Therefore, requiring $Z_r=i_r$ is equivalent to requiring $Y_{i_{r-1},r}=i_r$.

Since the states $i_0,\ldots,i_n$ are fixed, the event on the right-hand side involves the random variables
\[
Z_0,\,
Y_{i_0,1},\,
Y_{i_1,2},\,
\ldots,\,
Y_{i_{n-1},n}.
\]

These are independent by construction. In particular, the selected $Y$-variables are distinct because their time indices are different, even if some of the states coincide.

Consequently,
\[
\begin{aligned}
&\mathbb P(Z_0=i_0,Z_1=i_1,\ldots,Z_n=i_n)\\
&\quad=
\mathbb P\left(
\{Z_0=i_0\}
\cap
\bigcap_{r=1}^{n}
\{Y_{i_{r-1},r}=i_r\}
\right)\\
&\quad=
\mathbb P(Z_0=i_0)
\prod_{r=1}^{n}
\mathbb P(Y_{i_{r-1},r}=i_r)\\
&\quad=
\mu_{i_0}
\prod_{r=1}^{n}p_{i_{r-1}i_r}\\
&\quad=
\mu_{i_0}p_{i_0i_1}\cdots p_{i_{n-1}i_n}.
\end{aligned}
\]

The first equality follows from the representation of the event just established. The second follows from the independence of the selected random variables, the third from their prescribed distributions, and the last from expanding the finite product.

Since $n\in\mathbb N$ and the states $i_0,\ldots,i_n\in E$ were arbitrary, we have proved that
\[
\mathbb P(Z_0=i_0,Z_1=i_1,\ldots,Z_n=i_n)
=
\mu_{i_0}p_{i_0i_1}\cdots p_{i_{n-1}i_n}.
\]

Thus, the constructed process satisfies the required joint distribution and is a time-homogeneous Markov chain with initial distribution $\mu$ and transition matrix $P$.

Therefore, the statement is proved.
%
% !TEX root = ../main.tex
\subsection*{A Monkey Types Shakespeare}%
\manualtoc{subsection}{\texorpdfstring{A Monkey Types Shakespeare}{A Monkey Types Shakespeare}}{studiabile-nav-0280}%


\subsubsection*{Statement and construction}%
%


Consider a fixed text, such as the complete works of William Shakespeare (WS), consisting of $N\in\mathbb N$ symbols in a specified order. Write the target sequence as
\[
(w_1,w_2,\ldots,w_N).
\]
Let $(X_n)_{n\in\mathbb N}$ be a sequence of independent random variables, where $X_n$ is the symbol typed by the monkey at the $n$-th keystroke. We assume that each symbol needed to type WS has a probability uniformly bounded below by a constant $\varepsilon>0$ at every keystroke:
\[
\mathbb P(X_n=x)\ge\varepsilon>0
\qquad\text{for every }n\in\mathbb N\text{ and every symbol }x
\text{ appearing in WS}.
\]
This makes explicit the uniform-in-time meaning of the positive lower bound used in the lecture notes. No assumption that the $X_n$ have identical distributions is needed.

For $s\in\mathbb N\cup\{0\}$, let
\[
C_s:=\bigcap_{r=1}^{N}\{X_{s+r}=w_r\}.
\]
The event $C_s$ means that a complete copy of WS starts at keystroke $s+1$ and ends at keystroke $s+N$. Different occurrences may overlap; each occurrence is identified by its starting position.

Following the lecture notes, introduce the events
\[
\begin{aligned}
H&:=\{\text{the monkey types infinitely many copies of WS}\},\\
H_k&:=\{\text{the monkey types at least }k\text{ copies of WS}\},\\
H_{m,k}&:=\{\text{the monkey types at least }k\text{ copies of WS by time }m\},\\
H^{(m)}&:=\{\text{the monkey types infinitely many copies of WS after time }m\}.
\end{aligned}
\]
Here $H_{m,k}$ concerns copies completed no later than the $m$-th keystroke, whereas $H^{(m)}$ concerns copies whose starting positions are after that keystroke. In particular, $H_{m,k}$ depends only on $X_1,\ldots,X_m$, while $H^{(m)}$ depends only on $X_{m+1},X_{m+2},\ldots$. These are measurable events: each occurrence $C_s$ is a finite intersection of events determined by random variables, and the conditions involving finitely or infinitely many occurrences can be expressed through countable unions and intersections.

We want to prove
\[
\mathbb P(H)=1.
\]

\subsubsection*{Expanded proof}%
%


\textit{We first show that $\mathbb P(H)$ is either $0$ or $1$.}

Fix $m,k\in\mathbb N$. We obtain
\[
\mathbb P\bigl(H_{m,k}\cap H^{(m)}\bigr)
=\mathbb P(H_{m,k})\mathbb P\bigl(H^{(m)}\bigr).
\]
The equality follows from independence: $H_{m,k}$ is determined by the first $m$ symbols, whereas $H^{(m)}$ is determined by symbols typed strictly after time $m$; the corresponding families of random variables are independent.

As we will see in HW~40, deleting finitely many possible starting positions does not change whether infinitely many copies occur, so $H^{(m)}=H$. Consequently,
\[
\mathbb P(H_{m,k}\cap H)
=\mathbb P(H_{m,k})\mathbb P(H).
\]
This equality follows by substituting $H^{(m)}=H$ into the preceding identity.

As we will also see in HW~40,
\[
H_{m,k}\uparrow H_k
\quad\text{as }m\to\infty,
\qquad H_k\cap H=H.
\]
Therefore,
\[
\begin{aligned}
\mathbb P(H)
&=\mathbb P(H_k\cap H) =\lim_{m\to\infty}\mathbb P(H_{m,k}\cap H) =\lim_{m\to\infty}\bigl[\mathbb P(H_{m,k})\mathbb P(H)\bigr] =\mathbb P(H_k)\mathbb P(H).
\end{aligned}
\]
The first equality uses $H_k\cap H=H$. The second follows from Lemma~1.5 (continuity from below), since $H_{m,k}\cap H\uparrow H_k\cap H$. The third uses the independence identity already established, and the last follows from continuity from below applied to $H_{m,k}\uparrow H_k$.

As we will prove in HW~40,
\[
H_k\downarrow H
\quad\text{as }k\to\infty.
\]
Consequently,
\[
\mathbb P(H)=\mathbb P(H)^2.
\]
The equality follows by taking $k\to\infty$ in $\mathbb P(H)=\mathbb P(H_k)\mathbb P(H)$ and applying Lemma~1.6 (continuity from above). Since $\mathbb P(H)\in[0,1]$, we conclude that
\[
\mathbb P(H)\in\{0,1\}.
\]
Indeed, the equation $p=p^2$ is equivalent to $p(1-p)=0$.

\textit{We now prove that $\mathbb P(H)=1$.}

Divide the successive keystrokes into disjoint blocks of length $N$. For every $k\in\mathbb N$, define
\[
E_k:=\bigcap_{r=1}^{N}
\{X_{(k-1)N+r}=w_r\}.
\]
Thus $E_k$ is the event that the monkey types an entire copy of WS in the $k$-th block, consisting precisely of keystrokes $(k-1)N+1,\ldots,kN$.

We obtain
\[
\begin{aligned}
\mathbb P(E_k)
&=\prod_{r=1}^{N}\mathbb P(X_{(k-1)N+r}=w_r) \ge\varepsilon^N>0.
\end{aligned}
\]
The equality follows from independence of the $N$ random variables used in the block, while the inequality follows from the uniform lower bound on the probability of each required symbol. In particular, the same positive bound holds for every $k$.

Consequently,
\[
\sum_{k=1}^{\infty}\mathbb P(E_k)
\ge\sum_{k=1}^{\infty}\varepsilon^N
=+\infty.
\]
The inequality follows from the preceding bound, and the equality holds because the sum contains infinitely many copies of the fixed positive constant $\varepsilon^N$.

Moreover, the events $(E_k)_{k\in\mathbb N}$ are independent, because distinct blocks involve disjoint sets of the independent random variables $(X_n)_{n\in\mathbb N}$. Thus all the assumptions of Lemma~4.2, the Second Borel--Cantelli Lemma, are satisfied, and we obtain
\[
\mathbb P\left(\limsup_{k\to\infty}E_k\right)=1.
\]
As we will see in HW~40,
\[
\limsup_{k\to\infty}E_k\subseteq H.
\]
Therefore,
\[
1=\mathbb P\left(\limsup_{k\to\infty}E_k\right)
\le\mathbb P(H)\le1.
\]
The first equality is the conclusion of the Second Borel--Cantelli Lemma, the first inequality follows from the inclusion and monotonicity of probability, and the last inequality follows from $H\in\mathcal F$.

Hence,
\[
\mathbb P(H)=1.
\]
Therefore, the statement is proved. 

\label{hw-solution-40}
\subsection*{HW 40 --- Set-Theoretic Verification of the Events in the Monkey Example}%
\manualtoc{subsection}{\texorpdfstring{HW 40 --- Set-Theoretic Verification of the Events in the Monkey Example}{HW 40 --- Set-Theoretic Verification of the Events in the Monkey Example}}{studiabile-nav-0283}%


\subsubsection*{Statement}%
%


With the notation introduced in A Monkey Types Shakespeare, prove that
\begin{enumerate}[label=(\roman*)]
\item $H=H^{(m)}$;
\item $H_k\cap H=H$;
\item $\limsup_{k\to\infty}E_k\subseteq H$;
\item $H_k=\bigcup_{m=1}^{\infty}H_{m,k}$;
\item $H=\bigcap_{k=1}^{\infty}H_k$.
\end{enumerate}

\subsubsection*{Expanded proof of HW 40}%
%


Recall that
$C_s:=\bigcap_{r=1}^{N}\{X_{s+r}=w_r\},
\; s\in\mathbb N\cup\{0\},$
is the event that the monkey types one complete copy of WS starting at keystroke $s+1$.

\noindent\studyheading{\studylabel{(i)}} Recall that $H$ is the event of infinitely many occurrences of WS in the entire sequence, while $H^{(m)}$ concerns infinitely many occurrences whose first keystrokes are strictly after time $m$. We also allow $m=0$.

For every fixed threshold $q\ge0$, consider the event $\bigcup_{s\ge q}C_s$. This is the event that at least one complete copy of WS starts at a keystroke $s+1$ with $s\ge q$. Here $s$ labels a possible starting position and $q$ is a variable threshold; all indices are nonnegative integers.

For fixed $m\ge0$, we have
\[
\begin{aligned}
H^{(m)}&=\bigcap_{q=m}^{\infty}\bigcup_{s\ge q}C_s,\\
\omega\in H^{(m)}&\iff\forall q\ge m,\ \exists s\ge q:\ \omega\in C_s.
\end{aligned}
\]
The outer intersection requires an occurrence beyond every threshold $q\ge m$, which is precisely the meaning of infinitely many occurrences after time $m$.

In particular,
\[
H=H^{(0)}=\limsup_{s\to\infty}C_s.
\]
The first equality uses the case $m=0$, when the entire sequence is considered; the second is the definition of the limit superior of events.

Fix $m\in\mathbb N$. \emph{First, we prove that $H^{(m)}\subseteq H$.}

For every $0\le q<m$, we have
\[
H^{(m)}\subseteq\bigcup_{s\ge m}C_s
\subseteq\bigcup_{s\ge q}C_s.
\]
The first inclusion follows by taking the threshold $m$ in the intersection defining $H^{(m)}$; the second follows because $s\ge m$ implies $s\ge q$. Hence,
\[
H^{(m)}\subseteq H^{(0)}=H.
\]
The inclusion holds because the conditions with $q<m$ follow from the preceding inclusions, while those with $q\ge m$ already enter the definition of $H^{(m)}$; the equality was established above.

\noindent\textit{Conversely, we prove that $H\subseteq H^{(m)}$ by contradiction.}

Recall that
\[
H\subseteq H^{(m)}
\quad\Longleftrightarrow\quad
H\cap\bigl(H^{(m)}\bigr)^c=\emptyset.
\]

The equivalence follows because the inclusion means that no outcome can belong to $H$ without belonging to $H^{(m)}$.

Suppose, for a contradiction, that there exists
\[
\omega_0\in H\cap\bigl(H^{(m)}\bigr)^c.
\]

Since $\omega_0\in H$, the monkey types infinitely many copies of WS in the entire sequence. On the other hand, since $\omega_0\in\bigl(H^{(m)}\bigr)^c$, only finitely many copies start strictly after time $m$.

However, at most $m$ distinct copies can start at or before time $m$, corresponding to the starting indices $s=0,\ldots,m-1$. Therefore, infinitely many copies in total and only finitely many copies starting after time $m$ are incompatible.

We now express this contradiction formally.

We obtain
\[
\bigl(H^{(m)}\bigr)^c
=\left(
\bigcap_{q=m}^{\infty}
\bigcup_{s\ge q}C_s
\right)^c
=\bigcup_{q=m}^{\infty}
\bigcap_{s\ge q}C_s^c.
\]

The first equality uses the representation of $H^{(m)}$, while the second follows from De Morgan's laws.

Consequently, since $\omega_0\in\bigl(H^{(m)}\bigr)^c$, there exists $q_0\ge m$ such that
\[
\omega_0\notin C_s
\qquad\forall s\ge q_0.
\]

By the definitions of union and intersection, membership in the preceding event gives such a threshold $q_0$. This is the meaning of finitely often: no further copy starts beyond that threshold.

On the other hand, since $\omega_0\in H$, we obtain
\[
\omega_0\in C_{s_0}
\qquad\text{for some }s_0\ge q_0.
\]

This follows by taking $q=q_0$ in the membership characterisation of $H$, since $q_0\ge m\ge0$. It contradicts the fact that $\omega_0\notin C_s$ for every $s\ge q_0$.

Hence,
\[
H\cap\bigl(H^{(m)}\bigr)^c=\emptyset,
\]

and therefore
\[
H\subseteq H^{(m)}.
\]

Combining the two inclusions, we conclude that
\[
H=H^{(m)}.
\]
\hfill$\square$


\noindent\studyheading{\studylabel{(iv)}} We first express the two events in terms of the $C_s$:
\[
H_k
=
\bigcup_{0\le s_1<\cdots<s_k}
\bigcap_{j=1}^{k}C_{s_j}
\]
and
\[
H_{m,k}
=
\bigcup_{\substack{0\le s_1<\cdots<s_k\\s_k+N\le m}}
\bigcap_{j=1}^{k}C_{s_j}.
\]
In the first expression, the intersection requires $k$ distinct occurrences of WS, while the union considers every possible selection of their starting positions. In the second, the additional restriction $s_k+N\le m$ requires all selected copies to be completed by time $m$.

Consequently,
\[
\bigcup_{m=1}^{\infty}H_{m,k}
=
\bigcup_{m=1}^{\infty}
\bigcup_{\substack{0\le s_1<\cdots<s_k\\s_k+N\le m}}
\bigcap_{j=1}^{k}C_{s_j} 
=
\bigcup_{0\le s_1<\cdots<s_k}
\bigcap_{j=1}^{k}C_{s_j} 
=H_k.
\]
The first equality substitutes the definition of $H_{m,k}$. The second follows because, for every finite selection of starting positions, we can choose $m\ge s_k+N$, so that selection eventually appears in the union. The last equality uses the representation of $H_k$.

Moreover,
$H_{m,k}\subseteq H_{m+1,k},$
because increasing the deadline from $m$ to $m+1$ does not exclude any previously completed copies.

Therefore,
\[
H_{m,k}\uparrow H_k
\;\text{as }m\to\infty.
\]

\noindent\studyheading{\studylabel{(v)}} Using complements, we obtain
\[
\left(\bigcap_{k=1}^{\infty}H_k\right)^c
=
\bigcup_{k=1}^{\infty}H_k^c =
\bigcup_{q=0}^{\infty}
\bigcap_{s\ge q}C_s^c =
\left(
\bigcap_{q=0}^{\infty}
\bigcup_{s\ge q}C_s
\right)^c =H^c.
\]
The first equality follows from De Morgan's law. The second expresses the equivalence between having fewer than $k$ occurrences for some $k$ and having only finitely many occurrences, which means that from some position $q$ onward no further $C_s$ occurs. The third equality follows again from De Morgan's law, while the last uses the representation of $H$ established above.

Taking complements, we conclude that
\[
H=\bigcap_{k=1}^{\infty}H_k.
\]

Moreover,
$H_{k+1}\subseteq H_k,$
because having at least $k+1$ copies implies having at least $k$ copies. Hence,
\[
H_k\downarrow H
\;\text{as }k\to\infty.
\]

\noindent\studyheading{\studylabel{(ii)}} In particular,
\[
H\subseteq H_k,
\;
H_k\cap H=H.
\]
The inclusion follows because an intersection is contained in each of its constituent sets, while the equality follows from the identity $A\cap B=B$ whenever $B\subseteq A$.

\noindent\studyheading{\studylabel{(iii)}} Recall that
$E_k=C_{(k-1)N}.$
The equality follows because $E_k$ requires one complete copy of WS in the $k$-th block of $N$ keystrokes.

For every $r\in\mathbb N$,
\[
\bigcup_{k\ge r}E_k
=
\bigcup_{k\ge r}C_{(k-1)N}
\subseteq
\bigcup_{s\ge(r-1)N}C_s.
\]
The equality substitutes the definition of $E_k$, while the inclusion follows because the positions $(k-1)N$, for $k\ge r$, form a subset of all starting positions $s\ge(r-1)N$.

Therefore,
\[
\limsup_{k\to\infty}E_k
=
\bigcap_{r=1}^{\infty}
\bigcup_{k\ge r}E_k \subseteq
\bigcap_{r=1}^{\infty}
\bigcup_{s\ge(r-1)N}C_s =
\bigcap_{q=0}^{\infty}
\bigcup_{s\ge q}C_s =H.
\]
The first equality is the definition of the limit superior, and the inclusion follows from the preceding inclusion of unions. The next equality holds because every omitted threshold $q$ is below some retained threshold $(r-1)N$, whose union is smaller and hence imposes a stronger condition. The last equality follows from the representation of $H$.

%
% !TEX root = ../main.tex
\subsection*{HW 41 --- Independence of Products of Random Variables}%
\manualtoc{subsection}{\texorpdfstring{HW 41 --- Independence of Products of Random Variables}{HW 41 --- Independence of Products of Random Variables}}{studiabile-nav-0291}%

\subsubsection*{Statement}%
%


Let $(Y_n)_{n\in\mathbb N\cup\{0\}}$ be a sequence of i.i.d. random variables such that
\[
\mathbb P(Y_n=1)=\mathbb P(Y_n=-1)=\frac12.
\]

Define
\[
X_n:=\prod_{i=0}^{n}Y_i,
\qquad n\in\mathbb N\cup\{0\}.
\]

Prove that $(X_n)_{n\in\mathbb N\cup\{0\}}$ is a sequence of i.i.d. random variables.

\subsubsection*{Expanded Proof of HW 41}%
%


We first notice that
\[
X_0=Y_0,
\qquad
X_n=X_{n-1}Y_n,
\quad n\ge1.
\]

The first equality follows from the definition of $X_0$, while the second follows by separating the last factor in the product defining $X_n$.

Moreover, $X_n\in\{-1,1\}$ for every $n$, since each $Y_i$ takes values in $\{-1,1\}$.

We now use the same idea as in Example 3 on Markov chains: we fix a possible path of the process and calculate its probability.

Fix $n\in\mathbb N\cup\{0\}$ and arbitrary states $x_0,\ldots,x_n\in\{-1,1\}$. We obtain
\[
\begin{aligned}
&\{X_0=x_0,\ldots,X_n=x_n\}\\
&\quad=
\{Y_0=x_0,Y_1=x_0x_1,\ldots,
Y_n=x_{n-1}x_n\}.
\end{aligned}
\]

The equality follows from the recursive relation $X_k=X_{k-1}Y_k$. Indeed, since $X_{k-1}^2=1$, the conditions $X_{k-1}=x_{k-1}$ and $X_k=x_k$ are equivalent to requiring $Y_k=x_{k-1}x_k$, once the previous states have been fixed.

Therefore,
\[
\begin{aligned}
&\mathbb P(X_0=x_0,\ldots,X_n=x_n)\\
&\quad=
\mathbb P(Y_0=x_0,Y_1=x_0x_1,\ldots,
Y_n=x_{n-1}x_n)\\
&\quad=
\mathbb P(Y_0=x_0)
\prod_{k=1}^{n}
\mathbb P(Y_k=x_{k-1}x_k)\\
&\quad=
\left(\frac12\right)^{n+1}.
\end{aligned}
\]

The first equality follows from the identity of events just established, the second from independence of the $Y_k$, and the last from their common distribution, since every $x_{k-1}x_k$ belongs to $\{-1,1\}$.

We now relate this result to Example 3.

Consider the state space $E=\{-1,1\}$, with initial distribution and transition matrix given by
\[
\mu_{-1}=\mu_1=\frac12,
\qquad
P=
\begin{pmatrix}
\frac12&\frac12\\
\frac12&\frac12
\end{pmatrix}.
\]

The joint distribution obtained above can be written as
\[
\mathbb P(X_0=x_0,\ldots,X_n=x_n)
=
\mu_{x_0}\prod_{k=1}^{n}p_{x_{k-1}x_k}.
\]

This is precisely the joint-distribution formula for a Markov chain introduced in Example 3, with $\mu_i=1/2$ and $p_{ij}=1/2$ for every $i,j\in E$.

In this particular case, all rows of the transition matrix are identical: the probability of the next state does not depend on the current state.

We now determine the marginal distributions. For every $m\ge0$ and $a\in\{-1,1\}$,
\[
\begin{aligned}
\mathbb P(X_m=a)
&=
\sum_{x_0,\ldots,x_{m-1}\in E}
\mathbb P(X_0=x_0,\ldots,X_{m-1}=x_{m-1},X_m=a) =2^m\left(\frac12\right)^{m+1} =\frac12.
\end{aligned}
\]

The first equality follows by summing over all possible values of the preceding variables, while the second uses the joint distribution already established and the fact that there are $2^m$ possible choices of $x_0,\ldots,x_{m-1}$.

Consequently,
\[
\begin{aligned}
\mathbb P(X_0=x_0,\ldots,X_n=x_n)
&=\left(\frac12\right)^{n+1} =\prod_{k=0}^{n}\mathbb P(X_k=x_k).
\end{aligned}
\]

The first equality was proved above, and the second follows from the marginal distributions. The factorisation also holds for any finite selection of the $X_n$, by summing over the remaining coordinates. Since the variables take values only in $\{-1,1\}$, this proves their independence according to Definition 4.2.

Moreover, all the $X_n$ have the same distribution. Therefore,
\[
(X_n)_{n\in\mathbb N\cup\{0\}}
\text{ is an i.i.d. sequence.}
\]

Therefore, the statement is proved.
%
\clearpage
%
\section*{5. Integration}
\manualtoc{section}{\texorpdfstring{5. Integration}{5. Integration}}{studiabile-nav-0294}
% !TEX root = ../main.tex
\subsection*{Motivation --- Expectation as an Integral}%
\manualtoc{subsection}{\texorpdfstring{Motivation --- Expectation as an Integral}{Motivation --- Expectation as an Integral}}{studiabile-nav-0295}%


A key quantity in probability theory is the expected value $\mathbb E[X]$ of a random variable $X$. We already know how to calculate it in some familiar cases, but we now want to understand its general mathematical meaning.

If $X$ takes countably many possible values $x_1,x_2,\ldots$, then
\[
\mathbb E[X]
=\sum_{i=1}^{\infty}x_i\,\mathbb P(X=x_i).
\]
This is the usual formula for the expectation of a discrete random variable, whenever the sum is well-defined (for instance, when $X$ is integrable).

If $X$ is real-valued and admits a probability density function $f_X$, then
\[
\mathbb E[X]
=\int_{-\infty}^{\infty} z f_X(z)\,dz.
\]
This is the familiar formula for a random variable with a density, again provided the expectation is well-defined.

These two expressions do not cover every random variable: its distribution need not be discrete or admit a density. We therefore seek a general notion of integration that allows us to write
\[
\mathbb E[X]=\int_{\Omega}X(\omega)\,d\mathbb P(\omega).
\]
Here $\mathbb P$ is the probability measure on $(\Omega,\mathcal F)$, and the integral is understood whenever it is well-defined. The Lebesgue integral will provide the meaning of this expression.

\subsection*{Recall --- Riemann Integral}%
\manualtoc{subsection}{\texorpdfstring{Recall --- Riemann Integral}{Recall --- Riemann Integral}}{studiabile-nav-0296}%


Let $g\in C([a,b])$, and for each $n\in\mathbb N$ consider the dyadic partition
\[
z_i^{(n)}=a+\frac{i}{2^n}(b-a),
\qquad i=0,1,\ldots,2^n.
\]
The Riemann integral can be expressed as
\[
\int_a^b g(z)\,dz
=\lim_{n\to\infty}
\sum_{i=0}^{2^n-1}
 g\bigl(z_i^{(n)}\bigr)
 \bigl(z_{i+1}^{(n)}-z_i^{(n)}\bigr).
\]
The sum is formed by multiplying the value of $g$ at the left endpoint of each subinterval by its length. There are $2^n$ subintervals, so the last summation index is $2^n-1$.

However, a random variable $X:\Omega\to\mathbb R$ is only assumed to be $\mathcal F$-measurable. Measurability alone does not guarantee Riemann integrability, and the sample space $\Omega$ need not be an interval of real numbers. We therefore need an integral defined for measurable functions on general measure spaces. For continuous functions on a real interval, the new integral should agree with the familiar Riemann integral; recall that continuous functions are Borel-measurable.

\subsection*{Notation --- Lebesgue Integral with Respect to a Measure}%
\manualtoc{subsection}{\texorpdfstring{Notation --- Lebesgue Integral with Respect to a Measure}{Notation --- Lebesgue Integral with Respect to a Measure}}{studiabile-nav-0297}%


Let $(S,\Sigma,\mu)$ be a measure space. For a $\Sigma$-measurable function $f$ for which the integral is well-defined (in particular, for any non-negative measurable $f$), we use the notation
\[
\mu(f)
:=\int_S f(s)\,d\mu(s)
=\int_S f\,d\mu.
\]
The measure $\mu$ determines how the integral assigns weight to measurable subsets of $S$. The precise construction of this integral will be developed next.

For any $A\in\Sigma$, we also write
\[
\mu(f;A)
:=\mu(f\mathbf 1_A)
=\int_S f\mathbf 1_A\,d\mu
=\int_A f\,d\mu.
\]
Here $\mathbf 1_A$ is the indicator function of $A$: multiplying by $\mathbf 1_A$ keeps the values of $f$ on $A$ and makes the integrand zero outside $A$. Thus the last expression denotes integration restricted to $A$.

\subsection*{Example --- Counting Measure and Infinite Series}%
\manualtoc{subsection}{\texorpdfstring{Example --- Counting Measure and Infinite Series}{Example --- Counting Measure and Infinite Series}}{studiabile-nav-0298}%


Consider the measure space
\[
(S,\Sigma,\mu)
=\bigl(\mathbb N,\mathcal P(\mathbb N),\#\bigr),
\]
where $\#$ is the counting measure, so that $\mu(\{k\})=1$ for every $k\in\mathbb N$.

Let $f:\mathbb N\to[0,\infty)$ be given by $f(k)=a_k$. Then
\[
\begin{aligned}
\int_S f\,d\mu
&=\sum_{k=1}^{\infty}f(k)\,\mu(\{k\}) =\sum_{k=1}^{\infty}a_k.
\end{aligned}
\]
The first equality is the integral formula for counting measure: each point $k$ contributes its value $f(k)$ multiplied by the measure of the singleton $\{k\}$. The second follows from $\mu(\{k\})=1$ and $f(k)=a_k$. The identity also holds for signed $a_k$ when the series is absolutely convergent.

Thus, the Lebesgue integral must recover the usual infinite series in this special case, just as it must agree with the Riemann integral when both notions apply.
%
% !TEX root = ../main.tex
\subsection*{Integrals of Positive Simple Functions}%
\manualtoc{subsection}{\texorpdfstring{Integrals of Positive Simple Functions}{Integrals of Positive Simple Functions}}{studiabile-nav-0299}%


\subsubsection*{Definition --- Integral of an Indicator Function}%
\manualtoc{subsubsection}{\texorpdfstring{Definition --- Integral of an Indicator Function}{Definition --- Integral of an Indicator Function}}{studiabile-nav-0300}%

Let $(S,\Sigma,\mu)$ be a measure space. For any $A\in\Sigma$, recall that the indicator function of $A$ is
\[
\mathbf 1_A(s)=
\begin{cases}
1,&s\in A,\\
0,&s\notin A.
\end{cases}
\]
We first define its integral by
\[
\mu_0(\mathbf 1_A)
:=\int_S\mathbf 1_A\,d\mu
:=\mu(A).
\]
Thus, integrating an indicator function gives the measure of the set on which it equals $1$. The notation $\mu_0$ emphasizes that, at this stage, the integral has been defined only for indicator functions. The value $\mu(A)$ may be infinite.

\subsubsection*{Definition --- Positive Simple Function}%
\manualtoc{subsubsection}{\texorpdfstring{Definition --- Positive Simple Function}{Definition --- Positive Simple Function}}{studiabile-nav-0301}%

A non-negative function $f:S\to[0,\infty)$ is called a simple function if it can be written as
\[
f(s)=\sum_{k=1}^{m}a_k\mathbf 1_{A_k}(s),
\qquad s\in S,
\]
for some $m\in\mathbb N$, coefficients $a_k\in[0,\infty)$ and sets $A_k\in\Sigma$.

For each $s\in S$, only the terms with $s\in A_k$ contribute to the sum. The sets $A_k$ need not be disjoint, so their contributions may add up. Since there are only finitely many indicators, $f$ can take only finitely many distinct values; this explains the word simple.

\emph{We now verify that every function with this representation belongs to $(m\Sigma)^+$, the class of non-negative $\Sigma$-measurable functions.}

\subsubsection*{HW 42 --- Measurability of a Simple Function}%
\manualtoc{subsubsection}{\texorpdfstring{HW 42 --- Measurability of a Simple Function}{HW 42 --- Measurability of a Simple Function}}{studiabile-nav-0302}%

\studyheading{Statement.} Let
\[
f(s)=\sum_{k=1}^{m}a_k\mathbf 1_{A_k}(s),
\qquad a_k\in[0,\infty),\quad A_k\in\Sigma.
\]
Check that $f\in m\Sigma$.

\par\noindent\studyheading{Expanded Proof.}\par\nobreak
Fix $k\in\{1,\ldots,m\}$. For any Borel set $B\in\mathcal B(\mathbb R)$, we have
\[
(\mathbf 1_{A_k})^{-1}(B)
\in\{\emptyset,A_k,A_k^c,S\}\subseteq\Sigma.
\]
The membership follows because $\mathbf 1_{A_k}$ takes only the values $0$ and $1$: its pre-image of $B$ depends only on which of these two values belong to $B$. The inclusion follows from $A_k\in\Sigma$ and the properties of a $\sigma$-algebra. Hence $\mathbf 1_{A_k}\in m\Sigma$ by Definition 3.1.

We now obtain
\[
a_k\mathbf 1_{A_k}\in m\Sigma,
\qquad k=1,\ldots,m.
\]
This follows from Lemma 3.4\studylabel{(iii)}, since each $a_k$ is a real constant and $\mathbf 1_{A_k}$ is measurable.

Consequently,
\[
f=\sum_{k=1}^{m}a_k\mathbf 1_{A_k}\in m\Sigma.
\]
The membership follows by applying Lemma 3.4\studylabel{(i)} repeatedly to the finite sum of measurable functions. Moreover, $f(s)\ge0$ for all $s\in S$, because every coefficient and indicator is non-negative. Thus $f\in(m\Sigma)^+$.

Therefore, the statement is proved.
%
% !TEX root = ../main.tex
\subsection*{Different Representations and Canonical Form of Simple Functions}%
\manualtoc{subsection}{\texorpdfstring{Different Representations and Canonical Form of Simple Functions}{Different Representations and Canonical Form of Simple Functions}}{studiabile-nav-0305}%


A simple function may admit different representations, corresponding to different choices of coefficients $a_k$ and measurable sets $A_k$.

For example, let $S=[0,1]$ and $\Sigma=\mathcal B([0,1])$. Then
\[
f(s)=\mathbf{1}_{[1/2,1]}(s)
=\mathbf{1}_{[1/2,2/3)}(s)
+\mathbf{1}_{[2/3,1]}(s).
\]

The equality follows because the two intervals on the right are disjoint and their union is $[1/2,1]$. Thus, the same function can be represented using different measurable sets.

A simple function is expressed in canonical form if
\[
f(s)=\sum_{k=1}^{m}a_k\mathbf{1}_{A_k}(s),
\]
where
\[
a_1<a_2<\cdots<a_m,
\qquad
A_k=f^{-1}(\{a_k\}),
\quad k=1,\ldots,m.
\]

In this representation, the coefficients are precisely the distinct values attained by $f$, and the sets $A_k$ are their corresponding level sets. Consequently, these sets are pairwise disjoint and form a partition of $S$.

Every nonnegative simple function can be expressed in canonical form by collecting all points where it takes the same value. The corresponding level sets are measurable, pairwise disjoint, and form a finite partition of $S$. All values actually attained by $f$ must be included, including zero whenever it is attained.

\subsubsection*{Definition --- Integral of a Positive Simple Function}%
\manualtoc{subsubsection}{\texorpdfstring{Definition --- Integral of a Positive Simple Function}{Definition --- Integral of a Positive Simple Function}}{studiabile-nav-0306}%


We denote the collection of non-negative simple functions by $SF^+$. For $f\in SF^+$, represented as
\[
f=\sum_{k=1}^{m}a_k\mathbf{1}_{A_k},
\]
we define
\[
\mu_0(f):=\int_S f\,d\mu
:=\sum_{k=1}^{m}a_k\mu(A_k).
\]

This extends the definition of the integral of an indicator function, since for $f=\mathbf{1}_A$ we recover $\mu_0(f)=\mu(A)$.

\subsubsection*{Observation --- The Convention $0\cdot\infty=0$}%
\manualtoc{subsubsection}{\texorpdfstring{Observation --- The Convention $0\cdot\infty=0$}{Observation --- The Convention 0 x infinity=0}}{studiabile-nav-0307}%


Let $S=\mathbb R$, $\Sigma=\mathcal B(\mathbb R)$, $\mu=\mathrm{Leb}$, and consider the zero function $f=0$.

We obtain
\[
\int_S f\,d\mu
=0\cdot\mu(\mathbb R)
=0\cdot\infty.
\]

Since the integral of the zero function must equal zero, we adopt the convention
\[
0\cdot\infty:=0.
\]

This ensures that zero coefficients contribute nothing to the integral, even when the corresponding sets have infinite measure.

Finally, since a simple function may admit different representations, we must verify that the value of $\mu_0(f)$ does not depend on the chosen representation.

This is the well-definedness of the integral of simple functions, which will be established in the next homework.
%
% !TEX root = ../main.tex
\subsection*{HW 43 --- Well-Definedness of the Integral of Simple Functions}%
\manualtoc{subsection}{\texorpdfstring{HW 43 --- Well-Definedness of the Integral of Simple Functions}{HW 43 --- Well-Definedness of the Integral of Simple Functions}}{studiabile-nav-1001}%

\studyheading{Statement.} Fix $(S,\Sigma,\mu)$ and let $f\in SF^+$ have the following two representations:
\[
f=\sum_{i=1}^{n}a_i\mathbf1_{A_i}
\qquad\text{and}\qquad
f=\sum_{j=1}^{m}b_j\mathbf1_{B_j},
\]
for numbers $(a_i)_{i=1}^{n},(b_j)_{j=1}^{m}\subseteq[0,\infty)$ and sets
$(A_i)_{i=1}^{n},(B_j)_{j=1}^{m}\subseteq\Sigma$. Show that
\[
\mu_0(f)=\sum_{i=1}^{n}a_i\mu(A_i)
=\sum_{j=1}^{m}b_j\mu(B_j).
\]

\par\noindent\studyheading{Proof.}\par\nobreak

\par\noindent\studyheading{Reduction to disjoint partitions.}\par\nobreak

A finite measurable partition of $S$ is a family of nonempty measurable sets that are pairwise disjoint and whose union is $S$. For instance, two sets $A,B\in\Sigma$ give the pieces
\[
A\cap B,\qquad A\cap B^c,\qquad A^c\cap B,\qquad A^c\cap B^c,
\]
with empty pieces omitted. In particular,
\[
A=(A\cap B)\cup(A\cap B^c),\qquad
B=(A\cap B)\cup(A^c\cap B).
\]
Thus each original set can be recovered as a union of disjoint pieces, while the remaining piece completes the partition of $S$.

For the first representation, start with $S$ and consider $A_1,\ldots,A_n$ one at a time. At the $i$th step, split every current piece $E$ into the nonempty sets among
\[
E\cap A_i,\qquad E\cap A_i^c.
\]
Indeed,
\[
E=(E\cap A_i)\cup(E\cap A_i^c),\qquad
(E\cap A_i)\cap(E\cap A_i^c)=\emptyset.
\]
Each step preserves disjointness and the union $S$; measurability follows from closure of $\Sigma$ under complements and finite intersections. After finitely many steps we obtain a partition $C_1,\ldots,C_r$ of $S$. Each $C_k$ lies entirely inside or outside every $A_i$, so each $A_i$ is a union of some $C_k$ and $f$ is constant on each $C_k$.

Applying the same construction to $B_1,\ldots,B_m$ gives another partition $D_1,\ldots,D_t$ of $S$. Each $B_j$ is a union of some $D_\ell$, and $f$ is constant on each $D_\ell$.

Choose $s_k\in C_k$ and $t_\ell\in D_\ell$, and set
\[
c_k=f(s_k),\qquad d_\ell=f(t_\ell),
\qquad k=1,\ldots,r,\quad \ell=1,\ldots,t.
\]
From the two representations of $f$,
\[
c_k=\sum_{i=1}^{n}a_i\mathbf1_{A_i}(s_k),\qquad
d_\ell=\sum_{j=1}^{m}b_j\mathbf1_{B_j}(t_\ell),
\]
for every $k=1,\ldots,r$ and $\ell=1,\ldots,t$.

Since the $C_k$ partition $S$, we obtain
\[
\begin{aligned}
\sum_{i=1}^{n}a_i\mu(A_i)&=\sum_{i=1}^{n}a_i\mu\left(A_i\cap\bigcup_{k=1}^{r}C_k\right) =\sum_{k=1}^{r}\sum_{i=1}^{n}a_i\mu(A_i\cap C_k)\\
&=\sum_{k=1}^{r}\left(\sum_{i=1}^{n}a_i\mathbf1_{A_i}(s_k)\right)\mu(C_k) =\sum_{k=1}^{r}c_k\mu(C_k).
\end{aligned}
\]
The second equality uses finite additivity on the disjoint pieces $A_i\cap C_k$ and rearranges finite nonnegative sums. The third follows because each $C_k$ lies entirely inside or outside $A_i$, and the last uses the expression for $c_k$. Likewise,
\[
\sum_{j=1}^{m}b_j\mu(B_j)=\sum_{\ell=1}^{t}d_\ell\mu(D_\ell).
\]
These identities show that we may assume, without loss of generality, that the representing families are finite measurable partitions of $S$, since the replacement preserves both $f$ and the weighted sums. Moreover, by merging partition elements with equal coefficients, any such representation can be reduced to canonical form, again without changing the weighted sum.

\par\noindent\studyheading{Proof for partition representations.}\par\nobreak

Assume now that $(A_i)_{i=1}^{n}$ and $(B_j)_{j=1}^{m}$ are finite measurable partitions of $S$, with the corresponding coefficients $a_i$ and $b_j$. As proved above, this entails no loss of generality.

The nonempty intersections $A_i\cap B_j$ form a common finite measurable partition of $S$. Indeed, they cover $S$ because every $s\in S$ belongs to some $A_i$ and some $B_j$. Moreover, for $(i,j)\neq(k,\ell)$,
\[
\begin{aligned}
(A_i\cap B_j)\cap(A_k\cap B_\ell)&=(A_i\cap A_k)\cap(B_j\cap B_\ell) =\emptyset.
\end{aligned}
\]
The last equality follows because either $i\neq k$ or $j\neq\ell$, so at least one intersection is empty. Finally, $A_i\cap B_j\in\Sigma$ by closure under finite intersections.

For every $i=1,\ldots,n$ and $j=1,\ldots,m$, we distinguish two cases.

\par\noindent\studyheading{Case 1 --- Empty intersection.}\par\nobreak
 If $A_i\cap B_j=\emptyset$, then
$\mu(A_i\cap B_j)=0$.

\par\noindent\studyheading{Case 2 --- Nonempty intersection.}\par\nobreak
 If $A_i\cap B_j\neq\emptyset$,
choose $s\in A_i\cap B_j$. Since the two families are partitions, $s$
belongs to exactly one set in each family. Hence,
\[
\begin{aligned}
a_i&=\sum_{k=1}^{n}a_k\mathbf1_{A_k}(s) =f(s) =\sum_{\ell=1}^{m}b_\ell\mathbf1_{B_\ell}(s) =b_j.
\end{aligned}
\]
The first and last equalities follow because only the indicators of $A_i$
and $B_j$ equal $1$ in their respective sums; the middle equalities use the
two representations of $f$.

Thus $a_i=b_j$ holds for pairs whose intersection is nonempty. For the
remaining pairs the coefficients need not agree, but the intersection has
measure zero. In both cases, therefore,
\[
a_i\mu(A_i\cap B_j)=b_j\mu(A_i\cap B_j),
\qquad\forall\,i=1,\ldots,n,\quad j=1,\ldots,m.
\]
Using $S=\bigcup_{j=1}^{m}B_j=\bigcup_{i=1}^{n}A_i$, we obtain
\[
\begin{aligned}
\sum_{i=1}^{n}a_i\mu(A_i)&=\sum_{i=1}^{n}a_i\mu(A_i\cap S) =\sum_{i=1}^{n}a_i\mu\left(A_i\cap\bigcup_{j=1}^{m}B_j\right) =\sum_{i=1}^{n}\sum_{j=1}^{m}a_i\mu(A_i\cap B_j)\\
&=\sum_{j=1}^{m}\sum_{i=1}^{n}b_j\mu(A_i\cap B_j) =\sum_{j=1}^{m}b_j\mu\left(B_j\cap\bigcup_{i=1}^{n}A_i\right) =\sum_{j=1}^{m}b_j\mu(B_j).
\end{aligned}
\]
The third and fifth equalities use finite additivity on disjoint intersections. The fourth follows from the equality of the contributions in both cases and rearrangement of finite sums. All expressions remain valid when some measures are infinite, under the convention $0\cdot\infty=0$.

Thus the weighted sum is independent of the chosen representation of $f$, and $\mu_0(f)$ is well-defined. Therefore, the statement is proved.
%
% !TEX root = ../main.tex
\subsection*{HW 44 --- Properties of Integrals of Simple Functions}%
\manualtoc{subsection}{\texorpdfstring{HW 44 --- Properties of Integrals of Simple Functions}{HW 44 --- Properties of Integrals of Simple Functions}}{studiabile-nav-1002}%

\studyheading{Statement.} Fix a measure space $(S,\Sigma,\mu)$. Prove the following properties.

\studyheading{(i)} If $f,g\in SF^+$ and $\mu(\{s\in S:f(s)\ne g(s)\})=0$, then $\mu_0(f)=\mu_0(g)$.

\studyheading{(ii)} If $f,g\in SF^+$ and $c\ge0$, then $f+g,cf\in SF^+$; moreover,
\[
\mu_0(f+g)=\mu_0(f)+\mu_0(g),\qquad \mu_0(cf)=c\mu_0(f).
\]

\studyheading{(iii)} If $f,g\in SF^+$ and $f\le g$, then $\mu_0(f)\le\mu_0(g)$.

\studyheading{(iv)} If $f,g\in SF^+$, then $f\vee g:=\max\{f,g\}$ and $f\wedge g:=\min\{f,g\}$ also belong to $SF^+$.

\par\noindent\studyheading{Recall --- $\boldsymbol{\mu}$-almost everywhere.}%
\inlinetoc{subsubsection}{\texorpdfstring{Recall --- $\boldsymbol{\mu}$-almost everywhere.}{Recall --- mu-almost everywhere.}}{studiabile-nav-1003}%

A property $P(s)$ holds $\mu$-almost everywhere ($\mu$-a.e.) if it fails only on a set of $\mu$-measure zero. In particular, for measurable $f,g$,
\[
f=g\quad\mu\text{-a.e.}
\quad\Longleftrightarrow\quad
\mu\bigl(\{s\in S:f(s)\ne g(s)\}\bigr)=0.
\]
Thus $f$ and $g$ need not agree at every point.

When constructing a common partition as in HW~43, we discard only empty intersections, since partition elements must be nonempty. We do not discard nonempty sets of measure zero: for example, on $S=[0,1]$ with $\mu=\mathrm{Leb}$,
\[
D=\{\tfrac12\}\ne\emptyset,\qquad \mu(D)=0.
\]
Such a set remains a partition element, although its contribution to the integral is zero.

\par\noindent\studyheading{Proof.}\par\nobreak

By HW~43, both $f$ and $g$ admit canonical representations, whose level sets form finite measurable partitions of $S$. Intersecting these partitions and omitting empty intersections gives a common finite measurable partition $D_1,\ldots,D_r$, on which both functions are constant.
\[
f=\sum_{k=1}^{r}a_k\mathbf1_{D_k},\qquad
 g=\sum_{k=1}^{r}b_k\mathbf1_{D_k},
\]
where $a_k,b_k\ge0$. By HW~43, the integral is independent of the representation, so
\[
\mu_0(f)=\sum_{k=1}^{r}a_k\mu(D_k),\qquad
\mu_0(g)=\sum_{k=1}^{r}b_k\mu(D_k).
\]

\studyheading{(i) Equality almost everywhere.}
Suppose $f=g$ $\mu$-a.e., and let
\[
N:=\{s\in S:f(s)\ne g(s)\},\qquad \mu(N)=0.
\]
Fix $k$ such that $a_k\ne b_k$. For any $s\in D_k$, the partition representations give
\[
f(s)=a_k\ne b_k=g(s)\quad\Longrightarrow\quad s\in N.
\]
Since $s\in D_k$ was arbitrary, $D_k\subseteq N$. By monotonicity of $\mu$,
\[
0\le\mu(D_k)\le\mu(N)=0
\quad\Longrightarrow\quad \mu(D_k)=0.
\]
Consequently, for every $k=1,\ldots,r$,
\[
a_k\mu(D_k)=b_k\mu(D_k).
\]
Indeed, if $a_k=b_k$ the equality is immediate; otherwise $\mu(D_k)=0$ and both sides vanish. Summing gives
\[
\mu_0(f)=\sum_{k=1}^{r}a_k\mu(D_k)
=\sum_{k=1}^{r}b_k\mu(D_k)=\mu_0(g).
\]

\studyheading{(ii) Additivity and positive homogeneity.}
Using the common partition,
\[
f+g=\sum_{k=1}^{r}(a_k+b_k)\mathbf1_{D_k},\qquad
cf=\sum_{k=1}^{r}(ca_k)\mathbf1_{D_k}.
\]
Thus $f+g,cf\in SF^+$, since their coefficients are nonnegative and the $D_k$ are measurable. Moreover,
\[
\begin{aligned}
\mu_0(f+g)&=\sum_{k=1}^{r}(a_k+b_k)\mu(D_k) =\sum_{k=1}^{r}a_k\mu(D_k)
 +\sum_{k=1}^{r}b_k\mu(D_k) =\mu_0(f)+\mu_0(g),
\end{aligned}
\]
and
\[
\begin{aligned}
\mu_0(cf)&=\sum_{k=1}^{r}(ca_k)\mu(D_k) =c\sum_{k=1}^{r}a_k\mu(D_k)
=c\mu_0(f).
\end{aligned}
\]
The equalities follow from the definition of $\mu_0$ and distributivity of finite nonnegative sums, with $0\cdot\infty=0$ when necessary.

\studyheading{(iii) Monotonicity.}
Suppose $f\le g$. Since $D_k\ne\emptyset$ and both functions are constant on it, $a_k\le b_k$ for every $k=1,\ldots,r$. Hence
\[
\mu_0(f)=\sum_{k=1}^{r}a_k\mu(D_k)
\le\sum_{k=1}^{r}b_k\mu(D_k)=\mu_0(g).
\]
The inequality follows from $a_k\le b_k$ and $\mu(D_k)\ge0$ for every $k$.

\studyheading{(iv) Maximum and minimum.}
Since $f=a_k$ and $g=b_k$ on $D_k$,
\[
\begin{aligned}
f\vee g&=\sum_{k=1}^{r}\max\{a_k,b_k\}\mathbf1_{D_k},\\
f\wedge g&=\sum_{k=1}^{r}\min\{a_k,b_k\}\mathbf1_{D_k}.
\end{aligned}
\]
Both functions belong to $SF^+$, because their coefficients are nonnegative and their partition elements are measurable.

Therefore, the statement is proved.
%
% !TEX root = ../main.tex
\subsection*{Example --- Equality Lebesgue-Almost Everywhere}%
\manualtoc{subsection}{\texorpdfstring{Example --- Equality Lebesgue-Almost Everywhere}{Example --- Equality Lebesgue-Almost Everywhere}}{studiabile-nav-1004}%

Let $S=[0,1]$, $\Sigma=\mathcal B([0,1])$, $\mu=\mathrm{Leb}$, and consider
\[
f=\mathbf1_{[\frac12,1]},\qquad
 g=\mathbf1_{(\frac12,1]}.
\]
The functions differ only at $s=\frac12$, so
\[
\{s\in S:f(s)\ne g(s)\}=\bigl\{\tfrac12\bigr\},
\qquad \mathrm{Leb}\bigl(\{\tfrac12\}\bigr)=0.
\]
Hence $f=g$ $\mathrm{Leb}$-a.e. and, by HW~44(i),
\[
\mu_0(f)
=\mathrm{Leb}\bigl([\tfrac12,1]\bigr)
=\tfrac12
=\mathrm{Leb}\bigl((\tfrac12,1]\bigr)
=\mu_0(g).
\]
The equalities follow from the definition of the integral of an indicator function and the length of the two intervals.

\subsection*{Definition --- Integral of a Positive Measurable Function}%
\manualtoc{subsection}{\texorpdfstring{Definition --- Integral of a Positive Measurable Function}{Definition --- Integral of a Positive Measurable Function}}{studiabile-nav-1005}%

Let $(S,\Sigma,\mu)$ be a measure space. For any $f\in(m\Sigma)^+$, define
\[
\mu(f):=\int_S f\,d\mu
:=\sup\bigl\{\mu_0(h):h\in SF^+,\ h\le f\bigr\}.
\]
Thus the integral of $f$ is the supremum of the integrals of all nonnegative simple functions lying below $f$.

\subsection*{Observation --- Agreement with the Simple Integral}%
\manualtoc{subsection}{\texorpdfstring{Observation --- Agreement with the Simple Integral}{Observation --- Agreement with the Simple Integral}}{studiabile-nav-1006}%

If $f\in SF^+$, then
\[
\mu(f)=\mu_0(f).
\]
Indeed, $f$ itself belongs to the family over which the supremum is taken, so $\mu(f)\ge\mu_0(f)$. Conversely, for every $h\in SF^+$ with $h\le f$, monotonicity of $\mu_0$ (HW~44(iii)) gives $\mu_0(h)\le\mu_0(f)$. Taking the supremum yields $\mu(f)\le\mu_0(f)$, proving the equality.

\subsection*{HW 45 --- Monotonicity of the Positive Lebesgue Integral}%
\manualtoc{subsection}{\texorpdfstring{HW 45 --- Monotonicity of the Positive Lebesgue Integral}{HW 45 --- Monotonicity of the Positive Lebesgue Integral}}{studiabile-nav-1007}%

\studyheading{Statement.} Prove that if $f,g\in(m\Sigma)^+$ satisfy $f\le g$, then
\[
\mu(f)\le\mu(g).
\]
%
% !TEX root = ../main.tex
\par\noindent\studyheading{Proof.}\par\nobreak

By definition, for every $f\in(m\Sigma)^+$,
\[
\mu(f):=\sup\{\mu_0(h):h\in SF^+,\ h\le f\}.
\]
Let $f,g\in(m\Sigma)^+$ with $f\le g$. For any $h\in SF^+$ such that $h\le f$, we have
\[
h\le f\le g\quad\Longrightarrow\quad h\le g.
\]
Hence,
\[
\{h\in SF^+:h\le f\}\subseteq\{h\in SF^+:h\le g\}.
\]
The inclusion follows because every nonnegative simple function lying below $f$ also lies below $g$.

Taking the supremum over the two families, we obtain
\[
\begin{aligned}
\mu(f)&=\sup\{\mu_0(h):h\in SF^+,\ h\le f\} \le\sup\{\mu_0(h):h\in SF^+,\ h\le g\} =\mu(g).
\end{aligned}
\]
The inequality follows from the inclusion above, since the supremum over a smaller family cannot exceed the supremum over a larger one. Both families are nonempty because the zero function is an admissible simple function.

Therefore, the statement is proved.
%
% !TEX root = ../main.tex
\subsection*{Lemma 5.1 --- Zero Integral Implies Zero Almost Everywhere}%
\manualtoc{subsection}{\texorpdfstring{Lemma 5.1 --- Zero Integral Implies Zero Almost Everywhere}{Lemma 5.1 --- Zero Integral Implies Zero Almost Everywhere}}{studiabile-nav-1008}%

\studyheading{Statement.} Let $(S,\Sigma,\mu)$ be a measure space and let
$f\in(m\Sigma)^+$ satisfy $\mu(f)=0$. Then
\[
\mu(\{f>0\})=\mu(\{s\in S:f(s)>0\})=0.
\]

\par\noindent\studyheading{Proof.}\par\nobreak

We argue by contradiction. Suppose that $\mu(\{f>0\})>0$ and define
\[
A_n:=\left\{f>\frac1n\right\}\in\Sigma,\qquad n\in\mathbb N.
\]
The sets are measurable because $f$ is $\Sigma$-measurable. As we will see in
HW~46,
\[
\{f>0\}
=\bigcup_{n=1}^{\infty}\left\{f>\frac1n\right\}
=\lim_{n\to\infty}\left\{f>\frac1n\right\}.
\]
Moreover, $A_n\subseteq A_{n+1}$ because $1/(n+1)<1/n$. By continuity
from below (Lemma~1.5),
\[
0<\mu(\{f>0\})=\lim_{n\to\infty}\mu(A_n).
\]
Since $(\mu(A_n))_{n\in\mathbb N}$ is nondecreasing and its limit is
strictly positive, there exists $\bar n\in\mathbb N$ such that
$\mu(A_n)>0$ for every $n\ge\bar n$. Fix $n=\bar n$ and consider the
nonnegative simple function
\[
h:=\frac1{\bar n}\mathbf1_{A_{\bar n}}\in SF^+.
\]
We have $h\le f$: if $s\in A_{\bar n}$, then
$h(s)=1/\bar n<f(s)$; otherwise $h(s)=0\le f(s)$.
Therefore, by the definition of the positive Lebesgue integral,
\[
\begin{aligned}
0=\mu(f)&=\sup\{\mu_0(g):g\in SF^+,\ g\le f\} \ge\mu_0(h)
=\frac1{\bar n}\mu(A_{\bar n})>0.
\end{aligned}
\]
The inequality holds because $h$ is an admissible simple function; the
following equality uses the definition of $\mu_0$, and the last inequality
follows from $\mu(A_{\bar n})>0$. This is a contradiction. Thus
$\mu(\{f>0\})=0$, i.e., $f=0$ $\mu$-a.e.

Therefore, the statement is proved.
%
% !TEX root = ../main.tex
\subsection*{HW 46 --- Level-Set Representation of $\{f>0\}$}%
\manualtoc{subsection}{\texorpdfstring{HW 46 --- Level-Set Representation of $\{f>0\}$}{HW 46 --- Level-Set Representation of f>0}}{studiabile-nav-1009}%

\studyheading{Statement.} Let $f\in(m\Sigma)^+$. Prove that
\[
\{f>0\}
=\bigcup_{n=1}^{\infty}\left\{f>\frac1n\right\}
=\lim_{n\to\infty}\left\{f>\frac1n\right\}.
\]

\par\noindent\studyheading{Proof.}\par\nobreak

{\itshape We first prove that
\[
\{f>0\}=\bigcup_{n=1}^{\infty}\left\{f>\frac1n\right\}
\]
by showing both inclusions.}

\studyheading{($\subseteq$)} Let $s_0\in\{f>0\}$. Then $f(s_0)>0$, so there exists
$n_0\in\mathbb N$ such that
\[
\frac1{n_0}<f(s_0).
\]
This follows from the Archimedean property of $\mathbb R$, and is immediate
if $f(s_0)=\infty$. Therefore,
\[
s_0\in\left\{f>\frac1{n_0}\right\}
\subseteq\bigcup_{n=1}^{\infty}\left\{f>\frac1n\right\}.
\]
Since $s_0$ was arbitrary, the first inclusion follows.

\studyheading{($\supseteq$)} Conversely, let
\[
s_0\in\bigcup_{n=1}^{\infty}\left\{f>\frac1n\right\}.
\]
By definition of union, there exists $n_0\in\mathbb N$ such that
\[
f(s_0)>\frac1{n_0}>0.
\]
Hence $s_0\in\{f>0\}$, proving the reverse inclusion.

Finally, define
\[
A_n:=\left\{f>\frac1n\right\},\qquad n\in\mathbb N.
\]
Since
\[
\frac1{n+1}<\frac1n
\quad\Longrightarrow\quad A_n\subseteq A_{n+1},
\]
the sequence $(A_n)_{n\in\mathbb N}$ is increasing. Therefore, by the
definition of the limit of an increasing sequence of sets,
\[
\lim_{n\to\infty}A_n
=\bigcup_{n=1}^{\infty}A_n
=\{f>0\}.
\]
Combining the two identities, we obtain
\[
\{f>0\}
=\bigcup_{n=1}^{\infty}\left\{f>\frac1n\right\}
=\lim_{n\to\infty}\left\{f>\frac1n\right\}.
\]
Therefore, the statement is proved.
%
% !TEX root = ../main.tex
\subsection*{Example --- Application of Lemma 5.1}%
\manualtoc{subsection}{\texorpdfstring{Example --- Application of Lemma 5.1}{Example --- Application of Lemma 5.1}}{studiabile-nav-1010}%

Let $(S,\Sigma,\mu)=([0,1],\mathcal B([0,1]),\mathrm{Leb})$ and
\[
 f(s)=\mathbf 1_{\{\frac12\}}(s).
\]
Since $f\in SF^+$,
\[
 \mu(f)=\mu_0(f)=\mathrm{Leb}\!\left(\left\{\frac12\right\}\right)=0.
\]
The equalities follow from the agreement of $\mu$ with $\mu_0$ on simple functions, the integral of an indicator, and the fact that a singleton has zero Lebesgue measure. By Lemma~5.1,
\[
 \mu(\{f>0\})=\mathrm{Leb}\!\left(\left\{\frac12\right\}\right)=0.
\]
Thus $f=0$, $\mathrm{Leb}$-a.e., although $f(\frac12)=1$.

The next theorem is probably the most useful single tool in the whole course.

\subsection*{Theorem 5.2 --- Monotone Convergence Theorem (MON)}%
\manualtoc{subsection}{\texorpdfstring{Theorem 5.2 --- Monotone Convergence Theorem (MON)}{Theorem 5.2 --- Monotone Convergence Theorem (MON)}}{studiabile-nav-1011}%

\studyheading{Statement.} Let $(f_n)_{n\in\mathbb N}\subseteq(m\Sigma)^+$ and let $f\in(m\Sigma)^+$ satisfy $f_n(s)\uparrow f(s)$ for $\mu$-almost every $s\in S$. Then
\[
 \int_S f_n\,d\mu\;\uparrow\;\int_S f\,d\mu.
\]

\par\noindent\studyheading{Expanded Proof.}\par\nobreak
\par
Let $N\in\Sigma$ be a null set outside which $f_n(s)\uparrow f(s)$. We first justify the comparison of integrals when an inequality holds only outside $N$, using no results introduced later. If $u,v\in(m\Sigma)^+$ satisfy $u\le v$ on $S\setminus N$, take any $h\in SF^+$ with $h\le u$. Then
\[
 h\mathbf1_{S\setminus N}\le v,
 \qquad \mu_0(h)=\mu_0(h\mathbf1_{S\setminus N}).
\]
The first inequality holds because $h\le u\le v$ outside $N$ and the left-hand side is zero on $N$; the equality follows from HW~44(i), since the two simple functions differ only on a null set. Hence $\mu_0(h)\le\mu(v)$, and taking the supremum over $h\le u$ gives $\mu(u)\le\mu(v)$.

In particular, since $f_n\le f_{n+1}\le f$ outside $N$, we obtain
\[
 \mu(f_n)\le\mu(f_{n+1})\le\mu(f),\qquad n\in\mathbb N.
\]

Therefore $(\mu(f_n))$ is nondecreasing and has a limit
\[
 \alpha:=\lim_{n\to\infty}\mu(f_n)\le\mu(f).
\]
\textit{We now prove the reverse inequality.} Fix $c\in(0,1)$ and $h\in SF^+$ such that $h\le f$, and define
\[
 E_n:=\{s\in S:ch(s)\le f_n(s)\}\in\Sigma.
\]
Measurability follows from that of $h$ and $f_n$. Outside $N$,
$E_n\cap N^c\subseteq E_{n+1}\cap N^c$ by monotonicity of $(f_n)$.
These sets cover $N^c$: if $h(s)=0$, then $s\in E_n$ for every $n$;
if $h(s)>0$, then $f(s)\ge h(s)>ch(s)$, so $s\in E_n$ eventually
because $f_n(s)\uparrow f(s)$. The latter also holds if $f(s)=\infty$,
since $h(s)$ is finite. This is the argument proved in HW~47.

Write the simple function as
\[
 h=\sum_{k=1}^{m}a_k\mathbf1_{A_k},\qquad a_k\ge0,\quad A_k\in\Sigma.
\]
Since $f_n\ge ch$ on $E_n$ and $f_n\ge0$ elsewhere, we have
\[
 f_n\ge ch\mathbf1_{E_n}.
\]
By HW~45, agreement with the simple integral, and HW~44(ii),
\[
\begin{aligned}
\mu(f_n)
 &\ge\mu(ch\mathbf1_{E_n}) =\mu_0(ch\mathbf1_{E_n}) =c\sum_{k=1}^{m}a_k\mu(A_k\cap E_n).
\end{aligned}
\]
The last equality follows from the representation $h\mathbf1_{E_n}=\sum_{k=1}^{m}a_k\mathbf1_{A_k\cap E_n}$.

To apply continuity from below, define
\[
 \widetilde E_n:=(E_n\cap N^c)\cup N.
\]
The preceding properties imply $\widetilde E_n\uparrow S$. Moreover,
for every $k=1,\ldots,m$ and $n\in\mathbb N$,
\[
 \mu(A_k\cap E_n)=\mu(A_k\cap\widetilde E_n),
\]
since the two sets agree outside $N$ and their parts inside $N$
have measure zero. Thus, by continuity from below (Lemma~1.5)
applied to the increasing sets $A_k\cap\widetilde E_n$,
\[
 \mu(A_k\cap E_n)=\mu(A_k\cap\widetilde E_n)
 \uparrow\mu(A_k),\qquad k=1,\ldots,m.
\]
No measures are subtracted, so this remains valid when
$\mu(A_k)=\infty$; zero coefficients are handled by $0\cdot\infty=0$.

Letting $n\to\infty$ and using the finiteness of the sum, we obtain
\[
 \alpha\ge c\sum_{k=1}^{m}a_k\mu(A_k)=c\mu_0(h).
\]
Since this holds for every $h\in SF^+$ with $h\le f$, the definition of $\mu(f)$ yields
\[
 \alpha\ge c\sup\{\mu_0(h):h\in SF^+,\ h\le f\}=c\mu(f).
\]
Finally, letting $c\uparrow1$ gives $\alpha\ge\mu(f)$, including the case $\mu(f)=\infty$. Combined with $\alpha\le\mu(f)$, this proves
\[
 \lim_{n\to\infty}\mu(f_n)=\mu(f),
\]
and the convergence is increasing because $(\mu(f_n))$ is nondecreasing.

Therefore, the statement is proved.
%
% !TEX root = ../main.tex
\subsection*{HW 47 --- Increasing Sets in the Monotone Convergence Theorem}%
\manualtoc{subsection}{\texorpdfstring{HW 47 --- Increasing Sets in the Monotone Convergence Theorem}{HW 47 --- Increasing Sets in the Monotone Convergence Theorem}}{studiabile-nav-1014}%

\studyheading{Statement.} In the proof of Theorem 5.2, let $f_n\uparrow f$ $\mu$-a.e., fix $c\in(0,1)$, and choose $h\in SF^+$ such that $h\le f$. Set
\[
E_n:=\{s\in S:ch(s)\le f_n(s)\}\in\Sigma.
\]
Prove that $E_n\subseteq E_{n+1}$ with $E_n\uparrow S$ (when the convergence holds everywhere).

\par\noindent\studyheading{Proof.}\par\nobreak
\par
We first assume that $f_n(s)\uparrow f(s)$ for every $s\in S$.

Fix $n\in\mathbb N$ and let $s_0\in E_n$. Then
\[
ch(s_0)\le f_n(s_0)\le f_{n+1}(s_0),
\]
where the first inequality follows from the definition of $E_n$, and the second from the monotonicity of $(f_n)$. Thus $s_0\in E_{n+1}$, and therefore
\[
E_n\subseteq E_{n+1},\qquad n\in\mathbb N.
\]

\textit{We now show that $\bigcup_{n=1}^{\infty}E_n=S$.} Let $s_0\in S$. If $h(s_0)=0$, then $ch(s_0)=0\le f_n(s_0)$ for every $n$, since $f_n\ge0$.

If $h(s_0)>0$, then
\[
f(s_0)\ge h(s_0)>ch(s_0),
\]
where the first inequality uses $h\le f$, and the second uses $0<c<1$. Since $f_n(s_0)\uparrow f(s_0)$, there exists $n_0\in\mathbb N$ such that
\[
f_n(s_0)\ge ch(s_0),\qquad \forall n\ge n_0.
\]
This remains true when $f(s_0)=\infty$, because $ch(s_0)$ is finite and $f_n(s_0)\to\infty$.
In either case $s_0\in E_{n_0}$ for some $n_0$, so $S\subseteq\bigcup_{n=1}^{\infty}E_n$. The reverse inclusion is immediate because $E_n\subseteq S$. Hence
\[
\bigcup_{n=1}^{\infty}E_n=S.
\]
Combining this equality with $E_n\subseteq E_{n+1}$ gives
\[
E_n\uparrow S.
\]

\par\noindent\studyheading{Almost-everywhere qualification.}\par\nobreak
 Theorem 5.2 only assumes $f_n\uparrow f$ $\mu$-a.e. Let $N\in\Sigma$ be a null set outside which this convergence holds. The argument above then proves
\[
E_n\cap N^c\subseteq E_{n+1}\cap N^c,
\qquad
\bigcup_{n=1}^{\infty}(E_n\cap N^c)=N^c.
\]
In particular, the measurable sets
\[
\widetilde E_n:=(E_n\cap N^c)\cup N
\]
satisfy $\widetilde E_n\uparrow S$ and agree with $E_n$ outside $N$.
For every $A\in\Sigma$, the portions of $A\cap E_n$ and
$A\cap\widetilde E_n$ contained in $N$ have measure zero.
By finite additivity and continuity from below (Lemma~1.5),
\[
\mu(A\cap E_n)=\mu(A\cap\widetilde E_n)
\uparrow\mu(A).
\]
No subtraction of measures is involved, even when $\mu(A)=\infty$.
This is the precise statement needed in the proof of Theorem~5.2.

\par\noindent\studyheading{Remark.}\par\nobreak
 Without this qualification, $E_n\uparrow S$ need not hold: on $S=[0,1]$ with Lebesgue measure, take $f=h=1$ and $f_n=\mathbf1_{(0,1]}$ for every $n$. Then $f_n\uparrow f$ Lebesgue-a.e., but $E_n=(0,1]$ for every $n$.

Therefore, the statement is proved.
%
% !TEX root = ../main.tex
\subsection*{Approximation by Positive Simple Functions}%
\manualtoc{subsection}{\texorpdfstring{Approximation by Positive Simple Functions}{Approximation by Positive Simple Functions}}{studiabile-nav-1012}%

From the monotone convergence theorem (Theorem 5.2), the next question arises: given
$f\in(m\Sigma)^+$, how can we construct an increasing sequence $(f^{(r)})_{r\in\mathbb N}\subseteq\SF$
that approximates $f$?

For $k,r\in\mathbb N$, define
\[
A_k^r:=\left\{s\in S:\frac{k-1}{2^r}\le f(s)<\frac{k}{2^r}\right\}.
\]
Set
\begin{equation*}
\label{eq:simple_approx}
f^{(r)}(s)
=\sum_{k=1}^{r2^r}\frac{k-1}{2^r}\ind{A_k^r}(s)
  +r\ind{\{f\ge r\}}(s).
\tag{5.1}
\end{equation*}

\par\noindent\studyheading{Verification.}\par\nobreak

First, $A_k^r,\{f\ge r\}\in\Sigma$ by measurability of $f$. The sets
$A_1^r,\ldots,A_{r2^r}^r,\{f\ge r\}$ are pairwise disjoint and cover $S$: if
$0\le f(s)<r$, its value belongs to exactly one half-open interval
$[(k-1)/2^r,k/2^r)$; if $f(s)\ge r$, it belongs to the last set.
Hence $f^{(r)}\in\SF$, since (5.1) is a finite sum of measurable indicators with
nonnegative coefficients.

\textit{We now prove that $f^{(r)}\le f^{(r+1)}\le f$.} For any $s\in S$ with
$x:=f(s)<\infty$, formula (5.1) gives
\[
f^{(r)}(s)=\min\left\{r,\frac{\lfloor 2^r x\rfloor}{2^r}\right\}.
\]
Indeed, below $r$ the coefficient is the left endpoint of the interval
containing $x$, while at and above $r$ the value is truncated at $r$.
For every $x\ge0$,
\[
\frac{\lfloor 2^r x\rfloor}{2^r}
\le\frac{\lfloor 2^{r+1}x\rfloor}{2^{r+1}}
\le x.
\]
The first inequality follows from
$\lfloor2^{r+1}x\rfloor\ge2\lfloor2^r x\rfloor$;
the second is the defining property of the integer part.
Since the truncation level also increases from $r$ to $r+1$, we obtain
$f^{(r)}(s)\le f^{(r+1)}(s)\le f(s)$.
If $f(s)=\infty$, then $f^{(r)}(s)=r$, so the same inequalities hold.

Finally, for $f(s)<\infty$, choose $r>f(s)$. The truncation no longer applies and
\[
0\le f(s)-f^{(r)}(s)<\frac1{2^r},
\]
by the defining bounds of the interval $A_k^r$ containing $s$.
Thus $f^{(r)}(s)\to f(s)$. If $f(s)=\infty$, then
$f^{(r)}(s)=r\to\infty$. Consequently,
\[
f^{(r)}(s)\uparrow f(s),\qquad \forall s\in S.
\]
In particular, the monotone convergence theorem applies to this sequence.

\subsection*{Lemma 5.3 --- Equality Almost Everywhere}%
\manualtoc{subsection}{\texorpdfstring{Lemma 5.3 --- Equality Almost Everywhere}{Lemma 5.3 --- Equality Almost Everywhere}}{studiabile-nav-1013}%

\studyheading{Statement.} If $f,g\in(m\Sigma)^+$ and
$\mu(\{f\ne g\})=0$, then $\mu(f)=\mu(g)$.

\par\noindent\studyheading{Proof.}\par\nobreak

Let $(f^{(r)})_{r\in\mathbb N}$ and $(g^{(r)})_{r\in\mathbb N}$ be the
simple-function approximations given by (5.1). By the construction above,
\[
f^{(r)}\uparrow f,\qquad g^{(r)}\uparrow g.
\]

Set $N:=\{s\in S:f(s)\ne g(s)\}\in\Sigma$, so that $\mu(N)=0$.
For every $r\in\mathbb N$ and $s\notin N$,
\[
f(s)=g(s)\quad\Longrightarrow\quad f^{(r)}(s)=g^{(r)}(s),
\]
since (5.1) applies the same dyadic approximation to both values.
Therefore,
\[
\{f^{(r)}\ne g^{(r)}\}\subseteq N
\quad\Longrightarrow\quad
0\le\mu(\{f^{(r)}\ne g^{(r)}\})\le\mu(N)=0.
\]
The second implication follows from monotonicity of $\mu$. Thus
$f^{(r)}=g^{(r)}$ $\mu$-a.e., and HW 44(i) gives
\[
\mu(f^{(r)})=\mu_0(f^{(r)})
=\mu_0(g^{(r)})=\mu(g^{(r)}),\qquad \forall r\in\mathbb N.
\]
The first and last equalities use the agreement of the positive Lebesgue
integral with the simple integral.

Applying the monotone convergence theorem (Theorem 5.2), we conclude that
\[
\begin{aligned}
\mu(f)&=\lim_{r\to\infty}\mu(f^{(r)}) =\lim_{r\to\infty}\mu(g^{(r)}) =\mu(g).
\end{aligned}
\]
The first and last equalities follow from monotone convergence, and the middle
one from equality of the integrals for every $r$.

Therefore, the statement is proved.
%
% !TEX root = ../main.tex
\subsection*{Lemma 5.4 --- Fatou's Lemma for Functions}%
\manualtoc{subsection}{\texorpdfstring{Lemma 5.4 --- Fatou's Lemma for Functions}{Lemma 5.4 --- Fatou's Lemma for Functions}}{studiabile-nav-1015}%

\studyheading{Statement.} Let $(S,\Sigma,\mu)$ be a measure space. For any
$(f_n)_{n\in\mathbb N}\subseteq(m\Sigma)^+$, it holds
\[
\liminf_{n\to\infty}\int_S f_n\,d\mu
\ge \int_S\liminf_{n\to\infty}f_n\,d\mu.
\]

\par\noindent\studyheading{Proof.}\par\nobreak

Recall that
\[
\liminf_{n\to\infty}f_n
=\sup_{n\ge1}\inf_{m\ge n}f_m.
\]
Define
\[
g_n(s):=\inf_{m\ge n}f_m(s),\qquad n\in\mathbb N.
\]
First, $g_n\in(m\Sigma)^+$, since $g_n\ge0$ and, for every $a\in\mathbb R$,
\[
\{g_n<a\}=\bigcup_{m=n}^{\infty}\{f_m<a\}\in\Sigma.
\]
The equality follows from the definition of infimum; the membership follows from
measurability of the $f_m$ and closure of $\Sigma$ under countable unions.

Moreover,
\[
g_n\le g_{n+1},\qquad
 g_n\le f_m\quad\forall m\ge n.
\]
The first inequality follows because the infimum defining $g_{n+1}$ is taken
over a smaller family; the second follows from the definition of $g_n$.
Consequently,
\[
g_n(s)\uparrow\sup_{n\ge1}g_n(s)
=\liminf_{n\to\infty}f_n(s),\qquad\forall s\in S.
\]
By the monotone convergence theorem (Theorem 5.2),
\[
\int_S\liminf_{n\to\infty}f_n\,d\mu
=\lim_{n\to\infty}\int_Sg_n\,d\mu.
\]
Since $g_n\le f_m$ for every $m\ge n$, monotonicity of the positive
Lebesgue integral (HW 45) gives
\[
\int_Sg_n\,d\mu\le\int_Sf_m\,d\mu,
\qquad\forall m\ge n.
\]
Taking the infimum over $m\ge n$, we obtain
\[
\int_Sg_n\,d\mu
\le\inf_{m\ge n}\int_Sf_m\,d\mu.
\]
Therefore,
\[
\begin{aligned}
\int_S\liminf_{n\to\infty}f_n\,d\mu&=\lim_{n\to\infty}\int_Sg_n\,d\mu \le\lim_{n\to\infty}\inf_{m\ge n}\int_Sf_m\,d\mu =\liminf_{n\to\infty}\int_Sf_n\,d\mu.
\end{aligned}
\]
The first equality follows from MON, the inequality from the preceding bound,
and the last equality from the definition of the numerical $\liminf$.

Therefore, the statement is proved.
%
% !TEX root = ../main.tex
\Needspace{21\baselineskip}
\subsection*{Lemma 5.5 --- Reverse Fatou's Lemma for Functions}%
\manualtoc{subsection}{\texorpdfstring{Lemma 5.5 --- Reverse Fatou's Lemma for Functions}{Lemma 5.5 --- Reverse Fatou's Lemma for Functions}}{studiabile-nav-1016}%

\studyheading{Statement.} Let $(f_n)_{n\in\mathbb N}\subseteq (m\Sigma)^+$. Suppose there exists
$g\in(m\Sigma)^+$ such that $f_n\le g$ for all but finitely many $n\in\mathbb N$ and
$\mu(g)<\infty$. Then
\[
\limsup_{n\to\infty}\int_S f_n\,d\mu
\le \int_S\limsup_{n\to\infty}f_n\,d\mu.
\]

\par\noindent\studyheading{Observation.}%
\inlinetoc{subsubsection}{\texorpdfstring{Observation.}{Observation.}}{studiabile-nav-1017}%
\par\nobreak
 The hypothesis means that there exists $n_0\in\mathbb N$ such that
\[
0\le f_n\le g,\qquad\forall n\ge n_0.
\]
Removing the first $n_0-1$ terms does not change either limit superior.
The finite-integral assumption $\mu(g)<\infty$ is essential to the argument.

\par\noindent\studyheading{Proof.}\par\nobreak
 As we will prove in \textbf{HW 48}, the result follows by applying
Fatou's lemma (Lemma 5.4) to the nonnegative functions $g-f_n$ for $n\ge n_0$.
The proof is deferred to that homework, as in the lecture notes.
%
% !TEX root = ../main.tex
\subsection*{HW 48 --- Reverse Fatou's Lemma for Functions}%
\manualtoc{subsection}{\texorpdfstring{HW 48 --- Reverse Fatou's Lemma for Functions}{HW 48 --- Reverse Fatou's Lemma for Functions}}{studiabile-nav-2001}%

\studyheading{Statement.} Let $(f_n)_{n\in\mathbb N}\subseteq(m\Sigma)^+$ and suppose there is $g\in(m\Sigma)^+$ such that $f_n\le g$ for all but finitely many $n$, with $\mu(g)<\infty$. Prove that
\[
\limsup_{n\to\infty}\int_S f_n\,d\mu
\le \int_S\limsup_{n\to\infty}f_n\,d\mu.
\]

\par\noindent\studyheading{Proof.}\par\nobreak

Choose $n_0\in\mathbb N$ such that $f_n\le g$ for every $n\ge n_0$. Removing the first $n_0-1$ terms does not change either limit superior.

Since $\mu(g)<\infty$, the set $N:=\{g=\infty\}$ is null. Indeed, for every $M>0$,
\[
M\mathbf1_N\le g
\quad\Longrightarrow\quad
M\mu(N)=\mu(M\mathbf1_N)\le\mu(g)<\infty.
\]
Letting $M\to\infty$ yields $\mu(N)=0$. Since $f_n\le g$ for $n\ge n_0$, both $g$ and the tail functions $f_n$ are finite outside $N$. By Lemma~5.3, redefining them to be zero on $N$ leaves their integrals and the integral of their pointwise limit superior unchanged. We may therefore assume that $g$ and $f_n$, for $n\ge n_0$, are finite-valued.

For $n\ge n_0$, define
\[
h_n:=g-f_n\in(m\Sigma)^+.
\]
The membership follows because $g$ and $f_n$ are finite measurable functions and $f_n\le g$.

We justify the additivity needed below using only earlier results. For any nonnegative measurable $u,v$, let $u^{(r)}\uparrow u$ and $v^{(r)}\uparrow v$ be the simple approximations from~(5.1). By HW~44(ii) and the monotone convergence theorem (Theorem~5.2),
\[
\begin{aligned}
\mu(u+v)&=\lim_{r\to\infty}\mu_0\bigl(u^{(r)}+v^{(r)}\bigr) =\lim_{r\to\infty}\bigl(\mu_0(u^{(r)})+\mu_0(v^{(r)})\bigr)
=\mu(u)+\mu(v).
\end{aligned}
\]
In particular, $g=f_n+h_n$ gives
\[
\mu(h_n)=\mu(g)-\mu(f_n),\qquad n\ge n_0.
\]
The subtraction is well-defined because $0\le\mu(f_n)\le\mu(g)<\infty$ by HW~45.

Set $L:=\limsup_{n\to\infty}f_n$. Since $f_n\le g$ eventually, $0\le L\le g$, and, because $g$ is finite-valued,
\[
\liminf_{n\to\infty}h_n
=\liminf_{n\to\infty}(g-f_n)
=g-L.
\]
The equality follows from the numerical identity $\liminf(a-x_n)=a-\limsup x_n$ for finite $a$. Additivity also gives $\mu(g-L)=\mu(g)-\mu(L)$; all terms are finite since $0\le L\le g$.

Applying Fatou's lemma (Lemma~5.4) to the nonnegative sequence $(h_n)_{n\ge n_0}$, we obtain
\[
\begin{aligned}
\mu(g)-\limsup_{n\to\infty}\mu(f_n)&=\liminf_{n\to\infty}\mu(h_n) \ge\mu\!\left(\liminf_{n\to\infty}h_n\right) =\mu(g-L)
=\mu(g)-\mu(L).
\end{aligned}
\]
The first equality uses $\mu(h_n)=\mu(g)-\mu(f_n)$, the inequality is Fatou's lemma, and the last equalities use the identities above. Subtracting the finite number $\mu(g)$ and reversing signs yields
\[
\limsup_{n\to\infty}\mu(f_n)
\le\mu(L)
=\mu\!\left(\limsup_{n\to\infty}f_n\right).
\]
Therefore, the statement is proved.
%
% !TEX root = ../main.tex
\subsection*{Lemma 5.6 --- Linearity of the Positive Lebesgue Integral}%
\manualtoc{subsection}{\texorpdfstring{Lemma 5.6 --- Linearity of the Positive Lebesgue Integral}{Lemma 5.6 --- Linearity of the Positive Lebesgue Integral}}{studiabile-nav-2002}%

\studyheading{Statement.} Let $\alpha,\beta\in[0,\infty)$ and $f,g\in(\mS)^+$. Then
\[
\int_S(\alpha f+\beta g)\,d\mu
=\alpha\int_Sf\,d\mu+\beta\int_Sg\,d\mu.
\]

\par\noindent\studyheading{Expanded Proof.}\par\nobreak

Let $(f^{(r)})_{r\in\mathbb N}$ and $(g^{(r)})_{r\in\mathbb N}$ be the sequences of nonnegative simple functions constructed in (5.1), so that
\[
f^{(r)}\uparrow f,\qquad g^{(r)}\uparrow g.
\]
Since $\alpha,\beta\ge0$, we have
\[
\alpha f^{(r)}+\beta g^{(r)}\uparrow\alpha f+\beta g.
\]
Moreover, $\alpha f^{(r)}+\beta g^{(r)}\in SF^+$ by HW 44(ii).

Applying the monotone convergence theorem (Theorem 5.2) and HW 44(ii), we obtain
\[
\begin{aligned}
\int_S(\alpha f+\beta g)\,d\mu&=\lim_{r\to\infty}\int_S(\alpha f^{(r)}+\beta g^{(r)})\,d\mu =\lim_{r\to\infty}\left(\alpha\int_S f^{(r)}\,d\mu
+\beta\int_S g^{(r)}\,d\mu\right)\\
&=\alpha\lim_{r\to\infty}\int_S f^{(r)}\,d\mu
+\beta\lim_{r\to\infty}\int_S g^{(r)}\,d\mu =\alpha\int_S f\,d\mu+\beta\int_S g\,d\mu.
\end{aligned}
\]
The first and last equalities follow from MON. The second uses HW 44(ii), since the approximating functions are simple, and the third follows from the limit rules for nonnegative sequences. The convention $0\cdot\infty=0$ applies when necessary.

Therefore, the statement is proved. 

\subsection*{Positive and Negative Parts of a Measurable Function}%
\manualtoc{subsection}{\texorpdfstring{Positive and Negative Parts of a Measurable Function}{Positive and Negative Parts of a Measurable Function}}{studiabile-nav-2003}%

We now extend the definition of the Lebesgue integral to measurable functions that can change sign.

Given $f\in\mS$, define its positive and negative parts by
\[
f^+:=\max\{f,0\},\qquad f^-:=\max\{-f,0\}.
\]
Then
\[
f=f^+-f^-,\qquad |f|=f^++f^-.
\]
Both functions are nonnegative. Their measurability and the preceding identities are verified in HW 49.

\subsection*{HW 49 --- Positive and Negative Parts}%
\manualtoc{subsection}{\texorpdfstring{HW 49 --- Positive and Negative Parts}{HW 49 --- Positive and Negative Parts}}{studiabile-nav-2004}%

\studyheading{Statement.} Let $f\in\mS$. Prove that
\[
f^+,f^-\in(\mS)^+,\qquad |f|=f^++f^-.
\]

\par\noindent\studyheading{Expanded Proof.}\par\nobreak

By definition,
\[
f^+(s)=\max\{f(s),0\},\qquad f^-(s)=\max\{-f(s),0\},
\]
so $f^+,f^-\ge0$.

\textit{We first verify measurability.} For every $a\in\mathbb R$,
\[
\{f^+>a\}=
\begin{cases}
S,&a<0,\\
\{f>a\},&a\ge0,
\end{cases}
\qquad
\{f^->a\}=
\begin{cases}
S,&a<0,\\
\{f<-a\},&a\ge0.
\end{cases}
\]
All these sets belong to $\Sigma$, since $f$ is measurable. Thus
\[
f^+,f^-\in(\mS)^+.
\]

Finally, for every $s\in S$,
\[
(f^+(s),f^-(s))=
\begin{cases}
(f(s),0),&f(s)\ge0,\\
(0,-f(s)),&f(s)<0.
\end{cases}
\]
In both cases,
\[
f(s)=f^+(s)-f^-(s),\qquad |f(s)|=f^+(s)+f^-(s).
\]
Since $s\in S$ was arbitrary, we conclude that
\[
f=f^+-f^-,\qquad |f|=f^++f^-.
\]
Therefore, the statement is proved. 

\subsection*{Definition --- Lebesgue Integral of a Signed Measurable Function}%
\manualtoc{subsection}{\texorpdfstring{Definition --- Lebesgue Integral of a Signed Measurable Function}{Definition --- Lebesgue Integral of a Signed Measurable Function}}{studiabile-nav-2005}%

Let $f\in\mS$. If
\[
\min\{\mu(f^+),\mu(f^-)\}<\infty,
\]
the Lebesgue integral of $f$ is defined by
\[
\mu(f):=\int_S f\,d\mu
=\int_S f^+\,d\mu-\int_S f^-\,d\mu.
\]
The condition ensures that at least one of the two nonnegative integrals is finite, so the subtraction is well-defined in the extended real numbers.

If both integrals are infinite, the expression would be $\infty-\infty$, which is undefined. If both are finite, then $\mu(f)\in\mathbb R$.
%
% !TEX root = ../main.tex
\subsection*{Definition 5.1 --- $L^1$-Space}%
\manualtoc{subsection}{\texorpdfstring{Definition 5.1 --- $L^1$-Space}{Definition 5.1 --- L1-Space}}{studiabile-nav-2006}%

Let $(S,\Sigma,\mu)$ be a measure space. We denote by $L^1(S,\Sigma,\mu)$ the collection of measurable functions $f\in m\Sigma$ such that
\[
\mu(|f|)=\int_S f^+\,d\mu+\int_S f^-\,d\mu<\infty.
\]
Recall from HW~49 that
\[
f=f^+-f^-,\qquad |f|=f^++f^-,
\]
where $f^+,f^-\in(m\Sigma)^+$. The equality in the definition follows from Lemma~5.6.

In particular, if $f\in L^1(S,\Sigma,\mu)$, both $\mu(f^+)$ and $\mu(f^-)$ are finite. Therefore,
\[
\mu(f)=\mu(f^+)-\mu(f^-)\in\mathbb R.
\]
Such functions are called Lebesgue integrable.

\subsection*{HW 50 --- Absolute Value of the Integral}%
\manualtoc{subsection}{\texorpdfstring{HW 50 --- Absolute Value of the Integral}{HW 50 --- Absolute Value of the Integral}}{studiabile-nav-2007}%

\studyheading{Statement.} For $f\in L^1(S,\Sigma,\mu)$, prove that
\[
|\mu(f)|\le\mu(|f|).
\]
\par\noindent\studyheading{Expanded Proof.}\par\nobreak

Since $f\in L^1(S,\Sigma,\mu)$, we have
\[
\mu(f^+)+\mu(f^-)=\mu(|f|)<\infty.
\]
Thus both integrals are finite. By the definition of the signed Lebesgue integral,
\[
\begin{aligned}
|\mu(f)|&=|\mu(f^+)-\mu(f^-)| \le\mu(f^+)+\mu(f^-) =\mu(|f|).
\end{aligned}
\]
The inequality follows from the triangle inequality for real numbers, since $\mu(f^+),\mu(f^-)\ge0$, and the last equality follows from Definition~5.1.

Therefore, the statement is proved. 

\subsection*{HW 51 --- Closure of $L^1$ under Linear Combinations}%
\manualtoc{subsection}{\texorpdfstring{HW 51 --- Closure of $L^1$ under Linear Combinations}{HW 51 --- Closure of L1 under Linear Combinations}}{studiabile-nav-2008}%

\studyheading{Statement.} Let $f,g\in L^1(S,\Sigma,\mu)$ and $\alpha,\beta\in\mathbb R$. Prove that
\[
\alpha f+\beta g\in L^1(S,\Sigma,\mu).
\]
\par\noindent\studyheading{Expanded Proof.}\par\nobreak

Set $u:=\alpha f+\beta g$. By Lemma~3.4, $u\in m\Sigma$, since $f,g$ are measurable and $\alpha,\beta\in\mathbb R$.

For every $s\in S$, the triangle inequality gives
\[
|u(s)|=|\alpha f(s)+\beta g(s)|
\le|\alpha|\,|f(s)|+|\beta|\,|g(s)|.
\]
By monotonicity (HW~45) and nonnegative linearity (Lemma~5.6),
\[
\begin{aligned}
\mu(|u|)&\le\mu(|\alpha|\,|f|+|\beta|\,|g|) =|\alpha|\,\mu(|f|)+|\beta|\,\mu(|g|) <\infty.
\end{aligned}
\]
The last inequality follows from $f,g\in L^1(S,\Sigma,\mu)$ and the finiteness of $\alpha,\beta$. Therefore $u\in L^1(S,\Sigma,\mu)$.

Therefore, the statement is proved. 

\subsection*{Linearity of the Integral on $L^1$}%
\manualtoc{subsection}{\texorpdfstring{Linearity of the Integral on $L^1$}{Linearity of the Integral on L1}}{studiabile-nav-2009}%

Linearity extends to signed integrable functions. If $f,g\in L^1(S,\Sigma,\mu)$ and $\alpha,\beta\in\mathbb R$, then
\[
\int_S(\alpha f+\beta g)\,d\mu
=\alpha\int_S f\,d\mu+\beta\int_S g\,d\mu.
\]
\par\noindent\studyheading{Expanded Proof.}\par\nobreak

By HW~51, $u:=\alpha f+\beta g\in L^1(S,\Sigma,\mu)$. Define the positive and negative parts of the scalars by
\[
\alpha^+=\max\{\alpha,0\},\quad\alpha^-=\max\{-\alpha,0\},
\qquad
\beta^+=\max\{\beta,0\},\quad\beta^-=\max\{-\beta,0\}.
\]
Consider the nonnegative measurable functions
\[
\begin{aligned}
P&:=\alpha^+f^++\alpha^-f^-+\beta^+g^++\beta^-g^-,\\
Q&:=\alpha^+f^-+\alpha^-f^++\beta^+g^-+\beta^-g^+.
\end{aligned}
\]
Since $f,g\in L^1$ and all coefficients are nonnegative, Lemma~5.6 gives $\mu(P),\mu(Q)<\infty$. Moreover,
\[
P-Q=(\alpha^+-\alpha^-)(f^+-f^-)
+(\beta^+-\beta^-)(g^+-g^-)
=\alpha f+\beta g=u.
\]
Together with $u=u^+-u^-$, this yields the pointwise identity
\[
u^++Q=u^-+P.
\]
By Lemma~5.6,
\[
\mu(u^+)+\mu(Q)=\mu(u^-)+\mu(P).
\]
All terms are finite, since $u\in L^1$ and $\mu(P),\mu(Q)<\infty$. Thus,
\[
\begin{aligned}
\mu(\alpha f+\beta g)&=\mu(u^+)-\mu(u^-) =\mu(P)-\mu(Q)\\
&=(\alpha^+-\alpha^-)(\mu(f^+)-\mu(f^-)) \quad +(\beta^+-\beta^-)(\mu(g^+)-\mu(g^-)) =\alpha\mu(f)+\beta\mu(g).
\end{aligned}
\]
The second equality follows from the preceding integrated identity; the third uses Lemma~5.6 and rearrangement of finite quantities, while the last follows from the signed-integral definition.

Therefore, the statement is proved. 

We will often need to interchange limits and integrals. One result of this kind is MON; the next theorem provides another important tool for doing so.
%
% !TEX root = ../main.tex
\subsection*{Theorem 5.7 --- Dominated Convergence (DOM)}%
\manualtoc{subsection}{\texorpdfstring{Theorem 5.7 --- Dominated Convergence (DOM)}{Theorem 5.7 --- Dominated Convergence (DOM)}}{studiabile-nav-2010}%

\studyheading{Statement.} Suppose $(f_n)_{n\in\mathbb N}\subseteq m\Sigma$ and $f\in m\Sigma$ satisfy $f_n(s)\to f(s)$ for $\mu$-almost every $s\in S$. If
\[
|f_n|\le g,\qquad \forall n\in\mathbb N,
\]
for some $g\in L^1(S,\Sigma,\mu)$, then
\[
\lim_{n\to\infty}\int_S|f_n-f|\,d\mu=0,
\qquad
\lim_{n\to\infty}\int_S f_n\,d\mu=\int_S f\,d\mu.
\]

\par\noindent\studyheading{Proof.}\par\nobreak

Since $|f_n|\le g$ and $g\in L^1(S,\Sigma,\mu)$, monotonicity (HW 45) gives
\[
0\le\mu(|f_n|)\le\mu(g)<\infty,\qquad \forall n\in\mathbb N.
\]
Thus every $f_n$ is integrable. Let $N\in\Sigma$ be a null set outside which $f_n(s)\to f(s)$. Passing to the pointwise limit in $|f_n(s)|\le g(s)$ yields
\[
|f(s)|\le g(s),\qquad \forall s\in N^c.
\]
As in Lemma 5.3, modifying functions on $N$ does not affect their integrals. We may therefore replace $f_n$ and $f$ by $f_n\mathbf1_{N^c}$ and $f\mathbf1_{N^c}$, respectively. After this replacement, $f_n\to f$ everywhere and $|f_n|,|f|\le g$ everywhere. In particular,
\[
\mu(|f|)\le\mu(g)<\infty,
\]
so $f\in L^1(S,\Sigma,\mu)$ as well. The original functions have the same integrals as these modifications, by Lemma 5.3 applied to their positive and negative parts.

By the triangle inequality,
\[
0\le |f_n-f|\le |f_n|+|f|\le 2g,
\qquad \forall n\in\mathbb N.
\]
Since $g\ge0$ and $g\in L^1$, Lemma 5.6 gives $\mu(2g)=2\mu(g)<\infty$. We may therefore apply reverse Fatou's lemma (Lemma 5.5) to the sequence $|f_n-f|$:
\[
\begin{aligned}
0\le\limsup_{n\to\infty}\int_S|f_n-f|\,d\mu&\le\int_S\limsup_{n\to\infty}|f_n-f|\,d\mu =\int_S0\,d\mu=0.
\end{aligned}
\]
The first inequality is nonnegativity of the integrals, the second follows from reverse Fatou and the bound by $2g$, and the equality uses $f_n\to f$ pointwise. Thus
\[
\lim_{n\to\infty}\int_S|f_n-f|\,d\mu=0.
\]

Finally, by linearity of the integral on $L^1$ and HW 50,
\[
\begin{aligned}
\left|\int_Sf_n\,d\mu-\int_Sf\,d\mu\right|&=\left|\int_S(f_n-f)\,d\mu\right| \le\int_S|f_n-f|\,d\mu\longrightarrow0.
\end{aligned}
\]
The inequality is valid since $f_n-f\in L^1$ by HW 51. This proves the second assertion.

Therefore, the statement is proved. 

\subsection*{HW 52 --- Why the Domination Hypothesis Is Necessary}%
\manualtoc{subsection}{\texorpdfstring{HW 52 --- Why the Domination Hypothesis Is Necessary}{HW 52 --- Why the Domination Hypothesis Is Necessary}}{studiabile-nav-2011}%

\studyheading{Statement.} Let $(S,\Sigma,\mu)=([0,1],\mathcal B([0,1]),\mathrm{Leb})$ and define
\[
f_n(s)=n\mathbf1_{(0,\,1/n)}(s),\qquad n\in\mathbb N.
\]
Show that the conclusion of DOM fails without an integrable dominating function.

\par\noindent\studyheading{Proof.}\par\nobreak

Fix $s\in(0,1]$. For every $n\ge1/s$, we have $s\notin(0,1/n)$, so $f_n(s)=0$. Moreover, $f_n(0)=0$ for every $n$. Hence
\[
f_n(s)\longrightarrow f(s):=0,\qquad \forall s\in[0,1].
\]
Nevertheless,
\[
\begin{aligned}
\int_0^1 f_n\,d\mathrm{Leb}
&=n\,\mathrm{Leb}\bigl((0,1/n)\bigr)=1,\qquad \forall n\in\mathbb N,\\
\int_0^1 f\,d\mathrm{Leb}&=0.
\end{aligned}
\]
The first equality uses the integral of a nonnegative simple function; the second is the integral of the zero function. Consequently,
\[
\lim_{n\to\infty}\int_0^1 f_n\,d\mathrm{Leb}=1
\ne0=\int_0^1 f\,d\mathrm{Leb}.
\]

\textit{We now verify directly which hypothesis of DOM fails.} Suppose that a function 
\[
g\in L^1([0,1],\mathcal B([0,1]),\mathrm{Leb})
\]
 satisfies $f_n\le g$ for every $n$. For each $k\in\mathbb N$, set
\[
I_k:=\left(\frac1{k+1},\frac1k\right).
\]
On $I_k$ we have $f_k=k$, so $g\ge k$. Since the intervals $I_k$ are pairwise disjoint, for every $N\in\mathbb N$,
\[
\begin{aligned}
\mu(g)&\ge\mu\left(\sum_{k=1}^{N}k\mathbf1_{I_k}\right)
=\sum_{k=1}^{N}k\,\mathrm{Leb}(I_k) =\sum_{k=1}^{N}k\left(\frac1k-\frac1{k+1}\right)
=\sum_{k=1}^{N}\frac1{k+1}.
\end{aligned}
\]
The inequality follows from HW 45, and the first equality uses the definition of the simple integral. As $N\to\infty$, the last sum diverges, contradicting $\mu(g)<\infty$.

Thus no single integrable function dominates all $f_n$. The example does not contradict DOM: its domination hypothesis is not satisfied.

Therefore, the statement is proved.
%
% !TEX root = ../main.tex
\subsection*{The Standard Machinery}%
\manualtoc{subsection}{\texorpdfstring{The Standard Machinery}{The Standard Machinery}}{studiabile-nav-3001}%

Many properties of Lebesgue integrals can be established through the following four steps:
\begin{enumerate}
\item Prove the property for indicators $\mathbf1_A$, where $A\in\Sigma$.
\item Use linearity to extend it to nonnegative simple functions.
\item Approximate nonnegative measurable functions by increasing simple functions and apply the monotone convergence theorem (MON, Theorem~5.2).
\item Use $f=f^+-f^-$ to extend the result to signed measurable functions whenever the relevant integrals are well-defined.
\end{enumerate}
This is the standard machinery used throughout the course.

\subsection*{A Measure Defined by a Nonnegative Density}%
\manualtoc{subsection}{\texorpdfstring{A Measure Defined by a Nonnegative Density}{A Measure Defined by a Nonnegative Density}}{studiabile-nav-3002}%

Different measures on the same measurable space yield different Lebesgue integrals. Fix $f\in(m\Sigma)^+$ and define, for every $A\in\Sigma$,
\[
\lambda(A):=\int_S f\mathbf1_A\,d\mu=\int_A f\,d\mu.
\]
The following homework proves that $\lambda:\Sigma\to[0,\infty]$ is a measure. In addition, $\lambda$ is absolutely continuous with respect to $\mu$, written $\lambda\ll\mu$, meaning
\[
\mu(N)=0\quad\Longrightarrow\quad\lambda(N)=0,
\qquad\forall N\in\Sigma.
\]

\subsection*{HW 53 --- Countable Additivity of $\lambda$}%
\manualtoc{subsection}{\texorpdfstring{HW 53 --- Countable Additivity of $\lambda$}{HW 53 --- Countable Additivity of lambda}}{studiabile-nav-3003}%

\studyheading{Statement.} Let $f\in(m\Sigma)^+$ and define $\lambda(A)=\int_S f\mathbf1_A\,d\mu$ for every $A\in\Sigma$. Prove that $\lambda$ is countably additive and hence is a measure on $(S,\Sigma)$.

\par\noindent\studyheading{Proof.}\par\nobreak

First, $\lambda(A)\ge0$ for every $A\in\Sigma$, since $f\mathbf1_A\ge0$, and
\[
\lambda(\emptyset)=\int_S f\mathbf1_{\emptyset}\,d\mu
=\int_S0\,d\mu=0.
\]
The equality follows from $\mathbf1_{\emptyset}=0$ and the definition of the positive integral.

Let $(A_n)_{n\in\mathbb N}\subseteq\Sigma$ be pairwise disjoint. Set
\[
U_r:=\bigcup_{n=1}^{r}A_n,
\qquad U:=\bigcup_{n=1}^{\infty}A_n.
\]
Since the sets are pairwise disjoint and $f\ge0$,
\[
f\mathbf1_{U_r}=\sum_{n=1}^{r}f\mathbf1_{A_n}
\uparrow f\mathbf1_U\qquad\text{pointwise on }S.
\]
The equality follows from $\mathbf1_{U_r}=\sum_{n=1}^{r}\mathbf1_{A_n}$; the monotone convergence follows from $U_r\uparrow U$.

Therefore,
\[
\begin{aligned}
\lambda\left(\bigcup_{n=1}^{\infty}A_n\right)&=\int_S f\mathbf1_U\,d\mu =\lim_{r\to\infty}\int_S f\mathbf1_{U_r}\,d\mu =\lim_{r\to\infty}\sum_{n=1}^{r}\int_S f\mathbf1_{A_n}\,d\mu =\sum_{n=1}^{\infty}\lambda(A_n).
\end{aligned}
\]
The second equality follows from MON (Theorem~5.2), the third from nonnegative linearity (Lemma~5.6), and the last from the definition of an infinite nonnegative sum. Hence $\lambda$ is countably additive and therefore is a measure on $(S,\Sigma)$.

Moreover, let $N\in\Sigma$ satisfy $\mu(N)=0$. Define
\[
h_r:=(f\wedge r)\mathbf1_N,\qquad r\in\mathbb N.
\]
Then $h_r\in(m\Sigma)^+$, $h_r\uparrow f\mathbf1_N$, and $0\le h_r\le r\mathbf1_N$. By monotonicity (HW~45),
\[
0\le\mu(h_r)\le\mu(r\mathbf1_N)=r\mu(N)=0.
\]
The equality uses the integral of a simple indicator. Applying MON once more gives
\[
\lambda(N)=\int_S f\mathbf1_N\,d\mu
=\lim_{r\to\infty}\mu(h_r)=0.
\]
Thus $\lambda\ll\mu$.

Therefore, the statement is proved. 

\subsection*{Theorem 5.8 --- Radon--Nikodym Theorem}%
\manualtoc{subsection}{\texorpdfstring{Theorem 5.8 --- Radon--Nikodym Theorem}{Theorem 5.8 --- Radon--Nikodym Theorem}}{studiabile-nav-3004}%

\studyheading{Statement.} Let $\mu$ and $\lambda$ be $\sigma$-finite measures on $(S,\Sigma)$ such that $\lambda\ll\mu$. Then there exists $f\in(m\Sigma)^+$ satisfying
\[
\lambda(A)=\int_S f\mathbf1_A\,d\mu,
\qquad\forall A\in\Sigma.
\]
The function
\[
f=: \frac{d\lambda}{d\mu}
\]
is called the Radon--Nikodym derivative of $\lambda$ with respect to $\mu$.

\par\noindent\studyheading{Recall.}%
\inlinetoc{subsubsection}{\texorpdfstring{Recall.}{Recall.}}{studiabile-nav-3005}%
\par\nobreak
 A measure $\nu$ is $\sigma$-finite if there exist measurable sets $(E_n)_{n\in\mathbb N}$ such that
\[
S=\bigcup_{n=1}^{\infty}E_n,
\qquad\nu(E_n)<\infty\quad\forall n\in\mathbb N.
\]
HW~53 constructs a measure $\lambda$ from a given nonnegative density $f$; the Radon--Nikodym theorem gives the converse under $\sigma$-finiteness and absolute continuity. As specified in the lecture notes, the theorem is stated without proof.
%
\clearpage
%
\end{document}
